% Circle Geometry — Pure Mathematics with Statistics Upper Sixth — Cours signé OnBuch+
% Official MINESEC syllabus. Compiler avec : tectonic PMS01-circle-geometry.tex
\newcommand{\DOCMATIERE}{Pure Mathematics with Statistics}
\newcommand{\DOCNIVEAU}{Upper Sixth}
\newcommand{\DOCMODULE}{Upper Sixth — Pure mathematics with statistics}
\newcommand{\DOCLECON}{1}
\newcommand{\DOCTITRE}{Circle Geometry}
% OnBuch+ — préambule « cours signés » ENRICHI (compilé avec tectonic / XeLaTeX)
% Système visuel « Pop » d'OnBuch+ : fond crème (jamais blanc), bordures encre
% épaisses, ombres « dures » décalées, titres Archivo Black en capitales.
% Version MATHÉMATIQUES : graphes (pgfplots/tikz), boîtes pédagogiques
% (définition, propriété, méthode, attention, le savais-tu, exercice résolu,
% exercice, TP, bilan), page de garde et signature OnBuch+ ancrée en bas.
%
% Macros attendues dans main.tex AVANT \input{preamble} :
%   \DOCMATIERE  \DOCNIVEAU  \DOCMODULE  \DOCLECON (numéro)  \DOCTITRE
\documentclass[11pt]{article}

\usepackage{fontspec}
\usepackage[english]{babel}
\usepackage[a4paper,margin=2cm,top=3.4cm,bottom=2.8cm,headheight=1.4cm,footskip=1.3cm]{geometry}
\usepackage{amsmath,amssymb,amsfonts}
\usepackage{mathtools} % \xmapsto, \coloneqq... souvent utilisés par les modèles
\usepackage{cancel} % simplifications barrées (bilans de forces)
% \abs{z} (module, valeur absolue) et \norm{u} (norme), utilisés par les modèles
\providecommand{\abs}{}\let\abs\relax\DeclarePairedDelimiter{\abs}{\lvert}{\rvert}
\providecommand{\norm}{}\let\norm\relax\DeclarePairedDelimiter{\norm}{\lVert}{\rVert}
\usepackage[table]{xcolor} % [table] : fournit aussi \rowcolors (alternance de lignes)
\usepackage{pagecolor}
\usepackage{tikz}
\usetikzlibrary{arrows.meta,positioning,calc,shapes.geometric,shapes.misc,shapes.arrows,decorations.pathmorphing,decorations.pathreplacing,decorations.markings,patterns,shadows,mindmap,trees,fit,backgrounds,matrix,intersections,angles,quotes}
\usepackage{pgfplots}
\usepackage[version=4]{mhchem} % équations nucléaires \ce{^{235}_{92}U}
\usepackage[european,siunitx]{circuitikz} % schémas électriques
\pgfplotsset{compat=1.18}
\usepgfplotslibrary{fillbetween,groupplots} % aires entre courbes (calcul intégral)
\usepackage{enumitem}
\usepackage{fancyhdr}
% [most] charge toutes les bibliothèques tcolorbox (listings, minted, raster,
% documentation...) et gonfle inutilement la mémoire TeX sur un document long
% (~300 boîtes) : on ne charge que ce qui est réellement utilisé.
\usepackage[skins,breakable]{tcolorbox}
\usepackage{graphicx}
\usepackage{colortbl}
\usepackage{booktabs}
\usepackage{tabularx}
\usepackage{diagbox} % case d'en-tête diagonale des tableaux à double entrée
% Colonnes extensibles centrée / gauche / droite (C, L, R), souvent utilisées par les modèles
\newcolumntype{C}{>{\centering\arraybackslash}X}
\newcolumntype{L}{>{\raggedright\arraybackslash}X}
\newcolumntype{R}{>{\raggedleft\arraybackslash}X}
\usepackage{array}
\usepackage{multicol}
\usepackage{siunitx}
% parse-numbers=false : accepte aussi \SI{2,5 \times 10^{-3}}{...}
\sisetup{locale=UK,output-decimal-marker={.},group-minimum-digits=4,parse-numbers=false}
\DeclareSIUnit\ohm{\ensuremath{\symup{\Omega}}} % U+2126 absent de la police
\DeclareSIUnit\division{div} % réglages d'oscilloscope (ms/div, V/div)
\DeclareSIUnit\tr{tr} % tours (vitesse de rotation, ex. \SI{1500}{\tr\per\minute})
\usepackage{qrcode}
% \degree : commande fréquemment utilisée par les modèles pour un angle ou une
% température en dehors de \SI{}{} (non fournie par les paquets chargés).
\providecommand{\degree}{^{\circ}}
% \qed (fin de démonstration), utilisé par les modèles sans amsthm
\providecommand{\qed}{\hfill$\square$}
% \barre : double barre des valeurs interdites dans un tableau de variations (tkz-tab)
\providecommand{\barre}{\ensuremath{\|}}
% environnement proof (amsthm non chargé) utilisé par les modèles
\ifdefined\proof\else\newenvironment{proof}[1][Proof]{\par\noindent\textit{#1.}\ }{\qed\par}\fi
\usepackage{lastpage}
\usepackage{hyperref}
\hypersetup{hidelinks}

% ============================================================
% Palette « Pop » OnBuch+ — mêmes hex que la classe P (plus_ui.dart)
% ============================================================
\definecolor{poporange}{HTML}{F59321}
\definecolor{popdark}{HTML}{E07A0C}
\definecolor{popcream}{HTML}{FFF6E8}
\definecolor{popink}{HTML}{1C1714}
\definecolor{poppurple}{HTML}{7A5AE0}
\definecolor{popgreen}{HTML}{2E9E6A}
\definecolor{poppink}{HTML}{E0457B}
\definecolor{popblue}{HTML}{2F7DE1}
\definecolor{popgold}{HTML}{C98A12}
\definecolor{popmuted}{HTML}{8A7F76}
\definecolor{popline}{HTML}{EADFD2}
\definecolor{popgray}{HTML}{8A7F76}
\colorlet{popgrey}{popgray}
% Alias tolérants (le modèle nomme parfois les couleurs différemment)
\colorlet{popwhite}{white}
\colorlet{popblack}{popink}
\colorlet{popred}{poppink}
\colorlet{popyellow}{popgold}
% Teintes claires pour fonds de tableaux / zones
\colorlet{popblueL}{popblue!10!white}
\colorlet{poporangeL}{poporange!14!white}
\colorlet{popgreenL}{popgreen!12!white}
\colorlet{poppurpleL}{poppurple!12!white}
\colorlet{poppinkL}{poppink!12!white}
\colorlet{popgoldL}{popgold!14!white}

\pagecolor{popcream}

% ============================================================
% Typographie
% ============================================================
\defaultfontfeatures{Ligatures=TeX}
\setmainfont{PlusJakartaSans-Medium.ttf}[Path=../../../fonts/,BoldFont=PlusJakartaSans-Bold.ttf]
% Maths en sans-serif (Fira Math) avec chiffres et lettres droites de Plus
% Jakarta, pour que les formules se fondent dans le texte.
\usepackage[math-style=ISO,bold-style=ISO]{unicode-math}
\setmathfont{FiraMath-Regular.otf}[Path=../../../fonts/]
\setmathfont{PlusJakartaSans-Medium.ttf}[Path=../../../fonts/,range={up/num,up/{Latin,latin},\mathup}]
% (icomma retiré : décimales avec point en anglais)
% Caractères mathématiques Unicode absents des polices de texte (Plus Jakarta,
% Archivo Black) : rendus par leur équivalent mathématique, en texte comme en
% formule — sinon ils s'affichent en carré vide (ex. « Arithmétique dans ℤ »).
\usepackage{newunicodechar}
\newunicodechar{ℝ}{\ensuremath{\mathbb{R}}}
\newunicodechar{ℤ}{\ensuremath{\mathbb{Z}}}
\newunicodechar{ℂ}{\ensuremath{\mathbb{C}}}
\newunicodechar{ℕ}{\ensuremath{\mathbb{N}}}
\newunicodechar{ℚ}{\ensuremath{\mathbb{Q}}}
\newunicodechar{𝔻}{\ensuremath{\mathbb{D}}}
\newunicodechar{α}{\ensuremath{\alpha}}
\newunicodechar{β}{\ensuremath{\beta}}
\newunicodechar{γ}{\ensuremath{\gamma}}
\newunicodechar{δ}{\ensuremath{\delta}}
\newunicodechar{ε}{\ensuremath{\varepsilon}}
\newunicodechar{θ}{\ensuremath{\theta}}
\newunicodechar{λ}{\ensuremath{\lambda}}
\newunicodechar{μ}{\ensuremath{\mu}}
\newunicodechar{σ}{\ensuremath{\sigma}}
\newunicodechar{τ}{\ensuremath{\tau}}
\newunicodechar{φ}{\ensuremath{\varphi}}
\newunicodechar{ψ}{\ensuremath{\psi}}
\newunicodechar{ω}{\ensuremath{\omega}}
\newunicodechar{Σ}{\ensuremath{\Sigma}}
% majuscules produites par \MakeUppercase dans les titres (α-aminés -> Α)
\newunicodechar{Α}{\ensuremath{\alpha}}
\newunicodechar{Β}{\ensuremath{\beta}}
\newunicodechar{Γ}{\ensuremath{\Gamma}}
\newunicodechar{Δ}{\ensuremath{\Delta}}
\newunicodechar{Ε}{\ensuremath{\varepsilon}}
\newunicodechar{Θ}{\ensuremath{\Theta}}
\newunicodechar{Λ}{\ensuremath{\Lambda}}
\newunicodechar{Μ}{\ensuremath{\mu}}
\newunicodechar{Τ}{\ensuremath{\tau}}
\newunicodechar{Φ}{\ensuremath{\Phi}}
\newunicodechar{Ψ}{\ensuremath{\Psi}}
\newunicodechar{Ω}{\ensuremath{\Omega}}
\newunicodechar{↦}{\ensuremath{\mapsto}}
\newunicodechar{∈}{\ensuremath{\in}}
\newunicodechar{∉}{\ensuremath{\notin}}
\newunicodechar{∧}{\ensuremath{\wedge}}
\newunicodechar{⇔}{\ensuremath{\Leftrightarrow}}
\newunicodechar{⟺}{\ensuremath{\Longleftrightarrow}}
\newunicodechar{⇒}{\ensuremath{\Rightarrow}}
\newunicodechar{⟹}{\ensuremath{\Longrightarrow}}
\newunicodechar{∘}{\ensuremath{\circ}}
\newunicodechar{∪}{\ensuremath{\cup}}
\newunicodechar{∩}{\ensuremath{\cap}}
\newunicodechar{≡}{\ensuremath{\equiv}}
\newunicodechar{⊂}{\ensuremath{\subset}}
\newunicodechar{⊥}{\ensuremath{\perp}}
\newunicodechar{∃}{\ensuremath{\exists}}
\newunicodechar{∀}{\ensuremath{\forall}}
\newunicodechar{∛}{\ensuremath{\sqrt[3]{\,}}}
\newunicodechar{∜}{\ensuremath{\sqrt[4]{\,}}}
\newunicodechar{⃗}{\ensuremath{{}^{\to}}}
\newunicodechar{ⁿ}{\ifmmode^{n}\else\textsuperscript{\ensuremath{n}}\fi}
\newunicodechar{ˣ}{\ifmmode^{x}\else\textsuperscript{\ensuremath{x}}\fi}
\newunicodechar{ᵇ}{\ifmmode^{b}\else\textsuperscript{\ensuremath{b}}\fi}
\newunicodechar{ʳ}{\ifmmode^{r}\else\textsuperscript{\ensuremath{r}}\fi}
\newunicodechar{ᵘ}{\ifmmode^{u}\else\textsuperscript{\ensuremath{u}}\fi}
\newunicodechar{ʸ}{\ifmmode^{y}\else\textsuperscript{\ensuremath{y}}\fi}
\newunicodechar{ᵃ}{\ifmmode^{a}\else\textsuperscript{\ensuremath{a}}\fi}
\newunicodechar{ᵉ}{\ifmmode^{e}\else\textsuperscript{\ensuremath{e}}\fi}
\newunicodechar{ᵏ}{\ifmmode^{k}\else\textsuperscript{\ensuremath{k}}\fi}
\newunicodechar{ᵖ}{\ifmmode^{p}\else\textsuperscript{\ensuremath{p}}\fi}
\newunicodechar{ᴺ}{\ifmmode^{N}\else\textsuperscript{\ensuremath{N}}\fi}
\newunicodechar{ᶜ}{\ifmmode^{c}\else\textsuperscript{\ensuremath{c}}\fi}
\newunicodechar{ʰ}{\ifmmode^{h}\else\textsuperscript{\ensuremath{h}}\fi}
\newunicodechar{ᵠ}{\ifmmode^{q}\else\textsuperscript{\ensuremath{q}}\fi}
\newunicodechar{ₙ}{\ifmmode_{n}\else\textsubscript{\ensuremath{n}}\fi}
\newunicodechar{ₐ}{\ifmmode_{a}\else\textsubscript{\ensuremath{a}}\fi}
\newunicodechar{ᵢ}{\ifmmode_{i}\else\textsubscript{\ensuremath{i}}\fi}
\newunicodechar{ᵣ}{\ifmmode_{r}\else\textsubscript{\ensuremath{r}}\fi}
\newunicodechar{ₖ}{\ifmmode_{k}\else\textsubscript{\ensuremath{k}}\fi}
\newunicodechar{ᵦ}{\ifmmode_{\beta}\else\textsubscript{\ensuremath{\beta}}\fi}
\newunicodechar{ₓ}{\ifmmode_{x}\else\textsubscript{\ensuremath{x}}\fi}
\newfontfamily\popdisplay{ArchivoBlack-Regular.ttf}[Path=../../../fonts/]
\newfontfamily\poplogo{SpaceGrotesk-Bold.ttf}[Path=../../../fonts/]
\newfontfamily\popsemi{PlusJakartaSans-SemiBold.ttf}[Path=../../../fonts/]
\newfontfamily\popxbold{PlusJakartaSans-ExtraBold.ttf}[Path=../../../fonts/]
\newfontfamily\popxboldls{PlusJakartaSans-ExtraBold.ttf}[Path=../../../fonts/,LetterSpace=3.0]

\setlist[enumerate]{leftmargin=1.7em,itemsep=5pt,topsep=4pt}
\setlist[itemize]{leftmargin=1.7em,itemsep=4pt,topsep=4pt,label={\color{poporange}\textbullet}}
\setlength{\parindent}{0pt}
\setlength{\parskip}{5pt}
\renewcommand{\arraystretch}{1.25}

% pgfplots : style Pop par défaut
\pgfplotsset{
  every axis/.append style={axis line style={popink,line width=1pt},tick style={popink},
    grid style={popline,line width=0.5pt},label style={font=\small},tick label style={font=\footnotesize}},
  popcurve/.style={poporange,line width=1.8pt,no markers,smooth},
}
% Même style disponible dans un \draw TikZ simple (hors pgfplots)
\tikzset{popcurve/.style={poporange,line width=1.8pt}}
\tikzset{popgrid/.style={popline,very thin}}
\colorlet{popcurve}{poporange} % le nom du style est parfois utilisé comme couleur

% ============================================================
% Pill / squiggle
% ============================================================
\newcommand{\poppill}[3]{%
  \tikz[baseline=-0.55ex]\node[draw=popink,line width=1.2pt,rounded corners=2.6pt,%
    fill=#2,inner xsep=6.5pt,inner ysep=3pt]{%
    \popxboldls\fontsize{7}{7}\selectfont\color{#3}\MakeUppercase{#1}};%
}
\newcommand{\popsquiggle}[1]{%
  \tikz[baseline=-0.4ex]{
    \draw[#1,line width=2.6pt,line cap=round]
      plot [smooth] coordinates {(0,0.09) (0.34,0.01) (0.68,0.21) (1.02,0.01) (1.36,0.10)};
  }%
}
% Mot-clé surligné (marqueur orange sous le texte)
\newcommand{\cle}[1]{{\popxbold\color{popdark}#1}}

% ============================================================
% En-tête / pied
% ============================================================
\pagestyle{fancy}
\fancyhf{}
\renewcommand{\headrulewidth}{0pt}
\renewcommand{\footrulewidth}{0pt}
\lhead{\raisebox{-0.15ex}{\poplogo\fontsize{13}{13}\selectfont\color{popink}OnBuch\color{poporange}+}}
\rhead{\poppill{\DOCMATIERE\ \textbullet\ \DOCNIVEAU\ \textbullet\ Chapter \DOCLECON}{white}{popink}}
\lfoot{\popxbold\fontsize{7.6}{7.6}\selectfont\color{popmuted}\MakeUppercase{Signed course OnBuch+}\ \textbullet\ {\popsemi\DOCTITRE}}
\rfoot{\tikz[baseline=-0.6ex]\node[rounded corners=8.5pt,fill=popink,minimum height=17pt,minimum width=17pt,inner xsep=5pt,inner sep=0]{\popxbold\fontsize{8}{8}\selectfont\color{popcream}\thepage\,/\,\pageref*{LastPage}};}

% ============================================================
% Titres
% ============================================================
\newcommand{\chaptitre}[2]{%
  \begin{tcolorbox}[enhanced,colback=poporange,colframe=popink,boxrule=2.6pt,arc=11pt,
      drop shadow=popink,left=15pt,right=15pt,top=12pt,bottom=11pt,before skip=3pt,after skip=18pt]
    \poppill{#2}{popink}{popcream}\par\vspace{9pt}
    {\raggedright\hyphenpenalty=10000\exhyphenpenalty=10000\popdisplay\fontsize{22}{24}\selectfont\color{popcream}\MakeUppercase{#1}\par}
  \end{tcolorbox}
}
% Section numérotée : pastille orange avec le numéro + titre Archivo
\newcounter{popsec}
\newcommand{\coursec}[1]{%
  \stepcounter{popsec}\setcounter{popsub}{0}%
  \par\vspace{16pt}\noindent
  \begin{minipage}{\linewidth}
  \tikz[baseline=-0.7ex]\node[draw=popink,line width=1.6pt,fill=poporange,rounded corners=4pt,minimum size=22pt,inner sep=0]{\popdisplay\fontsize{12}{12}\selectfont\color{popcream}\thepopsec};%
  \hspace{8pt}{\popdisplay\fontsize{15}{17}\selectfont\color{popink}\MakeUppercase{#1}}\par
  \vspace{4pt}\noindent{\color{poporange}\rule{36pt}{3.6pt}}
  \end{minipage}\par\nopagebreak\vspace{8pt}
}
\newcounter{popsub}
\newcommand{\courssub}[1]{%
  \stepcounter{popsub}%
  \par\vspace{9pt}\noindent{\popxbold\fontsize{11.5}{13}\selectfont\color{popink}{\color{poporange}\thepopsec.\thepopsub}\hspace{6pt}#1}\par\nopagebreak\vspace{4pt}
}

% ============================================================
% Boîtes pédagogiques — même recette : fond blanc, bordure épaisse colorée,
% ombre dure encre, badge plein. Titre optionnel [#1].
% ============================================================
\tcbset{popbase/.style={breakable,enhanced,colback=white,boxrule=2.3pt,arc=9pt,
  shadow={3pt}{-3pt}{0pt}{popink},left=14pt,right=14pt,top=11pt,bottom=11pt,before skip=11pt,after skip=15pt,
  fontupper=\popsemi\fontsize{10.6}{14.5}\selectfont\color{popink},
  fontlower=\popsemi\fontsize{10.6}{14.5}\selectfont\color{popink}}}
% Coche dessinée (le glyphe ✓ n'existe pas dans les polices du document)
\newcommand{\popcheck}[1]{\tikz[baseline=-0.5ex]\draw[#1,line width=1.6pt,line cap=round,line join=round](-0.13,0.0)--(-0.04,-0.09)--(0.13,0.1);}
\providecommand{\coche}{\popcheck{popgreen}}
\providecommand{\resultat}[1]{\par\smallskip\noindent\fcolorbox{popgreen}{popgreenL}{\parbox{\dimexpr\linewidth-2\fboxsep-2\fboxrule\relax}{\textbf{Result.} #1}}\par\smallskip}
\newcommand{\popboxhead}[3]{\poppill{#1}{#2}{white}\ifx\relax#3\relax\else\hspace{6pt}{\popxbold\fontsize{10.6}{12}\selectfont\color{popink}#3}\fi\par\vspace{7pt}}

\newtcolorbox{aretenir}[1][]{popbase,colframe=popgreen,before upper={\popboxhead{Key points}{popgreen}{#1}}}
\newtcolorbox{exemplebox}[1][]{popbase,colframe=poporange,before upper={\popboxhead{Exemple}{poporange}{#1}}}
\newtcolorbox{definition}[1][]{popbase,colframe=popblue,before upper={\popboxhead{Definition}{popblue}{#1}}}
\newtcolorbox{propriete}[1][]{popbase,colframe=poppurple,before upper={\popboxhead{Property}{poppurple}{#1}}}
\newtcolorbox{methode}[1][]{popbase,colframe=popgold,before upper={\popboxhead{Method}{popgold}{#1}}}
\newtcolorbox{attention}[1][]{popbase,colframe=poppink,before upper={\popboxhead{Warning}{poppink}{#1}}}
\newtcolorbox{experience}[1][]{popbase,colframe=popblue,colback=popblueL,before upper={\popboxhead{Experiment}{popblue}{#1}}}
% « Le savais-tu ? » : carte violette pleine (contexte, culture, Cameroun)
\newtcolorbox{savaistu}[1][]{popbase,colback=poppurple,colframe=popink,
  fontupper=\popsemi\fontsize{10.4}{14.2}\selectfont\color{white},
  before upper={\poppill{Did you know?}{white}{poppurple}\ifx\relax#1\relax\else\hspace{6pt}{\popxbold\color{white}#1}\fi\par\vspace{7pt}}}
% Exercice résolu : énoncé en haut, \tcblower puis la solution sur fond crème
\newcounter{popexo}\newcounter{popexr}
\newtcolorbox{exoresolu}[1][]{popbase,colframe=poporange,colbacklower=popcream,
  segmentation style={popink,line width=1.2pt,dashed},
  before upper={\stepcounter{popexr}\popboxhead{Worked example \thepopexr}{poporange}{#1}},
  before lower={\poppill{Solution}{popink}{popcream}\par\vspace{6pt}}}
% Exercice (non résolu en place) — la correction est dans la partie Corrigés
\newtcolorbox{exercice}[1][]{popbase,colframe=popink,
  before upper={\stepcounter{popexo}\popboxhead{Exercise \thepopexo}{popink}{#1}}}
% Corrigé
\newtcolorbox{corrige}[1][]{popbase,colframe=popgreen,colback=popgreenL,
  before upper={\popboxhead{Solution}{popgreen}{#1}}}
% Objectifs / prérequis / bilan
\newtcolorbox{objectifs}{popbase,colframe=popink,colback=poporangeL,
  before upper={\popboxhead{Chapter objectives}{popink}{}}}
\newtcolorbox{prerequis}{popbase,colframe=popink,colback=popblueL,
  before upper={\popboxhead{Prerequisites}{popblue}{}}}
\newtcolorbox{bilan}{popbase,colframe=popink,colback=white,boxrule=2.8pt,
  before upper={\popboxhead{Fiche bilan}{poporange}{L'essentiel à connaître pour l'examen}}}
% Figure encadrée avec légende
\newtcolorbox{popfigure}[1][]{enhanced,breakable=false,colback=white,colframe=popline,boxrule=1.4pt,arc=8pt,
  left=10pt,right=10pt,top=10pt,bottom=8pt,before skip=10pt,after skip=12pt,halign=center,
  lower separated=false,halign lower=center,
  fontlower=\popsemi\fontsize{9}{11.5}\selectfont\color{popmuted},
  #1}
\newcommand{\legende}[1]{\par\vspace{6pt}{\popsemi\fontsize{9}{11.5}\selectfont\color{popmuted}#1}}

% ============================================================
% Signature OnBuch+
% ============================================================
\newcommand{\poplogoinline}[1]{{\poplogo\fontsize{#1}{#1}\selectfont\color{popink}OnBuch\color{poporange}+}}
\newcommand{\signatureonbuch}{%
  \begin{tcolorbox}[enhanced,colback=popink,colframe=popink,boxrule=2.2pt,arc=10pt,
      drop shadow=poporange,left=14pt,right=14pt,top=10pt,bottom=10pt,before skip=0pt,after skip=0pt]
    \tikz[baseline=-0.7ex]\node[circle,fill=popgreen,draw=popcream,line width=1.3pt,minimum size=22pt,inner sep=0]{\popcheck{white}};%
    \hspace{9pt}\begin{minipage}[c]{0.62\linewidth}
      {\popxbold\fontsize{12}{13}\selectfont\color{popcream}Signed\ \textbullet\ {\poplogo OnBuch\color{poporange}+}}\par\vspace{2pt}
      {\raggedright\popsemi\fontsize{8}{10.5}\selectfont\color{popline}\MakeUppercase{OnBuch+ teaching team}\\\MakeUppercase{Follows the official MINESEC syllabus}\par}
    \end{minipage}\hfill
    {\poplogo\fontsize{22}{22}\selectfont\color{popcream}OnBuch\color{poporange}+}
  \end{tcolorbox}%
}

% ============================================================
% Page de garde
% #1 titre · #2 matière · #3 niveau · #4 module · #5 n° leçon · #6 durée ·
% #7 description / objectifs (texte ou itemize)
% ============================================================
\newcommand{\pagegarde}[7]{%
  \thispagestyle{empty}
  \newgeometry{a4paper,margin=1.7cm,top=1.6cm,bottom=1.4cm}%
  \begin{tikzpicture}[remember picture,overlay]
    \fill[poporange!18] ([xshift=-3.2cm,yshift=-3.6cm]current page.north east) circle (4.6cm);
    \fill[poppurple!14] ([xshift=2.4cm,yshift=7.6cm]current page.south west) circle (3.8cm);
  \end{tikzpicture}%
  {\poplogo\fontsize{30}{30}\selectfont\color{popink}OnBuch\color{poporange}+}\hfill
  \raisebox{0.6ex}{\poppill{Signed course \textbullet\ Premium}{popink}{popcream}}\par
  \vspace{1pt}\popsquiggle{poppurple}\par
  \vspace{8pt}
  \begin{tcolorbox}[enhanced,colback=poporange,colframe=popink,boxrule=2.8pt,arc=13pt,
      drop shadow=popink,left=18pt,right=18pt,top=12pt,bottom=13pt,after skip=10pt]
    \poppill{#2 \textbullet\ #3}{popink}{popcream}\hspace{5pt}\poppill{#4}{white}{popink}\par\vspace{8pt}
    {\popdisplay\fontsize{14}{14}\selectfont\color{popink}CHAPTER #5}\par\vspace{4pt}
    {\raggedright\hyphenpenalty=10000\exhyphenpenalty=10000\popdisplay\fontsize{25}{27}\selectfont\color{popcream}\MakeUppercase{#1}\par}
    \vspace{8pt}
    {\popxbold\fontsize{9.5}{9.5}\selectfont\color{popink}\MakeUppercase{Suggested time: #6}}
  \end{tcolorbox}
  \begin{tcolorbox}[enhanced,colback=white,colframe=popink,boxrule=2.2pt,arc=10pt,
      drop shadow=popink,left=16pt,right=16pt,top=10pt,bottom=10pt,after skip=8pt]
    \poppill{In this chapter}{poppurple}{white}\par\vspace{5pt}
    \popsemi\fontsize{9.3}{12.6}\selectfont\color{popink}#7
  \end{tcolorbox}
  \noindent\begin{minipage}[t]{0.487\linewidth}
    \begin{tcolorbox}[enhanced,colback=popgreen,colframe=popink,boxrule=2.2pt,arc=10pt,
        drop shadow=popink,left=11pt,right=11pt,top=10pt,bottom=10pt,height=3.1cm,valign=center]
      \tcbox[colback=white,colframe=popink,boxrule=1.3pt,arc=5pt,boxsep=0pt,nobeforeafter,
        left=3pt,right=3pt,top=3pt,bottom=3pt,valign=center]{\qrcode[height=1.9cm]{https://play.google.com/store/apps/details?id=com.luvvix.onbuch&hl=en}}%
      \hspace{8pt}\begin{minipage}[c]{0.5\linewidth}
        {\popdisplay\fontsize{11}{12.5}\selectfont\color{white}\MakeUppercase{Download OnBuch}}\par\vspace{4pt}
        {\raggedright\popsemi\fontsize{8.4}{11}\selectfont\color{white}Free \textbullet\ Google Play\\AI tutor, past papers, results\par}
      \end{minipage}
    \end{tcolorbox}
  \end{minipage}\hfill
  \begin{minipage}[t]{0.487\linewidth}
    \begin{tcolorbox}[enhanced,colback=poppurple,colframe=popink,boxrule=2.2pt,arc=10pt,
        drop shadow=popink,left=11pt,right=11pt,top=10pt,bottom=10pt,height=3.1cm,valign=center]
      \tikz[baseline=-0.7ex]\node[circle,fill=white,draw=popink,line width=1.3pt,minimum size=1.6cm,inner sep=0]{\popdisplay\fontsize{24}{24}\selectfont\color{poppurple}+};%
      \hspace{8pt}\begin{minipage}[c]{0.6\linewidth}
        {\popdisplay\fontsize{11}{12.5}\selectfont\color{white}\MakeUppercase{Go OnBuch+}}\par\vspace{4pt}
        {\raggedright\popsemi\fontsize{8.4}{11}\selectfont\color{white}All signed courses, unlimited AI tutor, PDF sheets — OnBuch+ tab in the app\par}
      \end{minipage}
    \end{tcolorbox}
  \end{minipage}
  \vspace{8pt}
  \popcontenu
  \vfill
  \signatureonbuch
  \restoregeometry
  \newpage
}

% Rangée « Ce cours contient » (page de garde)
\newcommand{\poptile}[3]{% #1 couleur #2 gros texte #3 légende
  \begin{tcolorbox}[enhanced,colback=#1,colframe=popink,boxrule=2pt,arc=9pt,shadow={2.5pt}{-2.5pt}{0pt}{popink},
      left=6pt,right=6pt,top=9pt,bottom=9pt,halign=center,valign=center,height=1.75cm,width=\linewidth,nobeforeafter]
    {\popdisplay\fontsize{12.5}{13}\selectfont\color{white}\MakeUppercase{#2}}\par\vspace{3pt}
    {\centering\popsemi\fontsize{7.8}{9.5}\selectfont\color{white}#3\par}
  \end{tcolorbox}}
\newcommand{\popcontenu}{%
  \par\noindent{\popxboldls\fontsize{7.5}{7.5}\selectfont\color{popmuted}THIS SIGNED COURSE CONTAINS}\par\vspace{7pt}
  \noindent\begin{minipage}[t]{0.235\linewidth}\poptile{popblue}{Course}{complete, illustrated, step by step}\end{minipage}\hfill
  \begin{minipage}[t]{0.235\linewidth}\poptile{poporange}{Methods}{and worked examples}\end{minipage}\hfill
  \begin{minipage}[t]{0.235\linewidth}\poptile{poppink}{Exercises}{exam style + solutions}\end{minipage}\hfill
  \begin{minipage}[t]{0.235\linewidth}\poptile{popgreen}{Summary sheet}{the essentials to remember}\end{minipage}\par}

% Sommaire
\newtcolorbox{sommaire}{enhanced,colback=white,colframe=popink,boxrule=2.2pt,arc=10pt,
  drop shadow=popink,left=18pt,right=18pt,top=13pt,bottom=13pt,before skip=6pt,after skip=20pt,
  before upper={\poppill{Contents}{poppurple}{white}\par\vspace{9pt}\popsemi\fontsize{10.6}{14.5}\selectfont\color{popink}}}

% Signature de fin de cours — poussée tout en bas de la dernière page
\newcommand{\findecours}{%
  \par\vfill\nopagebreak
  \begin{center}\popsquiggle{poporange}\end{center}\vspace{4pt}
  \signatureonbuch
}
\AtBeginDocument{\def\llbracket{\mathopen{[\mkern-3.5mu[}}\def\rrbracket{\mathclose{]\mkern-3.5mu]}}} % intervalles d'entiers (glyphe ⟦ absent de la police)
% Glyphes absents de Fira Math : redéfinis à partir de glyphes disponibles
\AtBeginDocument{%
  \def\setminus{\mathbin{\backslash}}\let\smallsetminus\setminus
  \def\vdots{\mathord{\vbox{\baselineskip3.2pt\lineskiplimit0pt\kern4pt\hbox{.}\hbox{.}\hbox{.}}}}%
  \def\ddots{\mathinner{\mkern1mu\raise6.5pt\hbox{.}\mkern2mu\raise3.6pt\hbox{.}\mkern2mu\raise0.7pt\hbox{.}\mkern1mu}}%
  \def\triangle{\mathord{\scalebox{0.9}{$\symup{\Delta}$}}}%
  \def\nabla{\mathord{\raisebox{\depth}{\scalebox{1}[-1]{$\symup{\Delta}$}}}}%
  \def\checkmark{\popcheck{popdark}}%
  \def\star{\mathbin{\ast}}\def\bigstar{\mathord{\ast}}%
  \def\diamond{\mathbin{\scalebox{0.6}{$\Diamond$}}}%
  \def\bigcup{\mathop{\mathchoice{\raisebox{-0.3em}{\scalebox{1.7}{$\cup$}}}{\scalebox{1.25}{$\cup$}}{\cup}{\cup}}\displaylimits}%
  \def\bigcap{\mathop{\mathchoice{\raisebox{-0.3em}{\scalebox{1.7}{$\cap$}}}{\scalebox{1.25}{$\cap$}}{\cap}{\cap}}\displaylimits}%
  \def\lesseqqgtr{\mathrel{\lessgtr}}\def\bigoplus{\mathop{\oplus}\displaylimits}%
}
\providecommand{\ssi}{if and only if} % abréviation fréquente
\AtBeginDocument{\providecommand{\wideparen}{\overparen}} % arc de cercle

% Fonctions usuelles que les modèles utilisent sans que amsmath les fournisse
\providecommand{\arsinh}{\operatorname{arsinh}}\providecommand{\arcosh}{\operatorname{arcosh}}
\providecommand{\artanh}{\operatorname{artanh}}\providecommand{\arsech}{\operatorname{arsech}}
\providecommand{\arcsch}{\operatorname{arcsch}}\providecommand{\arcoth}{\operatorname{arcoth}}
\providecommand{\arcsinh}{\operatorname{arsinh}}\providecommand{\arccosh}{\operatorname{arcosh}}
\providecommand{\arctanh}{\operatorname{artanh}}\providecommand{\sech}{\operatorname{sech}}
\providecommand{\csch}{\operatorname{csch}}\providecommand{\arcsec}{\operatorname{arcsec}}
\providecommand{\arccsc}{\operatorname{arccsc}}\providecommand{\arccot}{\operatorname{arccot}}
\providecommand{\sgn}{\operatorname{sgn}}\providecommand{\lcm}{\operatorname{lcm}}
\providecommand{\Var}{\operatorname{Var}}\providecommand{\Cov}{\operatorname{Cov}}

%% FIGFIX-BEGIN v3 — correctifs globaux des figures (docs/onbuch-generation/tools/figfix)
\usepackage{adjustbox}\usepackage{etoolbox}
% toute figure TikZ/pgfplots plus large que la ligne est réduite à la largeur disponible
\BeforeBeginEnvironment{tikzpicture}{\begin{adjustbox}{max width=\linewidth}}
\AfterEndEnvironment{tikzpicture}{\end{adjustbox}}
% légendes pgfplots lisibles par-dessus les courbes
\pgfplotsset{every axis/.append style={legend style={fill=white,fill opacity=0.92,text opacity=1,draw=black!15,inner sep=3pt}}}
% étiquettes d'arêtes (syntaxe edge["..."]) avec fond blanc
\tikzset{every edge quotes/.append style={fill=white,inner sep=1.2pt}}
% bulles de cartes mentales : texte à la ligne dans la bulle, bulle un peu plus grande
\tikzset{every concept/.append style={text width=2.0cm,align=center,minimum size=2.3cm,inner sep=2pt}}
%% FIGFIX-END
\begin{document}

\pagegarde{Circle Geometry}{Pure Mathematics with Statistics}{Upper Sixth}{Upper Sixth — Pure mathematics with statistics}{1}{22 periods}{This chapter develops the complete analytic study of the circle in the Cartesian plane: its equations in all standard and general forms, tangents, intersections with lines and other circles, parametric representation, loci and the radical axis. It provides the essential toolkit for solving GCE Advanced Level problems on circle geometry.
\vspace{4pt}
\begin{multicols}{2}\small
\begin{itemize}
\item By the end of this chapter I can write the equation of a circle given its centre and radius, a diameter, or three points on it.
\item By the end of this chapter I can convert the general equation x² + y² + 2gx + 2fy + c = 0 into centre–radius form and identify the centre and radius.
\item By the end of this chapter I can find the equation of a tangent to a circle at a given point or from an external point, and state the condition for a line to touch a circle.
\item By the end of this chapter I can use parametric equations of a circle and convert between parametric and Cartesian forms.
\item By the end of this chapter I can find the points of intersection of a line and a circle, and of two circles.
\item By the end of this chapter I can determine the equation of a circle through the intersection of two circles or of a line and a circle using families of circles.
\end{itemize}\end{multicols}}

\begin{sommaire}
\begin{multicols}{2}
\begin{enumerate}
\item Section 1: Equations of a Circle — Centre at the Origin and General Centre
\item Section 2: The General Equation and Circles on a Given Diameter
\item Section 3: Circle Through Three Points and Circles Under Geometric Conditions
\item Section 4: Tangents to a Circle
\item Section 5: Parametric Equations of a Circle
\item Section 6: Intersection of a Line and a Circle
\item Section 7: Intersection of Two Circles, Families of Circles and the Radical Axis
\item Section 8: Conditions for Two Circles to Touch
\item Section 9: Locus Problems and Synthesis for the GCE Exam
\item Integration activity
\item Methods and worked examples
\item Exercises
\item Solutions to the exercises
\item Summary sheet
\end{enumerate}
\end{multicols}
\end{sommaire}

\begin{objectifs}
\begin{itemize}
\item By the end of this chapter I can write the equation of a circle given its centre and radius, a diameter, or three points on it.
\item By the end of this chapter I can convert the general equation x² + y² + 2gx + 2fy + c = 0 into centre–radius form and identify the centre and radius.
\item By the end of this chapter I can find the equation of a tangent to a circle at a given point or from an external point, and state the condition for a line to touch a circle.
\item By the end of this chapter I can use parametric equations of a circle and convert between parametric and Cartesian forms.
\item By the end of this chapter I can find the points of intersection of a line and a circle, and of two circles.
\item By the end of this chapter I can determine the equation of a circle through the intersection of two circles or of a line and a circle using families of circles.
\item By the end of this chapter I can state and apply the conditions for two circles to touch internally or externally.
\item By the end of this chapter I can find the radical axis of two circles and the length of a tangent from an external point.
\item By the end of this chapter I can solve locus problems and find circles satisfying mixed conditions (touching axes, touching a line, centre on a line).
\end{itemize}
\end{objectifs}

\begin{prerequis}
\begin{itemize}
\item Coordinate geometry of the straight line: distance between two points, midpoint, gradient, equations of lines, perpendicular gradients (Lower Sixth).
\item Pythagoras' theorem and basic circle theorems from O-Level geometry.
\item Solving simultaneous equations (linear and quadratic) and the discriminant of a quadratic.
\item Completing the square for a quadratic expression.
\item Basic trigonometry: sine and cosine of an angle (useful for parametric equations).
\end{itemize}
\end{prerequis}

\newpage

\coursec{Section 1: Equations of a Circle — Centre at the Origin and General Centre}

In Bamenda, a telecommunications engineer must position a relay antenna so that it is exactly the same distance from three villages; on a map, the villages are points and the antenna sits at the centre of a circle through them. In Douala, a roundabout is designed so that a straight avenue just touches its circular kerb — a tangent problem. When a potter in Bafoussam spins clay, every point of the rim traces a circle whose equation we can write down. This chapter gives you the algebraic tools to describe, construct and analyse all these situations: the circle, the simplest curve after the straight line, and a favourite of the GCE Advanced Level examiners.

The central problem of this first section is the following: \emph{how can we translate the geometric idea of a circle into an algebraic equation that we can manipulate?} Once we can do that, questions about centres, radii, and positions of points become routine algebra.

\courssub{The Circle as a Locus}

Think of the potter's wheel in Bafoussam. Fix your attention on one grain of clay at the rim. As the wheel spins, that grain moves, but its distance from the centre of the wheel never changes. The rim is exactly the collection of all positions whose distance from the centre equals the radius of the wheel. This everyday observation is the whole secret of the circle.

\begin{definition}[Circle as a locus]
A \cle{circle} is the \cle{locus} of a point which moves in a plane so that its distance from a fixed point, called the \cle{centre}, is constant. This constant distance is called the \cle{radius} of the circle.
\end{definition}

The word \emph{locus} (plural \emph{loci}) simply means ``the set of all points satisfying a given condition''. So a circle is not a drawn object first — it is a \emph{condition}: ``be at distance $r$ from $C$''. Our task is to rewrite this condition in coordinates. We shall return to locus problems in full generality in Section 9; here the locus condition is the simplest possible one.

\courssub{Circle Centred at the Origin}

Place the centre at the origin $O(0,0)$ of a Cartesian plane, and let $P(x,y)$ be any point on the circle of radius $r$. By the distance formula (itself a direct consequence of Pythagoras' theorem),
\[ OP = \sqrt{(x-0)^2 + (y-0)^2} = \sqrt{x^2+y^2}. \]
The locus condition $OP = r$ gives $\sqrt{x^2+y^2} = r$, and squaring both sides (both sides are non-negative, so this is safe) yields the equation of the circle.

\begin{propriete}[Circle centred at the origin]
The circle with centre $O(0,0)$ and radius $r$ has equation
\[ x^2 + y^2 = r^2. \]
This derivation is examinable: you may be asked to \emph{derive} it from the distance formula.
\end{propriete}

\begin{popfigure}
\begin{center}
\begin{tikzpicture}[scale=1.1]
\draw[->,gray] (-3.2,0) -- (3.2,0) node[right] {$x$};
\draw[->,gray] (0,-3.2) -- (0,3.2) node[above] {$y$};
\draw[popblue,thick] (0,0) circle (2.2);
\coordinate (P) at (1.32,1.76);
\fill[poporange] (P) circle (2pt) node[above right,popdark] {$P(x,y)$};
\draw[popdark,thick] (0,0) -- (P) node[midway,above left] {$r$};
\draw[popgreen,thick] (0,0) -- (1.32,0) node[midway,below] {$x$};
\draw[popgreen,thick] (1.32,0) -- (P) node[midway,right] {$y$};
\draw (1.12,0) -- (1.12,0.2) -- (1.32,0.2);
\fill[popblue] (0,0) circle (2pt) node[below left] {$O$};
\end{tikzpicture}
\end{center}
\legende{The right triangle inside the circle: $x^2 + y^2 = r^2$ is Pythagoras' theorem.}
\end{popfigure}

\begin{exemplebox}[Reading and writing equations centred at the origin]
\begin{enumerate}
\item The circle $x^2 + y^2 = 25$ has centre $(0,0)$ and radius $r = \sqrt{25} = 5$.
\item The circle with centre $(0,0)$ and radius $3$ has equation $x^2 + y^2 = 9$.
\item The circle $x^2 + y^2 = 7$ has radius $\sqrt{7}$: the radius need not be a whole number.
\end{enumerate}
\end{exemplebox}

\begin{attention}[The right-hand side is $r^2$]
In $x^2 + y^2 = r^2$, the number on the right is the \textbf{square} of the radius. For $x^2 + y^2 = 25$ the radius is $5$, not $25$. This is one of the most frequent errors in GCE scripts — always take the square root.
\end{attention}

\courssub{Circle with Centre $C(a,b)$: the Standard Form}

Now suppose the centre is no longer at the origin but at a general point $C(a,b)$. Geometrically, nothing has changed: the circle is just \emph{shifted} $a$ units horizontally and $b$ units vertically. Let $P(x,y)$ be any point on the circle. The locus condition is still $CP = r$, and the distance formula gives
\[ \sqrt{(x-a)^2 + (y-b)^2} = r. \]
Squaring both sides produces the standard form.

\begin{propriete}[Standard equation of a circle]
The circle with centre $C(a,b)$ and radius $r$ has equation
\[ (x-a)^2 + (y-b)^2 = r^2. \]
Conversely, any equation of this form represents a circle with centre $(a,b)$ and radius $\sqrt{r^2}$. The derivation from the distance formula is examinable.
\end{propriete}

\begin{popfigure}
\begin{center}
\begin{tikzpicture}[scale=0.85]
\draw[->,gray] (-1.5,0) -- (7.5,0) node[right] {$x$};
\draw[->,gray] (0,-1.5) -- (0,6) node[above] {$y$};
\coordinate (C) at (4,3);
\draw[popblue,thick] (C) circle (1.8);
\fill[popblue] (C) circle (2pt) node[above,popdark] {$C(a,b)$};
\coordinate (P) at (5.27,4.44);
\fill[poporange] (P) circle (2pt) node[above right,popdark] {$P(x,y)$};
\draw[popdark,thick] (C) -- (P) node[midway,above left] {$r$};
\draw[dashed,popgreen] (0,3) -- (C) node[midway,above] {$a$};
\draw[dashed,popgreen] (4,0) -- (C) node[midway,right] {$b$};
\draw[dashed,gray] (0,0) -- (C);
\fill[black] (0,0) circle (1.5pt) node[below left] {$O$};
\draw[popgreen,thick] (C) -- (5.27,3) node[midway,below] {$x-a$};
\draw[popgreen,thick] (5.27,3) -- (P) node[midway,right] {$y-b$};
\draw (5.07,3) -- (5.07,3.2) -- (5.27,3.2);
\end{tikzpicture}
\end{center}
\legende{Shifting the centre to $C(a,b)$: the horizontal leg has length $x-a$, the vertical leg $y-b$.}
\end{popfigure}

\begin{attention}[Watch the signs of the centre]
In $(x-a)^2 + (y-b)^2 = r^2$ the centre is $(a,b)$, \textbf{not} $(-a,-b)$. The formula contains \emph{minus} signs, so:
\begin{itemize}
\item $(x-3)^2 + (y-2)^2 = 16$ has centre $(3,2)$;
\item $(x+3)^2 + (y-1)^2 = 16$, i.e. $(x-(-3))^2 + (y-1)^2 = 16$, has centre $(-3,1)$.
\end{itemize}
Rewrite any ``$+$'' as ``$-(-\,\cdot\,)$'' before reading off the coordinates.
\end{attention}

\begin{exemplebox}[Reading centre and radius from the standard form]
\begin{enumerate}
\item $(x-2)^2 + (y+5)^2 = 9$: centre $(2,-5)$, radius $3$.
\item $(x+1)^2 + (y+4)^2 = 20$: centre $(-1,-4)$, radius $\sqrt{20} = 2\sqrt{5}$.
\item $x^2 + (y-3)^2 = 4$: this is $(x-0)^2 + (y-3)^2 = 4$, centre $(0,3)$ on the $y$-axis, radius $2$.
\end{enumerate}
\end{exemplebox}

\begin{methode}[Sketching a circle from its equation]
\begin{enumerate}
\item Read the centre $C(a,b)$ and compute $r = \sqrt{r^2}$.
\item Plot $C$ in the plane.
\item Mark the four points $(a\pm r, b)$ and $(a, b\pm r)$ — the extreme points of the circle.
\item Draw a smooth circle through these four points and label the centre and radius.
\end{enumerate}
\end{methode}

\courssub{Circle Given its Centre and a Point on it}

Suppose we know the centre $C$ and one point $P$ lying \emph{on} the circle. Then the radius is simply the distance $CP$: the point on the circle ``tells us'' how big the circle is. Since the equation only needs $r^2$, we compute $r^2 = CP^2$ directly — no square root is ever needed.

\begin{methode}[Equation given the centre and a point on the circle]
\begin{enumerate}
\item Compute $r^2 = CP^2 = (x_P - a)^2 + (y_P - b)^2$.
\item Substitute into the standard form $(x-a)^2 + (y-b)^2 = r^2$.
\item Expand only if the question asks for the general form.
\end{enumerate}
\end{methode}

\begin{exoresolu}[GCE-style: centre and a point on the circle]
Find the equation of the circle with centre $C(2,-1)$ which passes through the point $P(5,3)$.
\tcblower
\textbf{Solution.} The radius squared is the squared distance $CP$:
\[ r^2 = (5-2)^2 + (3-(-1))^2 = 3^2 + 4^2 = 9 + 16 = 25. \]
Substituting into the standard form with centre $(2,-1)$:
\[ (x-2)^2 + (y+1)^2 = 25. \]
\emph{Check:} does $P(5,3)$ satisfy it? $(5-2)^2 + (3+1)^2 = 9 + 16 = 25$. \checkmark\ The equation is correct.
\end{exoresolu}

\begin{exoresolu}[Circle centred on an axis, passing through the origin]
A circle has its centre on the positive $x$-axis and passes through the points $O(0,0)$ and $A(6,0)$. Find its equation.
\tcblower
\textbf{Solution.} Since the centre lies on the $x$-axis, write $C(a,0)$. The centre is equidistant from $O$ and $A$:
\[ CO^2 = a^2, \qquad CA^2 = (6-a)^2. \]
Setting $CO^2 = CA^2$: $a^2 = (6-a)^2 = 36 - 12a + a^2$, so $12a = 36$ and $a = 3$. Then $r^2 = CO^2 = 9$, and the circle is
\[ (x-3)^2 + y^2 = 9. \]
Note that the origin lies on the circle: $(0-3)^2 + 0^2 = 9$. \checkmark
\end{exoresolu}

\begin{attention}[Circles through the origin]
A circle passes through the origin if and only if $r^2 = a^2 + b^2$ (substitute $(0,0)$ into the standard form). For example $(x-3)^2 + (y-4)^2 = 25$ passes through the origin since $9 + 16 = 25$. Recognising this saves time in the exam.
\end{attention}

\courssub{Position of a Point Relative to a Circle}

The locus condition $CP = r$ splits the plane into three natural regions. For a point $P$ and a circle of centre $C$ and radius $r$:
\begin{itemize}
\item $CP < r$: $P$ lies \cle{inside} the circle;
\item $CP = r$: $P$ lies \cle{on} the circle;
\item $CP > r$: $P$ lies \cle{outside} the circle.
\end{itemize}
In practice we compare $CP^2$ with $r^2$ to avoid square roots: substitute the coordinates of $P$ into the left-hand side $(x-a)^2 + (y-b)^2$ and compare the result with the right-hand side.

\begin{popfigure}
\begin{center}
\begin{tikzpicture}
\begin{axis}[
  width=0.72\linewidth, axis equal image,
  axis lines=middle, xlabel={$x$}, ylabel={$y$},
  xmin=-4, xmax=8, ymin=-4, ymax=8,
  xtick={-2,0,2,4,6}, ytick={-2,0,2,4,6},
  legend style={at={(0.02,0.98)}, anchor=north west, font=\small}
]
\addplot[popblue, thick, domain=0:360, samples=100] ({2+3*cos(x)},{2+3*sin(x)});
\addlegendentry{$(x-2)^2+(y-2)^2=9$}
\addplot[only marks, mark=*, mark size=2pt, popgreen] coordinates {(3,3)};
\addlegendentry{$A(3,3)$: inside}
\addplot[only marks, mark=*, mark size=2pt, poporange] coordinates {(5,2)};
\addlegendentry{$B(5,2)$: on}
\addplot[only marks, mark=*, mark size=2pt, poppink] coordinates {(6,6)};
\addlegendentry{$D(6,6)$: outside}
\addplot[only marks, mark=*, mark size=1.5pt, popdark] coordinates {(2,2)};
\node[font=\small, popdark] at (axis cs:1.2,1.4) {$C$};
\end{axis}
\end{tikzpicture}
\end{center}
\legende{The distance test: $CA^2 = 2 < 9$, $CB^2 = 9$, $CD^2 = 32 > 9$.}
\end{popfigure}

\begin{exoresolu}[Inside, on, or outside?]
Determine the position of each point relative to the circle $(x-2)^2 + (y-2)^2 = 9$:
$A(3,3)$, $B(5,2)$, $D(6,6)$.
\tcblower
\textbf{Solution.} Here $C(2,2)$ and $r^2 = 9$. We compute the squared distance from $C$ for each point.
\begin{align*}
A(3,3):& \quad (3-2)^2 + (3-2)^2 = 1 + 1 = 2 < 9 \quad\Rightarrow\quad A \text{ is inside.}\\
B(5,2):& \quad (5-2)^2 + (2-2)^2 = 9 + 0 = 9 = r^2 \quad\Rightarrow\quad B \text{ is on the circle.}\\
D(6,6):& \quad (6-2)^2 + (6-2)^2 = 16 + 16 = 32 > 9 \quad\Rightarrow\quad D \text{ is outside.}
\end{align*}
\end{exoresolu}

\begin{exemplebox}[A Cameroonian application]
A transmitter at the point $T(1,2)$ (coordinates in kilometres on a map) covers every receiver within \SI{5}{km}. Is a receiver at $R(5,6)$ covered? The coverage zone is the circle $(x-1)^2 + (y-2)^2 = 25$. Compute
\[ TR^2 = (5-1)^2 + (6-2)^2 = 16 + 16 = 32 > 25. \]
Since $TR > 5$, the receiver lies \emph{outside} the coverage zone: it is not covered.
\end{exemplebox}

\courssub{Consolidation}

\begin{aretenir}[Key points of Section 1]
\begin{itemize}
\item A circle is the locus of points at a fixed distance $r$ (radius) from a fixed point $C$ (centre).
\item Centre at the origin: $x^2 + y^2 = r^2$.
\item Centre $C(a,b)$: $(x-a)^2 + (y-b)^2 = r^2$ — the \textbf{standard form}.
\item Given the centre and a point $P$ on the circle: $r^2 = CP^2$.
\item Position of a point $P$: compare $CP^2$ with $r^2$ (inside, on, outside).
\item The circle passes through the origin exactly when $a^2 + b^2 = r^2$.
\end{itemize}
\end{aretenir}

\begin{savaistu}[Why the GCE loves the circle]
The equation $(x-a)^2 + (y-b)^2 = r^2$ is nothing more than Pythagoras' theorem wearing coordinate clothing. Nearly every question in this chapter — tangents, intersections, loci — ultimately reduces to a right-angled triangle and the distance formula. Master this section and the rest of the chapter becomes a sequence of variations on one idea.
\end{savaistu}

\begin{exercice}[Practice]
\begin{enumerate}
\item Write down the centre and radius of each circle:
\begin{enumerate}
\item $x^2 + y^2 = 49$;
\item $(x-4)^2 + (y+1)^2 = 36$;
\item $(x+2)^2 + (y+7)^2 = 12$.
\end{enumerate}
\item Find the equation of the circle with centre $(-3,2)$ passing through $(1,5)$.
\item Find the equation of the circle with centre on the $y$-axis, passing through $(0,0)$ and $(0,8)$.
\item Determine whether $(4,-2)$ lies inside, on or outside the circle $(x-1)^2 + (y+2)^2 = 10$.
\end{enumerate}
\end{exercice}

\begin{corrige}[Exercise 1]
\begin{enumerate}
\item (a) Centre $(0,0)$, $r=7$. (b) Centre $(4,-1)$, $r=6$. (c) Centre $(-2,-7)$, $r=\sqrt{12}=2\sqrt{3}$.
\item $r^2 = (1+3)^2 + (5-2)^2 = 16 + 9 = 25$, so $(x+3)^2 + (y-2)^2 = 25$.
\item Centre $(0,b)$ with $b^2 = (8-b)^2$, giving $b = 4$ and $r^2 = 16$: $x^2 + (y-4)^2 = 16$.
\item $(4-1)^2 + (-2+2)^2 = 9 < 10$: the point lies \textbf{inside} the circle.
\end{enumerate}
\end{corrige}

\coursec{Section 2: The General Equation and Circles on a Given Diameter}

In Section 1 we established the \cle{centre--radius form} of a circle: a circle with centre $C(a,b)$ and radius $r$ has equation
\[ (x-a)^2 + (y-b)^2 = r^2. \]
This form is perfect when the centre and radius are known. But in many GCE questions the equation of a circle is handed to you \emph{already expanded}, with all brackets removed, and you must work \emph{backwards} to recover the centre and radius. In other problems you are told that a segment $AB$ is a \emph{diameter} of the circle, and you must build the equation from the two endpoints. This section develops both techniques: the \cle{general equation} of a circle and the \cle{diameter form}.

\courssub{From Centre--Radius Form to the General Equation}

\textbf{An observation to start with.}\par
Take the circle with centre $(2,-3)$ and radius $5$. Its equation is
\[ (x-2)^2 + (y+3)^2 = 25. \]
Expand the brackets:
\[ x^2 - 4x + 4 + y^2 + 6y + 9 = 25, \]
and move everything to one side:
\[ x^2 + y^2 - 4x + 6y - 12 = 0. \]
Look carefully at the result. The $x^2$ and $y^2$ terms both have coefficient $1$; there is a term in $x$, a term in $y$, and a constant. There is \emph{no} $xy$ term. This shape is completely general: expanding $(x-a)^2+(y-b)^2=r^2$ always produces an equation of the form
\[ x^2 + y^2 + 2gx + 2fy + c = 0, \]
where $g = -a$, $f = -b$ and $c = a^2 + b^2 - r^2$. (The factors $2$ in $2g$ and $2f$ are a convenient convention: they make the completion of the square come out cleanly.)

\begin{definition}[General equation of a circle]
The \cle{general equation} of a circle is
\[ x^2 + y^2 + 2gx + 2fy + c = 0, \]
where $g$, $f$ and $c$ are constants. The circle has
\[ \text{centre } (-g,\,-f) \qquad \text{and} \qquad \text{radius } r = \sqrt{g^2 + f^2 - c}. \]
\end{definition}

\begin{attention}[Equal coefficients of $x^2$ and $y^2$]
For an equation of the second degree to represent a \emph{circle}, the coefficients of $x^2$ and $y^2$ must be \textbf{equal} (and there must be no $xy$ term). If you are given, say, $3x^2 + 3y^2 + 6x - 12y + 4 = 0$, \textbf{divide through by $3$ first} before reading off $g$, $f$ and $c$. An equation such as $2x^2 + 5y^2 = 20$ is \emph{not} a circle (it is an ellipse), because the coefficients of $x^2$ and $y^2$ differ.
\end{attention}

\courssub{Recovering Centre and Radius: Completing the Square}

Why does the centre come out as $(-g,-f)$? The cleanest way to see it is to reverse the expansion, that is, to \emph{complete the square} in $x$ and in $y$. Starting from
\[ x^2 + y^2 + 2gx + 2fy + c = 0, \]
group the $x$-terms and the $y$-terms:
\[ \bigl(x^2 + 2gx\bigr) + \bigl(y^2 + 2fy\bigr) = -c. \]
Recall the identity $X^2 + 2kX = (X+k)^2 - k^2$. Applying it to each group:
\[ (x+g)^2 - g^2 + (y+f)^2 - f^2 = -c, \]
\[ (x+g)^2 + (y+f)^2 = g^2 + f^2 - c. \]
Comparing with the centre--radius form $(x-a)^2+(y-b)^2=r^2$, we read off at once:
\[ a = -g, \qquad b = -f, \qquad r^2 = g^2 + f^2 - c. \]
This derivation is short, elegant and fully examinable — a GCE examiner may well ask you to \emph{show} that a given equation represents a circle and to state its centre and radius.

\begin{methode}[Centre and radius from the general form]
Given $x^2 + y^2 + 2gx + 2fy + c = 0$:
\begin{enumerate}
\item Check that the coefficients of $x^2$ and $y^2$ are equal; if they are not $1$, divide the whole equation by that common coefficient.
\item \textbf{Quick route:} identify $g$ and $f$ as \emph{half} the coefficients of $x$ and $y$; then centre $=(-g,-f)$ and $r=\sqrt{g^2+f^2-c}$.
\item \textbf{Safe route (recommended in the exam):} complete the square in $x$ and in $y$ to rewrite the equation as $(x+g)^2+(y+f)^2 = g^2+f^2-c$, then read off the centre and radius directly.
\end{enumerate}
\end{methode}

\begin{exemplebox}[Reading off the centre and radius]
Find the centre and radius of the circle $x^2 + y^2 - 6x + 8y + 9 = 0$.\par
\textbf{Quick route:} $2g = -6 \Rightarrow g=-3$; $2f = 8 \Rightarrow f=4$; $c = 9$.
\[ \text{Centre } = (-g,-f) = (3,-4), \qquad r = \sqrt{(-3)^2 + 4^2 - 9} = \sqrt{9+16-9} = \sqrt{16} = 4. \]
\textbf{Check by completing the square:}
\[ (x^2 - 6x) + (y^2 + 8y) = -9 \;\Longrightarrow\; (x-3)^2 - 9 + (y+4)^2 - 16 = -9, \]
\[ (x-3)^2 + (y+4)^2 = 16. \quad\checkmark \]
\end{exemplebox}

\begin{popfigure}
\begin{tikzpicture}[font=\small, >=stealth]
\node[draw=popblue, rounded corners, fill=popblueL, text width=3.6cm, align=center] (gen) at (0,0) {\textbf{General form}\\[2pt] $x^2+y^2+2gx$\\ $+\,2fy+c=0$};
\node[draw=poporange, rounded corners, fill=white, text width=3.4cm, align=center] (grp) at (5.6,0) {\textbf{Group \& complete}\\[2pt] $(x+g)^2 - g^2$\\ $+\,(y+f)^2 - f^2 = -c$};
\node[draw=poppurple, rounded corners, fill=white, text width=3.6cm, align=center] (std) at (11.2,0) {\textbf{Standard form}\\[2pt] $(x+g)^2+(y+f)^2$\\ $= g^2+f^2-c$};
\node[draw=popgreen, rounded corners, fill=popgreenL, text width=3.4cm, align=center] (read) at (11.2,-2.6) {\textbf{Read off}\\[2pt] centre $(-g,-f)$\\ $r=\sqrt{g^2+f^2-c}$};
\draw[->, very thick, popblue] (gen) -- (grp);
\draw[->, very thick, poporange] (grp) -- (std);
\draw[->, very thick, poppurple] (std) -- (read);
\end{tikzpicture}
\legende{Converting the general equation to centre--radius form by completing the square.}
\end{popfigure}

\courssub{When Does the General Equation Represent a Real Circle?}

The formula $r^2 = g^2 + f^2 - c$ hides a subtlety that the GCE Board loves to test: the quantity under the square root may be positive, zero or negative.

\begin{propriete}[Existence condition — proof examinable]
The equation $x^2 + y^2 + 2gx + 2fy + c = 0$ represents:
\begin{itemize}
\item a \textbf{real circle} of radius $\sqrt{g^2+f^2-c}$ if $g^2 + f^2 - c > 0$;
\item a \textbf{point circle} (the single point $(-g,-f)$) if $g^2 + f^2 - c = 0$;
\item \textbf{no real locus at all} if $g^2 + f^2 - c < 0$.
\end{itemize}
\emph{Proof.} Completing the square gives $(x+g)^2+(y+f)^2 = g^2+f^2-c$. The left-hand side, being a sum of two squares of real numbers, is always $\geq 0$. If the right-hand side is negative, no real point $(x,y)$ can satisfy the equation. If it is zero, we need both squares to vanish, forcing $x=-g$ and $y=-f$: a single point. If it is positive, we obtain a genuine circle. $\blacksquare$
\end{propriete}

\begin{exemplebox}[Testing the three cases]
Examine the nature of each locus:
\begin{itemize}
\item $x^2+y^2-2x+4y-4=0$: $g^2+f^2-c = 1+4+4 = 9 > 0$: real circle, centre $(1,-2)$, radius $3$.
\item $x^2+y^2+6x-2y+10=0$: $g^2+f^2-c = 9+1-10 = 0$: the single point $(-3,1)$.
\item $x^2+y^2+2x+2y+5=0$: $g^2+f^2-c = 1+1-5 = -3 < 0$: no real locus.
\end{itemize}
\end{exemplebox}

\begin{exoresolu}[Finding $k$ for a real circle]
Find the range of values of $k$ for which the equation
\[ x^2 + y^2 + 4x - 6y + k = 0 \]
represents a real circle, and state what happens at the boundary value.
\tcblower
Here $2g = 4$ so $g=2$; $2f=-6$ so $f=-3$; $c = k$. The existence condition is
\[ g^2 + f^2 - c > 0 \;\Longleftrightarrow\; 4 + 9 - k > 0 \;\Longleftrightarrow\; k < 13. \]
Hence the equation represents a real circle for all $k < 13$, with centre $(-2,3)$ and radius $\sqrt{13-k}$. When $k = 13$ the radius is $0$: the locus collapses to the single point $(-2,3)$ (a point circle). For $k > 13$ there is no real locus.
\end{exoresolu}

\courssub{Circle on a Given Diameter}

Suppose now that instead of the centre you are given the two endpoints of a \emph{diameter}, $A(x_1,y_1)$ and $B(x_2,y_2)$. Two facts from Section 1 and from O-Level geometry combine beautifully here:
\begin{itemize}
\item the centre of the circle is the \textbf{midpoint} $M$ of $AB$;
\item the radius is \textbf{half the length} of $AB$.
\end{itemize}

\begin{methode}[Circle with diameter $AB$ — midpoint route]
\begin{enumerate}
\item Compute the centre $M = \left(\dfrac{x_1+x_2}{2}, \dfrac{y_1+y_2}{2}\right)$.
\item Compute the radius $r = \tfrac{1}{2}|AB| = \tfrac{1}{2}\sqrt{(x_2-x_1)^2+(y_2-y_1)^2}$ (or find $|MA|$ directly).
\item Write the equation $(x - h)^2 + (y - k)^2 = r^2$ where $M=(h,k)$, and expand if the general form is required.
\end{enumerate}
\end{methode}

There is, however, a faster and very elegant route, based on a classical theorem: the \textbf{angle in a semicircle is a right angle}. If $P(x,y)$ is any point of the circle with diameter $AB$, then $\angle APB = 90^\circ$, so the vectors $\vec{PA}$ and $\vec{PB}$ are perpendicular, and their scalar product is zero.

\begin{propriete}[Diameter form of the equation of a circle — proof examinable]
The circle having $A(x_1,y_1)$ and $B(x_2,y_2)$ as endpoints of a diameter has equation
\[ (x - x_1)(x - x_2) + (y - y_1)(y - y_2) = 0. \]
\emph{Proof.} Let $P(x,y)$ be a point of the circle. Then
\[ \vec{PA} = \begin{pmatrix} x_1 - x \\ y_1 - y \end{pmatrix}, \qquad \vec{PB} = \begin{pmatrix} x_2 - x \\ y_2 - y \end{pmatrix}. \]
Since the angle in a semicircle is $90^\circ$, $\vec{PA} \perp \vec{PB}$, so $\vec{PA}\cdot\vec{PB} = 0$:
\[ (x_1-x)(x_2-x) + (y_1-y)(y_2-y) = 0, \]
which is equivalent to $(x-x_1)(x-x_2) + (y-y_1)(y-y_2) = 0$ (each bracket is multiplied twice by $-1$). Conversely, any point satisfying this equation either coincides with $A$ or $B$, or satisfies $\vec{PA}\cdot\vec{PB}=0$, hence lies on the circle with diameter $AB$. $\blacksquare$
\end{propriete}

\begin{popfigure}
\begin{tikzpicture}[scale=1.15, font=\small]
\coordinate (M) at (0,0);
\coordinate (A) at (-2,0);
\coordinate (B) at (2,0);
\coordinate (P) at (1.2,1.6);
\draw[thick, popblue] (M) circle (2);
\draw[thick, popdark] (A) -- (B);
\draw[thick, poporange] (A) -- (P) -- (B);
\draw[popmuted, dashed] (M) -- (P);
\draw[popgreen] (1.2,1.6) ++(-0.55,-0.18) -- ++(-0.14,-0.42) -- ++(0.42,-0.14);
\fill[popink] (A) circle (1.6pt) node[below left] {$A(x_1,y_1)$};
\fill[popink] (B) circle (1.6pt) node[below right] {$B(x_2,y_2)$};
\fill[popink] (M) circle (1.6pt) node[below] {$M$};
\fill[poporange] (P) circle (1.8pt) node[above right] {$P(x,y)$};
\node[popgreen] at (0.62,1.05) {$90^\circ$};
\node[popmuted] at (-1.15,0.95) {$r$};
\end{tikzpicture}
\legende{Circle with diameter $AB$: the centre $M$ is the midpoint of $AB$, and for every point $P$ on the circle, $\angle APB = 90^\circ$.}
\end{popfigure}

\begin{exoresolu}[Circle on a given diameter — both routes]
Find the equation of the circle having the segment joining $A(1,2)$ and $B(5,-4)$ as a diameter.
\tcblower
\textbf{Route 1 (midpoint).} The centre is
\[ M = \left(\frac{1+5}{2}, \frac{2+(-4)}{2}\right) = (3,-1). \]
The radius is half of $|AB|$:
\[ |AB| = \sqrt{(5-1)^2 + (-4-2)^2} = \sqrt{16+36} = \sqrt{52}, \qquad r^2 = \frac{52}{4} = 13. \]
Hence $(x-3)^2 + (y+1)^2 = 13$, which expands to $x^2 + y^2 - 6x + 2y - 3 = 0$.

\textbf{Route 2 (diameter form — the GCE shortcut).} Apply $(x-x_1)(x-x_2)+(y-y_1)(y-y_2)=0$:
\[ (x-1)(x-5) + (y-2)(y+4) = 0, \]
\[ x^2 - 6x + 5 + y^2 + 2y - 8 = 0 \;\Longrightarrow\; x^2 + y^2 - 6x + 2y - 3 = 0. \quad\checkmark \]
Both routes agree. Notice how the diameter form produces the \emph{general equation} in a single line — a valuable time-saver in the examination room.
\end{exoresolu}

\begin{attention}[Sign traps with the diameter form]
Two errors recur every year:
\begin{itemize}
\item Writing $(x-x_1)(y-y_2)+(y-y_1)(x-x_2)=0$: the $x$'s stay together and the $y$'s stay together.
\item Forgetting that a negative coordinate changes the sign inside the bracket: with $B(5,-4)$ the factor is $(y-(-4)) = (y+4)$, not $(y-4)$.
\end{itemize}
Also remember: the diameter form gives the circle for which $AB$ is a \emph{diameter}, not merely a chord. If the problem says ``$AB$ is a chord'', this formula does not apply.
\end{attention}

\courssub{Mixed Techniques: A Consolidation Example}

Many GCE questions force you to combine everything seen so far: expand, complete the square, test the existence condition, and use a diameter.

\begin{exoresolu}[Synthesis]
A circle $C_1$ has equation $x^2 + y^2 + 2x - 8y + 8 = 0$.
\begin{enumerate}
\item Find the centre and radius of $C_1$.
\item A second circle $C_2$ has the line joining $P(-1,3)$ and $Q(3,5)$ as a diameter. Find the centre and radius of $C_2$, and determine whether $C_1$ and $C_2$ are concentric.
\end{enumerate}
\tcblower
\textbf{1.} Complete the square:
\[ (x^2+2x) + (y^2-8y) = -8 \;\Longrightarrow\; (x+1)^2 - 1 + (y-4)^2 - 16 = -8, \]
\[ (x+1)^2 + (y-4)^2 = 9. \]
Centre $(-1,4)$, radius $3$. (Check: $g^2+f^2-c = 1+16-8 = 9 > 0$, so a real circle indeed.)

\textbf{2.} The centre of $C_2$ is the midpoint of $PQ$:
\[ M = \left(\frac{-1+3}{2}, \frac{3+5}{2}\right) = (1,4). \]
This is \emph{not} the same centre as $C_1$, which is $(-1,4)$ — so the two circles are \textbf{not} concentric. The radius of $C_2$ is
\[ r = \tfrac{1}{2}|PQ| = \tfrac{1}{2}\sqrt{(3-(-1))^2 + (5-3)^2} = \tfrac{1}{2}\sqrt{16+4} = \tfrac{1}{2}\sqrt{20} = \sqrt{5}. \]
Using the diameter form directly:
\[ (x+1)(x-3) + (y-3)(y-5) = 0 \;\Longrightarrow\; x^2 - 2x - 3 + y^2 - 8y + 15 = 0, \]
\[ x^2 + y^2 - 2x - 8y + 12 = 0. \]
\end{exoresolu}

\begin{exercice}[Practice for the GCE]
\begin{enumerate}
\item Find the centre and radius of each circle:
\begin{itemize}
\item $x^2 + y^2 - 10x + 4y + 20 = 0$;
\item $2x^2 + 2y^2 + 6x - 10y - 1 = 0$ (divide through first!).
\end{itemize}
\item Determine whether $x^2 + y^2 - 4x + 6y + 14 = 0$ represents a real circle, a point, or no real locus.
\item Find the values of $\lambda$ for which $x^2+y^2+8x-2y+\lambda=0$ is a real circle.
\item Find the equation of the circle with diameter endpoints $A(-2,1)$ and $B(4,7)$, giving your answer in general form.
\item A circle has equation $x^2+y^2-6x+2y+6=0$. Show that the point $(5,-3)$ lies on it, and find the equation of the circle having the segment from $(5,-3)$ to the centre as diameter.
\end{enumerate}
\end{exercice}

\begin{corrige}[Exercise 1]
\textbf{1.} (a) $g=-5,\ f=2,\ c=20$: centre $(5,-2)$, $r=\sqrt{25+4-20}=\sqrt{9}=3$. (b) Divide by $2$: $x^2+y^2+3x-5y-\tfrac12=0$; $g=\tfrac32,\ f=-\tfrac52$: centre $\left(-\tfrac32,\tfrac52\right)$, $r=\sqrt{\tfrac94+\tfrac{25}{4}+\tfrac12}=\sqrt{\tfrac{36}{4}}=3$.\par
\textbf{2.} $g^2+f^2-c = 4+9-14 = -1 < 0$: no real locus.\par
\textbf{3.} $16+1-\lambda > 0 \Rightarrow \lambda < 17$.\par
\textbf{4.} $(x+2)(x-4)+(y-1)(y-7)=0 \Rightarrow x^2+y^2-2x-8y-1=0$.\par
\textbf{5.} Centre $(3,-1)$, radius $2$; $(5)^2+(-3)^2-6(5)+2(-3)+6 = 25+9-30-6+6 = 4 = r^2$: yes, $(5,-3)$ lies on the circle. Diameter endpoints $(5,-3)$ and $(3,-1)$: $(x-5)(x-3)+(y+3)(y+1)=0 \Rightarrow x^2+y^2-8x+4y+18=0$.
\end{corrige}

\begin{aretenir}[Key points of Section 2]
\begin{itemize}
\item General equation: $x^2+y^2+2gx+2fy+c=0$, with \textbf{centre} $(-g,-f)$ and \textbf{radius} $\sqrt{g^2+f^2-c}$.
\item Real circle requires $g^2+f^2-c>0$; equality gives a point circle; a negative value gives no real locus.
\item Complete the square in $x$ and in $y$ to pass from the general form to the centre--radius form — and always divide through first if the coefficients of $x^2$ and $y^2$ are not $1$.
\item Circle with diameter $A(x_1,y_1)$, $B(x_2,y_2)$: centre $=$ midpoint of $AB$, radius $=\tfrac12|AB|$, or use the shortcut $(x-x_1)(x-x_2)+(y-y_1)(y-y_2)=0$, justified by the right angle in a semicircle.
\end{itemize}
\end{aretenir}

With the general equation now mastered, we are ready for Section 3, where we shall determine the circle passing through \emph{three given points} and circles subject to other geometric conditions — problems in which the general form $x^2+y^2+2gx+2fy+c=0$, with its three unknown constants, becomes our most powerful tool.

\coursec{Section 3: Circle Through Three Points and Circles Under Geometric Conditions}

In Section 2 we learnt to recognise and manipulate the \cle{general equation} of a circle, $x^{2}+y^{2}+2gx+2fy+c=0$, whose centre is $(-g,-f)$ and whose radius is $\sqrt{g^{2}+f^{2}-c}$ (provided $g^{2}+f^{2}-c>0$). We also saw that a circle is completely determined once we know its diameter. But in many GCE problems the data are given differently: we are told \emph{points the circle must pass through}, or \emph{lines it must touch}, or \emph{a line on which its centre must lie}. This section develops the algebraic and geometric machinery for all such situations. The central idea is simple: \textbf{three independent conditions determine a circle}, because the general equation contains exactly three unknown constants $g$, $f$ and $c$.

\courssub{Circle Through Three Non-Collinear Points}

\begin{experience}[How many points fix a circle?]
Take a sheet of paper. Through \emph{one} given point you can draw infinitely many circles. Through \emph{two} given points you can still draw infinitely many circles (their centres all lie on the perpendicular bisector of the segment joining the points). But through \emph{three} non-collinear points --- think of three towns on a map of the Northwest Region --- exactly \emph{one} circle passes. If the three points are collinear, no circle passes through them at all. This is why three points are the natural data for determining a circle.
\end{experience}

\begin{propriete}[Circle through three points]
Through any three \cle{non-collinear} points there passes exactly one circle. Its equation is found by substituting the coordinates of the three points into the general equation
\[ x^{2}+y^{2}+2gx+2fy+c=0 \]
and solving the resulting system of three linear equations in the three unknowns $g$, $f$ and $c$. (The algebraic technique is examinable; it is simply simultaneous equations.)
\end{propriete}

\begin{methode}[Circle through three given points]
\begin{enumerate}
\item Write the general equation $x^{2}+y^{2}+2gx+2fy+c=0$.
\item Substitute each of the three points in turn to obtain three linear equations in $g$, $f$, $c$.
\item Solve the system (elimination is usually fastest).
\item Write the equation, and if required state the centre $(-g,-f)$ and radius $\sqrt{g^{2}+f^{2}-c}$.
\item \textbf{Check:} verify that each given point satisfies the final equation.
\end{enumerate}
\end{methode}

\begin{exoresolu}[Circle through three points]
Find the equation of the circle passing through the points $A(1,2)$, $B(3,-4)$ and $C(5,-6)$. State its centre and radius.
\tcblower
Substitute each point into $x^{2}+y^{2}+2gx+2fy+c=0$.

Through $A(1,2)$: $\;1+4+2g+4f+c=0 \;\Rightarrow\; 2g+4f+c=-5 \quad (1)$

Through $B(3,-4)$: $\;9+16+6g-8f+c=0 \;\Rightarrow\; 6g-8f+c=-25 \quad (2)$

Through $C(5,-6)$: $\;25+36+10g-12f+c=0 \;\Rightarrow\; 10g-12f+c=-61 \quad (3)$

Subtract $(1)$ from $(2)$: $\;4g-12f=-20 \;\Rightarrow\; g-3f=-5 \quad (4)$

Subtract $(2)$ from $(3)$: $\;4g-4f=-36 \;\Rightarrow\; g-f=-9 \quad (5)$

Subtracting $(5)$ from $(4)$: $-2f=4$, so $f=-2$, and then $g=-9+f=-11$.

From $(1)$: $c=-5-2g-4f=-5+22+8=25$.

Hence the circle is
\[ x^{2}+y^{2}-22x-4y+25=0, \]
with centre $(11,2)$ and radius $\sqrt{121+4-25}=\sqrt{100}=10$.

\emph{Check:} for $A(1,2)$: $1+4-22-8+25=0$ \checkmark; for $B(3,-4)$: $9+16-66+16+25=0$ \checkmark; for $C(5,-6)$: $25+36-110+24+25=0$ \checkmark.
\end{exoresolu}

\courssub{The Geometric Alternative: Perpendicular Bisectors}

There is a beautiful geometric reason why three points determine a circle, and it gives a second method.

\begin{propriete}[Centre on the perpendicular bisector of a chord]
The centre of a circle lies on the \cle{perpendicular bisector} of every chord of the circle. Consequently, the centre of the circle through three non-collinear points $A$, $B$, $C$ is the intersection of the perpendicular bisectors of any two of the segments $AB$, $BC$, $CA$. This point is called the \cle{circumcentre} of triangle $ABC$.
\end{propriete}

\textbf{Justification.} Let $O$ be the centre and $PQ$ a chord with midpoint $M$. Then $OP=OQ$ (radii), so triangles $OMP$ and $OMQ$ are congruent (SSS, with $OM$ common and $PM=MQ$). Hence $\angle OMP=\angle OMQ$, and since they sum to $180^{\circ}$, each is $90^{\circ}$. So $OM\perp PQ$: the line from the centre to the midpoint of a chord is perpendicular to the chord, i.e.\ the centre lies on the chord's perpendicular bisector. $\blacksquare$

\begin{popfigure}
\begin{tikzpicture}[scale=0.9]
\coordinate (O) at (0,0);
\coordinate (A) at (120:2.5);
\coordinate (B) at (210:2.5);
\coordinate (C) at (-20:2.5);
\draw[popblue, thick] (O) circle (2.5);
\draw[popink, thick] (A)--(B)--(C)--cycle;
\coordinate (M1) at ($(A)!0.5!(B)$);
\coordinate (M2) at ($(B)!0.5!(C)$);
\draw[poporange, dashed, thick] ($(M1)!1.6!(O)$)--($(O)!1.6!(M1)$);
\draw[popgreen, dashed, thick] ($(M2)!1.7!(O)$)--($(O)!1.7!(M2)$);
\fill[popdark] (O) circle (2pt) node[below right, text=popdark] {$O$};
\fill[popink] (A) circle (2pt) node[above left] {$A$};
\fill[popink] (B) circle (2pt) node[below left] {$B$};
\fill[popink] (C) circle (2pt) node[below right] {$C$};
\fill[poporange] (M1) circle (1.5pt) node[left, text=poporange] {$M_1$};
\fill[popgreen] (M2) circle (1.5pt) node[below, text=popgreen] {$M_2$};
\end{tikzpicture}
\legende{The perpendicular bisectors of chords $AB$ and $BC$ meet at the circumcentre $O$, the centre of the circle through $A$, $B$, $C$.}
\end{popfigure}

\begin{methode}[Circumcentre method]
\begin{enumerate}
\item Find the midpoint and gradient of one chord, hence the equation of its perpendicular bisector.
\item Repeat for a second chord.
\item Solve the two bisector equations simultaneously: the intersection is the centre $(h,k)$.
\item The radius is the distance from $(h,k)$ to any of the three points.
\item Write $(x-h)^{2}+(y-k)^{2}=r^{2}$.
\end{enumerate}
\end{methode}

\begin{attention}[Which method should you choose?]
Both methods are accepted at GCE. The \emph{general-equation} method is purely mechanical and safe. The \emph{perpendicular bisector} method is often shorter when two of the points share a coordinate (the bisector is then horizontal or vertical). Whichever you use, \textbf{always check} your final equation against all three points.
\end{attention}

\courssub{Circle Through Two Points with Centre on a Given Line}

Suppose the circle must pass through two given points $A$ and $B$, and its centre must lie on a given line $L$. Geometrically, the centre lies on the perpendicular bisector of $AB$ \emph{and} on $L$: the centre is their intersection, so the circle is unique (unless the lines are parallel, in which case no such circle exists).

Algebraically, write the centre as $(-g,-f)$ and impose:
\begin{itemize}
\item the two points satisfy the general equation (two equations);
\item the centre satisfies the equation of $L$ (one equation).
\end{itemize}
Three equations in $g$, $f$, $c$: the circle is determined.

\begin{exoresolu}[Centre on a given line]
Find the equation of the circle passing through $A(2,3)$ and $B(-1,1)$ whose centre lies on the line $x-3y-11=0$.
\tcblower
Let the circle be $x^{2}+y^{2}+2gx+2fy+c=0$, centre $(-g,-f)$.

Through $A$: $4+9+4g+6f+c=0 \Rightarrow 4g+6f+c=-13 \quad (1)$

Through $B$: $1+1-2g+2f+c=0 \Rightarrow -2g+2f+c=-2 \quad (2)$

Centre on the line: $-g-3(-f)-11=0 \Rightarrow -g+3f=11 \quad (3)$

$(1)-(2)$: $6g+4f=-11 \quad (4)$. From $(3)$, $g=3f-11$. Substituting in $(4)$:
$6(3f-11)+4f=-11 \Rightarrow 22f=55 \Rightarrow f=\tfrac{5}{2}$, $g=-\tfrac{7}{2}$.

From $(2)$: $c=-2+2g-2f=-2-7-5=-14$.

\[ \boxed{x^{2}+y^{2}-7x+5y-14=0} \]
Centre $\left(\tfrac{7}{2},-\tfrac{5}{2}\right)$, radius $\sqrt{\tfrac{49}{4}+\tfrac{25}{4}+14}=\sqrt{\tfrac{130}{4}}=\tfrac{\sqrt{130}}{2}$.

\emph{Check:} $A$: $4+9-14+15-14=0$ \checkmark; $B$: $1+1+7+5-14=0$ \checkmark; centre on line: $\tfrac{7}{2}-3\left(-\tfrac{5}{2}\right)-11=\tfrac{7}{2}+\tfrac{15}{2}-11=0$ \checkmark.
\end{exoresolu}

\begin{popfigure}
\begin{center}
\begin{tikzpicture}
\begin{axis}[width=11cm, axis lines=middle, xlabel={$x$}, ylabel={$y$},
xmin=-4, xmax=9, ymin=-7, ymax=6, axis equal image,
xtick={-2,0,2,4,6,8}, ytick={-6,-4,-2,0,2,4},
legend style={font=\small, at={(0.02,0.98)}, anchor=north west}]
\addplot[popmuted, dashed, domain=-4:9] {(x-11)/3};
\addlegendentry{centre line $x-3y=11$}
\addplot[popblue, thick, domain=0:360, samples=120] ({3.5+5.701*cos(x)},{-2.5+5.701*sin(x)});
\addlegendentry{the circle}
\addplot[only marks, mark=*, popink, mark size=1.8pt] coordinates {(2,3) (-1,1) (3.5,-2.5)};
\node[popink, anchor=south west] at (axis cs:2,3) {$A(2,3)$};
\node[popink, anchor=south east] at (axis cs:-1,1) {$B(-1,1)$};
\node[popink, anchor=north west] at (axis cs:3.5,-2.5) {$C\!\left(\tfrac72,-\tfrac52\right)$};
\end{axis}
\end{tikzpicture}
\end{center}
\legende{The centre is forced to lie on the line $x-3y-11=0$; together with the two points $A$ and $B$ this fixes a unique circle.}
\end{popfigure}

\courssub{Circles Touching a Given Straight Line}

\begin{propriete}[Touching condition]
A circle with centre $(h,k)$ and radius $r$ \cle{touches} the line $ax+by+c=0$ if and only if the perpendicular distance from the centre to the line equals the radius:
\[ \frac{|ah+bk+c|}{\sqrt{a^{2}+b^{2}}}=r. \]
Because of the modulus, this condition usually produces \textbf{two} possible values, one for each side of the line.
\end{propriete}

\begin{methode}['Circle touches line' problems]
\begin{enumerate}
\item Express the centre $(h,k)$ using the other given conditions (e.g.\ centre on a line, or circle through a point).
\item Write the touching condition $\dfrac{|ah+bk+c|}{\sqrt{a^{2}+b^{2}}}=r$.
\item Remove the modulus by considering both signs; solve.
\item For each solution, write the equation and \textbf{check} it satisfies \emph{all} the given data.
\end{enumerate}
\end{methode}

\begin{exoresolu}[Through two points and touching a line]
Find the equations of the circles which pass through $A(1,4)$ and $B(3,2)$ and touch the $x$-axis.
\tcblower
The centre lies on the perpendicular bisector of $AB$. Midpoint of $AB$: $(2,3)$; gradient of $AB$: $\frac{2-4}{3-1}=-1$, so the bisector has gradient $1$:
\[ y-3=1(x-2) \;\Rightarrow\; y=x+1. \]
Write the centre as $(h,h+1)$. Touching the $x$-axis means the radius equals the distance from the centre to $y=0$, i.e.\ $r=|h+1|$.

The circle passes through $A(1,4)$:
\[ (1-h)^{2}+(4-(h+1))^{2}=(h+1)^{2} \]
\[ (1-h)^{2}+(3-h)^{2}=(h+1)^{2} \]
\[ 1-2h+h^{2}+9-6h+h^{2}=h^{2}+2h+1 \]
\[ h^{2}-10h+9=0 \;\Rightarrow\; (h-1)(h-9)=0. \]

\textbf{Case 1:} $h=1$, centre $(1,2)$, $r=2$:
\[ (x-1)^{2}+(y-2)^{2}=4. \]

\textbf{Case 2:} $h=9$, centre $(9,10)$, $r=10$:
\[ (x-9)^{2}+(y-10)^{2}=100. \]

\emph{Check (Case 1):} $A$: $0+4=4$ \checkmark; $B$: $4+0=4$ \checkmark; distance from $(1,2)$ to the $x$-axis is $2=r$ \checkmark. \emph{Check (Case 2):} $A$: $64+36=100$ \checkmark; $B$: $36+64=100$ \checkmark; distance from $(9,10)$ to the $x$-axis is $10=r$ \checkmark. Both circles are valid answers.
\end{exoresolu}

\begin{attention}[Two solutions are the rule, not the exception]
Problems involving \emph{touching a line} almost always yield \textbf{two} circles, because the centre can lie on either side of the line (or the radius can correspond to two positions). Never stop after the first solution: solve the equation completely, and only reject a solution if it contradicts a stated condition (e.g.\ ``the circle lies in the first quadrant'').
\end{attention}

\courssub{Circles Touching the Coordinate Axes}

\begin{propriete}[Circle touching both axes]
A circle of radius $r$ touching \emph{both} coordinate axes has its centre at distance $r$ from each axis, so the centre is $(\pm r,\pm r)$, the signs depending on the quadrant. Its equation is
\[ (x\pm r)^{2}+(y\pm r)^{2}=r^{2}. \]
In the first quadrant: centre $(r,r)$, equation $(x-r)^{2}+(y-r)^{2}=r^{2}$.
\end{propriete}

\begin{popfigure}
\begin{tikzpicture}[scale=1.1]
\draw[->, popmuted] (-0.5,0)--(5,0) node[right] {$x$};
\draw[->, popmuted] (0,-0.5)--(0,5) node[above] {$y$};
\draw[popblue, thick] (2,2) circle (2);
\fill[popdark] (2,2) circle (2pt) node[above right, text=popdark] {$(r,r)$};
\draw[poporange, thick, <->] (2,2)--(2,0) node[midway, right, text=poporange] {$r$};
\draw[popgreen, thick, <->] (2,2)--(0,2) node[midway, above, text=popgreen] {$r$};
\fill[poporange] (2,0) circle (1.5pt);
\fill[popgreen] (0,2) circle (1.5pt);
\node[popmuted, below left] at (0,0) {$O$};
\end{tikzpicture}
\legende{A circle of radius $r$ touching both axes in the first quadrant has centre $(r,r)$: each radius to an axis is perpendicular to that axis.}
\end{popfigure}

\begin{exemplebox}[Touching both axes with an extra condition]
Find the equation of the circle in the first quadrant which touches both axes and passes through the point $(4,2)$.

\emph{Solution.} Centre $(r,r)$, radius $r$. Substituting $(4,2)$:
\[ (4-r)^{2}+(2-r)^{2}=r^{2} \;\Rightarrow\; 16-8r+r^{2}+4-4r+r^{2}=r^{2} \]
\[ r^{2}-12r+20=0 \;\Rightarrow\; (r-2)(r-10)=0. \]
Both $r=2$ and $r=10$ give first-quadrant circles through $(4,2)$:
\[ (x-2)^{2}+(y-2)^{2}=4 \qquad\text{and}\qquad (x-10)^{2}+(y-10)^{2}=100. \]
\end{exemplebox}

\courssub{Strategy: Standard Form or General Form?}

\begin{aretenir}[Choosing the right form]
\begin{itemize}
\item Use the \textbf{standard form} $(x-h)^{2}+(y-k)^{2}=r^{2}$ when the data involve the \emph{centre} or \emph{radius} directly: centre on a line, touching a line (distance condition), touching the axes.
\item Use the \textbf{general form} $x^{2}+y^{2}+2gx+2fy+c=0$ when the data are \emph{points on the circle}: substitution gives \emph{linear} equations in $g$, $f$, $c$, which are easy to solve.
\item Mixed data (points + centre condition): either form works; the general form with the centre condition $-g,-f$ substituted into the line equation is usually cleanest.
\item Remember: \textbf{three independent conditions} are always needed to determine a circle.
\end{itemize}
\end{aretenir}

\begin{exercice}[Consolidation]
\begin{enumerate}
\item Find the equation of the circle through $(0,0)$, $(4,0)$ and $(2,6)$, and state its centre and radius.
\item Find the equation of the circle through $(3,1)$ and $(5,3)$ whose centre lies on the line $y=x+2$.
\item Find the equations of the circles which touch the line $3x+4y=10$, have radius $5$, and whose centres lie on the $x$-axis.
\item A circle touches both coordinate axes and passes through $(1,2)$. Find its possible equations.
\end{enumerate}
\end{exercice}

\begin{corrige}[Consolidation Exercises]
\textbf{1.} Substituting into $x^{2}+y^{2}+2gx+2fy+c=0$: $(0,0)$ gives $c=0$; $(4,0)$ gives $16+8g=0$, so $g=-2$; $(2,6)$ gives $40+4g+12f=0$, so $12f=-40+8=-32$, $f=-\tfrac{8}{3}$. Equation: $x^{2}+y^{2}-4x-\tfrac{16}{3}y=0$; centre $\left(2,\tfrac{8}{3}\right)$, radius $\sqrt{4+\tfrac{64}{9}}=\tfrac{10}{3}$.

\textbf{2.} Centre $(-g,-f)$ on $y=x+2$: $-f=-g+2 \Rightarrow f=g-2$. Points give $6g+2f+c=-10$ and $10g+6f+c=-34$; subtracting: $4g+4f=-24$, i.e.\ $g+f=-6$. With $f=g-2$: $2g-2=-6 \Rightarrow g=-2$, $f=-4$, then $c=-10+12+8=10$. Equation: $x^{2}+y^{2}-4x-8y+10=0$; centre $(2,4)$, radius $\sqrt{10}$.

\textbf{3.} Centre $(h,0)$, radius $5$: $\dfrac{|3h-10|}{5}=5 \Rightarrow |3h-10|=25 \Rightarrow h=\tfrac{35}{3}$ or $h=-5$. Circles: $(x+5)^{2}+y^{2}=25$ and $\left(x-\tfrac{35}{3}\right)^{2}+y^{2}=25$.

\textbf{4.} Centre $(r,r)$ (the point $(1,2)$ is in the first quadrant): $(1-r)^{2}+(2-r)^{2}=r^{2} \Rightarrow r^{2}-6r+5=0 \Rightarrow r=1$ or $r=5$. Circles: $(x-1)^{2}+(y-1)^{2}=1$ and $(x-5)^{2}+(y-5)^{2}=25$.
\end{corrige}

\begin{savaistu}[Circumcircles in real life]
The circle through three points is exactly what surveyors and engineers use to reconstruct a circular arc from three measured points --- for instance when designing a roundabout or a curved section of road. The circumcentre of a triangle is also one of the four classical ``centres'' of a triangle studied since Euclid, alongside the centroid, the orthocentre and the incentre.
\end{savaistu}

\coursec{Section 4: Tangents to a Circle}

In Section 3 we learnt how to determine a circle from geometric data: three points, a diameter, or conditions involving lines and other circles. We now change our point of view. Instead of \emph{building} circles, we study one of the most important lines attached to a circle: the \cle{tangent}. Tangents are everywhere in the GCE Advanced Level examination: they appear in coordinate geometry questions, in locus problems, and in the study of families of circles. Master this section and a large family of exam questions becomes routine.

\courssub{What is a Tangent?}

\begin{experience}[From a secant to a tangent]
Draw a circle and a straight line. Depending on the position of the line, three things can happen:
\begin{itemize}
\item the line misses the circle completely (no intersection);
\item the line cuts the circle in \emph{two} distinct points --- such a line is called a \cle{secant};
\item the line meets the circle in \emph{exactly one} point.
\end{itemize}
Now imagine a secant through two points $A$ and $B$ of the circle, and let $B$ slide along the circle towards $A$. The two intersection points merge into one, and the secant settles into a limiting position: the tangent at $A$.
\end{experience}

\begin{definition}[Tangent to a circle]
A \cle{tangent} to a circle is a straight line which meets the circle in exactly one point. That point is called the \cle{point of contact}. A line which cuts the circle in two distinct points is a \cle{secant}; a line which does not meet the circle is said to be \emph{exterior} to it.
\end{definition}

\begin{propriete}[Radius--tangent perpendicularity]
The tangent at any point of a circle is perpendicular to the radius drawn to the point of contact. (This result is examinable and its proof is instructive.)
\end{propriete}

\textbf{Proof.} Let the circle have centre $C$ and radius $r$, and let $T$ be the point of contact of a tangent line $\ell$. Take any other point $Q$ on $\ell$, $Q \neq T$. Since $\ell$ meets the circle only at $T$, the point $Q$ lies \emph{outside} the circle, so $CQ > r = CT$. Thus $CT$ is the \emph{shortest} distance from $C$ to the line $\ell$. But the shortest distance from a point to a line is the perpendicular distance. Hence $CT \perp \ell$. $\blacksquare$

\begin{popfigure}
\begin{center}
\begin{tikzpicture}[scale=1.1]
\coordinate (C) at (0,0);
\coordinate (P) at (1.5,2);
\draw[popblue, thick] (C) circle (2.5);
\draw[popdark, thick] (C) -- (P) node[midway, above left, popdark]{$r$};
\draw[poporange, thick] ($(P)!1.6cm!90:(C)$) -- ($(P)!1.6cm!-90:(C)$);
\fill[popink] (C) circle (1.5pt) node[below left]{$C$};
\fill[popink] (P) circle (1.5pt) node[above right]{$P(x_1,y_1)$};
\draw[popgreen] ($(P)!0.28cm!(C)$) -- ($(P)!0.28cm!(C)!0.28cm!90:(C)$) -- ($(P)!0.28cm!-90:(C)!0.28cm!90:(P)$);
\node[poporange] at (3.1,1.2) {tangent};
\end{tikzpicture}
\end{center}
\legende{The tangent at $P$ is perpendicular to the radius $CP$.}
\end{popfigure}

\courssub{Tangent to the Circle $x^2+y^2=r^2$ at a Given Point}

We now turn the geometric property into an algebraic formula. Consider the circle $x^2+y^2=r^2$ with centre the origin $O$, and a point $P(x_1,y_1)$ \emph{on} the circle, so that
\[ x_1^2+y_1^2=r^2. \]
The radius $OP$ has gradient $\dfrac{y_1}{x_1}$ (provided $x_1\neq 0$; the cases $x_1=0$ or $y_1=0$ give vertical or horizontal tangents which are easily handled separately). Since the tangent is perpendicular to $OP$, its gradient is $-\dfrac{x_1}{y_1}$. The tangent through $P(x_1,y_1)$ therefore has equation
\[ y-y_1=-\frac{x_1}{y_1}(x-x_1). \]
Multiplying by $y_1$ and rearranging:
\[ y y_1-y_1^2=-x x_1+x_1^2 \quad\Longrightarrow\quad x x_1+y y_1=x_1^2+y_1^2=r^2. \]

\begin{propriete}[Tangent to $x^2+y^2=r^2$ at $(x_1,y_1)$]
The equation of the tangent to the circle $x^2+y^2=r^2$ at the point $P(x_1,y_1)$ on the circle is
\[ \boxed{\,x x_1+y y_1=r^2\,} \]
(The derivation above, using the gradient of the radius and the perpendicular gradient of the tangent, is examinable.)
\end{propriete}

Notice the beautiful pattern: to obtain the tangent at $(x_1,y_1)$, replace one $x$ by $x_1$ and one $y$ by $y_1$ in $x^2+y^2=r^2$. This is the simplest case of the \emph{replace rule} that we shall generalise below.

\begin{attention}[Check that the point is on the circle!]
The formula $xx_1+yy_1=r^2$ is valid \emph{only} when $(x_1,y_1)$ actually lies on the circle. Always verify that $x_1^2+y_1^2=r^2$ before applying it. If the point is not on the circle, the line $xx_1+yy_1=r^2$ is still a perfectly good line (it is in fact the \emph{polar} of the point with respect to the circle), but it is \textbf{not} a tangent through that point.
\end{attention}

\begin{exoresolu}[Tangent at a point on $x^2+y^2=25$]
Find the equation of the tangent to the circle $x^2+y^2=25$ at the point $P(3,4)$.
\tcblower
\textbf{Solution.} First check: $3^2+4^2=9+16=25$, so $P$ lies on the circle. Applying $xx_1+yy_1=r^2$ with $(x_1,y_1)=(3,4)$ and $r^2=25$:
\[ 3x+4y=25. \]
\emph{Verification by gradients:} the radius $OP$ has gradient $\tfrac{4}{3}$; the line $3x+4y=25$ has gradient $-\tfrac{3}{4}$, and $\tfrac{4}{3}\times(-\tfrac{3}{4})=-1$. Perpendicular, as required.
\end{exoresolu}

\courssub{Tangent to a General Circle: the Replace Rule}

For the circle $(x-a)^2+(y-b)^2=r^2$ with centre $C(a,b)$, the same gradient argument applies. The radius from $C$ to $P(x_1,y_1)$ has gradient $\dfrac{y_1-b}{x_1-a}$, so the tangent has gradient $-\dfrac{x_1-a}{y_1-b}$, giving
\[ (x-x_1)(x_1-a)+(y-y_1)(y_1-b)=0, \]
which simplifies, using $(x_1-a)^2+(y_1-b)^2=r^2$, to
\[ (x-a)(x_1-a)+(y-b)(y_1-b)=r^2. \]

For the \emph{general} circle
\[ x^2+y^2+2gx+2fy+c=0 \quad\text{(centre $(-g,-f)$, radius $r=\sqrt{g^2+f^2-c}$)}, \]
expanding the previous formula with $a=-g$, $b=-f$ and $r^2=g^2+f^2-c$ leads to the following memorable rule.

\begin{propriete}[The replace rule for tangents]
The tangent at $P(x_1,y_1)$ to the circle $x^2+y^2+2gx+2fy+c=0$ is
\[ \boxed{\;x x_1+y y_1+g(x+x_1)+f(y+y_1)+c=0\;} \]
obtained from the equation of the circle by the replacements
\[ x^2\to x x_1,\qquad y^2\to y y_1,\qquad 2gx\to g(x+x_1),\qquad 2fy\to f(y+y_1),\qquad c\to c. \]
\end{propriete}

\begin{exemplebox}[Using the replace rule]
Find the tangent to the circle $x^2+y^2-4x+6y-12=0$ at the point $(5,1)$.\\
\textbf{Check that the point lies on the circle:}
\[ 5^2+1^2-4(5)+6(1)-12 = 25+1-20+6-12 = 0. \]
Good. Here $2g=-4\Rightarrow g=-2$, $2f=6\Rightarrow f=3$, $c=-12$. Apply the replace rule:
\[ 5x+1\cdot y + (-2)(x+5) + 3(y+1) -12 = 0 \]
\[ 5x+y -2x-10 +3y+3 -12 = 0 \]
\[ 3x+4y -19 = 0. \]
\end{exemplebox}

\begin{savaistu}[The polar line]
If $P(x_1,y_1)$ is \emph{outside} the circle, the line $xx_1+yy_1+g(x+x_1)+f(y+y_1)+c=0$ is called the \emph{polar} of $P$ with respect to the circle. It passes through the two points of contact of the tangents drawn from $P$ to the circle. This is a powerful tool in projective geometry and appears in advanced locus problems.
\end{savaistu}

\courssub{Condition for a Line to be Tangent to a Circle}

Suppose now that we are given a \emph{line} and a \emph{circle} and we ask: is the line a tangent? Geometrically the answer is immediate.

\begin{methode}[Tangency test: distance from centre equals radius]
A line $\ell$ is tangent to a circle if and only if the \cle{perpendicular distance} from the centre of the circle to $\ell$ equals the radius.
\begin{enumerate}
\item Write the line in the form $Ax+By+C=0$ and identify the centre $(a,b)$ and radius $r$ of the circle.
\item Compute $d=\dfrac{|Aa+Bb+C|}{\sqrt{A^2+B^2}}$.
\item If $d=r$: tangent (one intersection). If $d<r$: secant (two intersections). If $d>r$: no intersection.
\end{enumerate}
\end{methode}

Apply this to the line $y=mx+c$, i.e. $mx-y+c=0$, and the circle $x^2+y^2=r^2$ centred at the origin:
\[ \frac{|c|}{\sqrt{m^2+1}}=r \quad\Longrightarrow\quad \boxed{\,c^2=r^2(1+m^2)\,}. \]

\begin{propriete}[Tangency condition for $y=mx+c$ and $x^2+y^2=r^2$]
The line $y=mx+c$ is tangent to the circle $x^2+y^2=r^2$ if and only if $c^2=r^2(1+m^2)$. Consequently the lines
\[ y=mx\pm r\sqrt{1+m^2} \]
are the two tangents of gradient $m$ to the circle.
\end{propriete}

\begin{propriete}[Tangency condition for a general circle]
The line $Ax+By+C=0$ is tangent to the circle $x^2+y^2+2gx+2fy+c=0$ (centre $(-g,-f)$, radius $r=\sqrt{g^2+f^2-c}$) if and only if
\[ \frac{|A(-g)+B(-f)+C|}{\sqrt{A^2+B^2}} = \sqrt{g^2+f^2-c}. \]
Equivalently, substituting $y=mx+c$ into the circle's equation yields a quadratic in $x$ whose discriminant is zero.
\end{propriete}

\begin{popfigure}
\begin{center}
\begin{tikzpicture}
\begin{axis}[width=9.5cm, height=7cm, axis lines=middle, xlabel={$x$}, ylabel={$y$},
xmin=-5.5, xmax=5.5, ymin=-5.5, ymax=5.5, samples=100, axis equal image,
xtick={-4,-2,2,4}, ytick={-4,-2,2,4}, tick label style={font=\small}]
\addplot[popblue, thick, domain=0:360] ({3*cos(x)},{3*sin(x)});
\addplot[poporange, thick, domain=-5:5] {x+3*sqrt(max(0,2))};
\addplot[poporange, thick, domain=-5:5] {x-3*sqrt(max(0,2))};
\addplot[popmuted, dashed, domain=-5:3] {x+6};
\addplot[popmuted, dashed, domain=-3:5] {x};
\end{axis}
\end{tikzpicture}
\end{center}
\legende{Lines $y=x+c$ and the circle $x^2+y^2=9$. Only $c=\pm3\sqrt{2}$ (orange) give tangents, since $c^2=9(1+1^2)=18$.}
\end{popfigure}

\begin{exoresolu}[Showing a line is a tangent]
Show that the line $3x+4y=30$ is a tangent to the circle $x^2+y^2+2x-4y-20=0$, and find the point of contact.
\tcblower
\textbf{Solution.} Complete the square for the circle:
\[ (x+1)^2+(y-2)^2 = 1+4+20 = 25, \]
so the centre is $C(-1,2)$ and the radius is $r=5$.
Distance from $C$ to the line $3x+4y-30=0$:
\[ d=\frac{|3(-1)+4(2)-30|}{\sqrt{3^2+4^2}}=\frac{|-3+8-30|}{5}=\frac{25}{5}=5 = r. \]
Hence the line is a tangent.

\textbf{Point of contact:} it is the foot of the perpendicular from $C(-1,2)$ to the line. The perpendicular through $C$ has direction vector $(3,4)$ (the normal of the line). Its parametric form: $(-1+3t,\,2+4t)$. Substitute into $3x+4y=30$:
\[ 3(-1+3t)+4(2+4t)=30 \;\Longrightarrow\; -3+9t+8+16t=30 \;\Longrightarrow\; 25t+5=30 \;\Longrightarrow\; t=1. \]
Point of contact: $(2,6)$. Check: $2^2+6^2+2(2)-4(6)-20=4+36+4-24-20=0$. \checkmark
\end{exoresolu}

\courssub{Tangents from an External Point}

Place a point $P$ outside a circle. How many tangents can you draw from $P$ to the circle? A sketch convinces us immediately: exactly \cle{two}, and they are symmetric about the line $CP$.

\begin{attention}[Two tangents, not one!]
From an external point there are \emph{exactly two} tangents to a circle. When an exam question asks for ``the tangents from $P$'', do not stop after finding one --- the second one carries marks too. Algebraically this appears as a quadratic in the gradient $m$; both roots must be reported.
\end{attention}

\begin{methode}[Finding the tangents from an external point (gradient method)]
To find the tangents from an external point $P(x_1,y_1)$ to a circle with centre $C$ and radius $r$:
\begin{enumerate}
\item Write a variable line through $P$: $y-y_1=m(x-x_1)$, rearranged as $mx-y+(y_1-mx_1)=0$.
\item Impose the tangency condition: perpendicular distance from $C$ to the line equals $r$.
\item Solve the resulting quadratic in $m$; each root gives one tangent.
\end{enumerate}
\end{methode}

\begin{exoresolu}[Tangents from an external point]
Find the equations of the tangents from $P(5,5)$ to the circle $x^2+y^2=10$.
\tcblower
\textbf{Solution.} Note $5^2+5^2=50>10$, so $P$ is external. A line through $P$: $y-5=m(x-5)$, i.e. $mx-y+5-5m=0$. Distance from $O(0,0)$ must equal $\sqrt{10}$:
\[ \frac{|5-5m|}{\sqrt{m^2+1}}=\sqrt{10}
\;\Longrightarrow\; 25(1-m)^2=10(m^2+1)
\;\Longrightarrow\; 15m^2-50m+15=0
\;\Longrightarrow\; 3m^2-10m+3=0. \]
Factorising: $(3m-1)(m-3)=0$, so $m=3$ or $m=\tfrac13$. The two tangents are
\[ y-5=3(x-5)\;\Rightarrow\; 3x-y-10=0, \qquad
y-5=\tfrac13(x-5)\;\Rightarrow\; x-3y+10=0. \]
\end{exoresolu}

\courssub{Length of the Tangent from an External Point}

Let $P$ be an external point, $T$ the point of contact of a tangent from $P$, and $C$ the centre. By the radius--tangent property, $\angle CTP=90^\circ$, so triangle $CPT$ is right-angled at $T$. Pythagoras gives
\[ CP^2=CT^2+PT^2=r^2+PT^2. \]

\begin{propriete}[Length of tangent from an external point]
The length of the tangent from an external point $P$ to a circle with centre $C$ and radius $r$ is
\[ \boxed{\,PT=\sqrt{CP^2-r^2}\,} \]
For the general circle $x^2+y^2+2gx+2fy+c=0$ and $P(x_1,y_1)$, since $CP^2-r^2=(x_1+g)^2+(y_1+f)^2-(g^2+f^2-c)$, this becomes
\[ PT=\sqrt{x_1^2+y_1^2+2gx_1+2fy_1+c}. \]
So the tangent length is simply the square root of the result of substituting $P$ into the left-hand side of the circle's equation (the \emph{power} of the point).
\end{propriete}

\begin{popfigure}
\begin{center}
\begin{tikzpicture}[scale=1.05]
\coordinate (C) at (0,0);
\coordinate (P) at (5,0);
\coordinate (T) at (1.2,1.6);
\draw[popblue, thick] (C) circle (2);
\draw[popdark] (C) -- (P) node[midway, below]{$d$};
\draw[popgreen, thick] (C) -- (T) node[midway, above left]{$r$};
\draw[poporange, thick] (P) -- (T) node[midway, above right]{$\sqrt{d^2-r^2}$};
\draw[poporange, thick] (P) -- ($(T)!2!($(C)!(T)!(P)$)$);
\fill[popink] (C) circle (1.5pt) node[below left]{$C$};
\fill[popink] (P) circle (1.5pt) node[right]{$P$};
\fill[popink] (T) circle (1.5pt) node[above left]{$T$};
\coordinate (A) at ($(T)!0.25cm!(C)$);
\coordinate (B) at ($(T)!0.25cm!(P)$);
\draw[popgreen] (A) -- ($(A)+(B)-(T)$) -- (B);
\end{tikzpicture}
\end{center}
\legende{The right triangle $CPT$: $PT=\sqrt{CP^2-r^2}$, and both tangents from $P$ have the same length.}
\end{popfigure}

\begin{exemplebox}[Length of a tangent]
Find the length of the tangent from $P(7,1)$ to the circle $x^2+y^2-4x+2y-4=0$.\\
\textbf{Solution.} Substitute $P$ into the left-hand side:
\[ 7^2+1^2-4(7)+2(1)-4=49+1-28+2-4=20. \]
Hence $PT=\sqrt{20}=2\sqrt{5}$. \emph{Check via centre and radius:} $C=(2,-1)$, $r^2=4+1+4=9$; $CP^2=(7-2)^2+(1+1)^2=25+4=29$; $PT=\sqrt{29-9}=\sqrt{20}$. \checkmark
\end{exemplebox}

\begin{aretenir}[Key results of Section 4]
\begin{itemize}
\item The tangent at $P$ is perpendicular to the radius $CP$.
\item Tangent to $x^2+y^2=r^2$ at $(x_1,y_1)$: $xx_1+yy_1=r^2$.
\item Replace rule for $x^2+y^2+2gx+2fy+c=0$ at $(x_1,y_1)$:
      $x^2\to xx_1$, $y^2\to yy_1$, $2gx\to g(x+x_1)$, $2fy\to f(y+y_1)$, $c\to c$.
\item Tangency condition: distance from centre to line $=$ radius.
      For $y=mx+c$ and $x^2+y^2=r^2$: $c^2=r^2(1+m^2)$.
\item From an external point: exactly two tangents, of equal length $\sqrt{CP^2-r^2} = \sqrt{x_1^2+y_1^2+2gx_1+2fy_1+c}$.
\end{itemize}
\end{aretenir}

\begin{exercice}[Practice]
\begin{enumerate}
\item Find the tangent to $x^2+y^2=29$ at $(5,-2)$.
\item Show that $y=2x+5\sqrt{5}$ touches $x^2+y^2=25$.
\item Find the tangents from $(0,5)$ to $x^2+y^2=5$ and their lengths.
\item Find the tangent to $x^2+y^2-6x-8y=0$ at the origin.
\item A circle has centre $(2,-3)$ and radius $4$. Find the equations of the tangents of gradient $1$ to this circle.
\end{enumerate}
\end{exercice}

\begin{corrige}[Exercise 1]
\begin{enumerate}
\item Check: $5^2+(-2)^2=25+4=29$. Tangent: $5x-2y=29$.
\item $c^2=(5\sqrt{5})^2=125$; $r^2(1+m^2)=25(1+4)=125$. Condition satisfied $\Rightarrow$ tangent.
\item Line through $(0,5)$: $y=mx+5$. Distance from $(0,0)$ to $mx-y+5=0$ is $\frac{5}{\sqrt{m^2+1}}=\sqrt{5} \Rightarrow m^2+1=5 \Rightarrow m=\pm2$. Tangents: $y=2x+5$ and $y=-2x+5$. Length $=\sqrt{0^2+5^2-5}=\sqrt{20}=2\sqrt{5}$.
\item Circle: $x^2+y^2-6x-8y=0$. At $(0,0)$: replace rule gives $-3x-4y=0$, i.e. $3x+4y=0$.
\item Circle: $(x-2)^2+(y+3)^2=16$. Tangents of gradient $1$: $y=x+c$. Distance from $(2,-3)$ to $x-y+c=0$ is $\frac{|2+3+c|}{\sqrt{2}}=4 \Rightarrow |c+5|=4\sqrt{2} \Rightarrow c=-5\pm4\sqrt{2}$. Tangents: $y=x-5+4\sqrt{2}$ and $y=x-5-4\sqrt{2}$.
\end{enumerate}
\end{corrige}

\coursec{Section 5: Parametric Equations of a Circle}

In Section 4 we studied tangents to a circle, and you will have noticed that the algebra quickly became heavy: substituting a line into the circle's equation, solving quadratics, imposing discriminant conditions. Much of that heaviness comes from the fact that the Cartesian equation $x^{2}+y^{2}=r^{2}$ describes the circle \emph{implicitly}: $x$ and $y$ are tangled together in one relation. In this section we learn a second, very powerful way of describing a circle: we give $x$ and $y$ \emph{separately}, each as a function of a single variable $\theta$ called a \cle{parameter}. This \cle{parametric form} will allow us to write down a ``general point'' on the circle and will turn many difficult problems (maximum and minimum distances, tangents, loci) into simple trigonometry.

\courssub{The Idea of a Parameter: From Observation to Definition}

\begin{experience}[A point moving around a circle]
Imagine a small bead sliding at constant speed around a circular wire of radius $r$ centred at the origin $O$, like the hand of a clock at the Biyem-Assi stadium scoreboard, but turning anticlockwise. At any instant, the position of the bead is completely known once we know the \emph{angle} $\theta$ that the radius $OP$ makes with the positive $x$-axis. Drop a perpendicular from $P$ to the $x$-axis, meeting it at $M$. In the right-angled triangle $OMP$:
\[
\cos\theta=\frac{OM}{OP}=\frac{x}{r}, \qquad \sin\theta=\frac{MP}{OP}=\frac{y}{r}.
\]
Hence $x=r\cos\theta$ and $y=r\sin\theta$. So the single number $\theta$ determines the whole position $(x,y)$ of the bead.
\end{experience}

This observation is the heart of the section. Instead of one equation linking $x$ and $y$, we describe the circle by \emph{two} equations, each giving one coordinate in terms of $\theta$.

\begin{definition}[Parametric equations of a circle]
The \cle{parametric equations} of the circle with centre $(a,b)$ and radius $r$ are
\[
x=a+r\cos\theta, \qquad y=b+r\sin\theta,
\]
where $\theta$ is the \cle{parameter}. In particular, for the circle centred at the origin,
\[
x=r\cos\theta, \qquad y=r\sin\theta.
\]
As $\theta$ varies from $0$ to $2\pi$, the point $(x,y)$ describes the whole circle exactly once (anticlockwise).
\end{definition}

\begin{attention}[Where is $\theta$ measured?]
The parameter $\theta$ is the angle measured \textbf{at the centre of the circle}, from the direction parallel to the positive $x$-axis, to the radius joining the centre to the point. It is \emph{not} measured at the origin, unless the centre happens to be the origin. For the circle $(x-2)^{2}+(y+1)^{2}=9$, the angle $\theta$ is measured at the point $(2,-1)$.
\end{attention}

\begin{popfigure}
\begin{center}
\begin{tikzpicture}[scale=1.5]
\coordinate (C) at (1.6,1.1);
\draw[popline,->] (-0.6,0) -- (4.4,0) node[below left] {$x$};
\draw[popline,->] (0,-0.6) -- (0,3.4) node[below left] {$y$};
\draw[popblue,thick] (C) circle (1.5);
\fill[popdark] (C) circle (0.045) node[below left] {$C(a,b)$};
\coordinate (P) at ($(C)+(40:1.5)$);
\fill[poporange] (P) circle (0.05) node[above right] {$P(a+r\cos\theta,\;b+r\sin\theta)$};
\draw[popdark,thick] (C) -- (P) node[midway,above left] {$r$};
\draw[popmuted,dashed] (C) -- ($(C)+(1.9,0)$);
\draw[popgreen,thick] ($(C)+(0.55,0)$) arc (0:40:0.55);
\node[popgreen] at ($(C)+(20:0.85)$) {$\theta$};
\draw[popmuted,dashed] (P) -- ($(P)-(0,1.1)$);
\draw[popmuted,dashed] (C) -- (0,1.1);
\end{tikzpicture}
\end{center}
\legende{The parameter $\theta$ is the angle at the centre $C(a,b)$ between the positive $x$-direction and the radius $CP$.}
\end{popfigure}

\courssub{From Parametric to Cartesian Form}

Suppose we are given parametric equations and we want the ordinary (Cartesian) equation of the circle. The tool is the fundamental identity $\cos^{2}\theta+\sin^{2}\theta=1$.

\begin{methode}[Eliminating the parameter]
To convert parametric equations to Cartesian form:
\begin{enumerate}
\item Rearrange the two equations to isolate $\cos\theta$ and $\sin\theta$:
\[
\cos\theta=\frac{x-a}{r}, \qquad \sin\theta=\frac{y-b}{r}.
\]
\item Substitute into the identity $\cos^{2}\theta+\sin^{2}\theta=1$:
\[
\left(\frac{x-a}{r}\right)^{2}+\left(\frac{y-b}{r}\right)^{2}=1.
\]
\item Multiply through by $r^{2}$ to obtain $(x-a)^{2}+(y-b)^{2}=r^{2}$.
\end{enumerate}
\end{methode}

This calculation also \emph{proves} that the parametric equations really do describe the circle: every point given by the parametric form satisfies the Cartesian equation, and conversely every point of the circle arises for some $\theta$.

\begin{exemplebox}[Parametric to Cartesian]
Convert to Cartesian form: $x=3+2\cos\theta$, $y=-1+2\sin\theta$.

\textbf{Solution.} Isolate the trigonometric functions: $\cos\theta=\dfrac{x-3}{2}$ and $\sin\theta=\dfrac{y+1}{2}$. Using $\cos^{2}\theta+\sin^{2}\theta=1$:
\[
\frac{(x-3)^{2}}{4}+\frac{(y+1)^{2}}{4}=1
\quad\Longrightarrow\quad
(x-3)^{2}+(y+1)^{2}=4.
\]
This is the circle with centre $(3,-1)$ and radius $2$.
\end{exemplebox}

\courssub{From Cartesian to Parametric Form}

The reverse journey is just as important. Given $(x-a)^{2}+(y-b)^{2}=r^{2}$, divide by $r^{2}$:
\[
\left(\frac{x-a}{r}\right)^{2}+\left(\frac{y-b}{r}\right)^{2}=1.
\]
Since the two squares sum to $1$, we may call the first $\cos^{2}\theta$ and the second $\sin^{2}\theta$, giving $x=a+r\cos\theta$, $y=b+r\sin\theta$.

\begin{exemplebox}[Cartesian to parametric]
Write parametric equations for the circle $x^{2}+y^{2}-4x+6y-12=0$.

\textbf{Solution.} Complete the square:
\[
(x-2)^{2}-4+(y+3)^{2}-9-12=0
\quad\Longrightarrow\quad
(x-2)^{2}+(y+3)^{2}=25.
\]
The centre is $(2,-3)$ and the radius is $5$. Hence parametric equations are
\[
x=2+5\cos\theta, \qquad y=-3+5\sin\theta.
\]
\end{exemplebox}

\begin{popfigure}
\begin{center}
\begin{tikzpicture}
\begin{axis}[
  width=9.5cm, height=9.5cm,
  axis lines=middle,
  xlabel={$x$}, ylabel={$y$},
  xmin=-4.5, xmax=4.5, ymin=-4.5, ymax=4.5,
  axis equal image,
  xtick={-4,-2,2,4}, ytick={-4,-2,2,4},
  tick label style={font=\small},
]
\addplot[popblue, thick, domain=0:360, samples=200, variable=\t] ({3*cos(t)},{3*sin(t)});
\addplot[only marks, mark=*, mark size=1.8pt, poporange] coordinates {(3,0) (2.121,2.121) (0,3) (-2.121,2.121) (-3,0) (0,-3)};
\node[popdark, font=\small] at (axis cs:3.35,-0.55) {$\theta=0$};
\node[popdark, font=\small] at (axis cs:2.6,2.75) {$\theta=\frac{\pi}{4}$};
\node[popdark, font=\small] at (axis cs:0.45,3.35) {$\theta=\frac{\pi}{2}$};
\node[popdark, font=\small] at (axis cs:-2.7,2.7) {$\theta=\frac{3\pi}{4}$};
\node[popdark, font=\small] at (axis cs:-3.4,-0.55) {$\theta=\pi$};
\node[popdark, font=\small] at (axis cs:0.55,-3.4) {$\theta=\frac{3\pi}{2}$};
\end{axis}
\end{tikzpicture}
\end{center}
\legende{The circle $x=3\cos\theta$, $y=3\sin\theta$ drawn parametrically, with points labelled by their values of $\theta$.}
\end{popfigure}

\courssub{The General Point on a Circle: A Powerful Technique}

\begin{methode}[Using a general point on the circle]
Any point on the circle with centre $(a,b)$ and radius $r$ may be written as
\[
P=(a+r\cos\theta,\; b+r\sin\theta).
\]
When a problem involves ``a point moving on the circle'', or asks for a maximum or minimum, or a locus, write the point in this form and work with the single variable $\theta$. Distances, gradients and midpoints then become trigonometric expressions, which are often much easier to handle than Cartesian algebra.
\end{methode}

\begin{exoresolu}[Maximum and minimum distance from a point to a circle]
A point $P$ moves on the circle $x^{2}+y^{2}=4$. Find the greatest and least values of the distance from $P$ to the fixed point $A(5,0)$.
\tcblower
\textbf{Solution.} Write $P=(2\cos\theta,\,2\sin\theta)$. Then
\begin{align*}
AP^{2}&=(2\cos\theta-5)^{2}+(2\sin\theta)^{2}\\
&=4\cos^{2}\theta-20\cos\theta+25+4\sin^{2}\theta\\
&=4(\cos^{2}\theta+\sin^{2}\theta)-20\cos\theta+25\\
&=29-20\cos\theta.
\end{align*}
Since $-1\leq\cos\theta\leq1$, the expression $29-20\cos\theta$ is greatest when $\cos\theta=-1$ and least when $\cos\theta=1$:
\[
AP^{2}_{\max}=29+20=49 \;\Rightarrow\; AP_{\max}=7, \qquad
AP^{2}_{\min}=29-20=9 \;\Rightarrow\; AP_{\min}=3.
\]
Geometrically, these are the distances from $A$ to the far and near ends of the diameter through $A$. Notice how the parametric form reduced the problem to the boundedness of $\cos\theta$ — no calculus, no discriminants.
\end{exoresolu}

\begin{exoresolu}[A locus found parametrically]
$P=(3\cos\theta,\,3\sin\theta)$ moves on the circle $x^{2}+y^{2}=9$, and $M$ is the midpoint of $P$ and the fixed point $B(7,0)$. Find the locus of $M$.
\tcblower
\textbf{Solution.} The coordinates of $M$ are
\[
M=\left(\frac{3\cos\theta+7}{2},\; \frac{3\sin\theta+0}{2}\right).
\]
Write $M=(X,Y)$, so $X=\dfrac{7+3\cos\theta}{2}$ and $Y=\dfrac{3\sin\theta}{2}$. Isolate the trigonometric functions:
\[
\cos\theta=\frac{2X-7}{3}, \qquad \sin\theta=\frac{2Y}{3}.
\]
Applying $\cos^{2}\theta+\sin^{2}\theta=1$:
\[
\frac{(2X-7)^{2}}{9}+\frac{4Y^{2}}{9}=1
\quad\Longrightarrow\quad
(2X-7)^{2}+4Y^{2}=9
\quad\Longrightarrow\quad
\left(X-\tfrac{7}{2}\right)^{2}+Y^{2}=\tfrac{9}{4}.
\]
The locus of $M$ is the circle with centre $\left(\dfrac{7}{2},0\right)$ and radius $\dfrac{3}{2}$ — exactly what we would expect, since a midpoint construction shrinks the original circle by a factor of $\tfrac{1}{2}$ about $B$.
\end{exoresolu}

\begin{exoresolu}[Tangent at a parametric point]
Find the equation of the tangent to the circle $x^{2}+y^{2}=25$ at the point $P=(5\cos\theta,\,5\sin\theta)$, and deduce the tangent at the point where $\theta=\dfrac{\pi}{3}$.
\tcblower
\textbf{Solution.} The radius $OP$ has gradient $\dfrac{5\sin\theta}{5\cos\theta}=\tan\theta$, so the tangent, being perpendicular to $OP$, has gradient $-\cot\theta$ (for $\sin\theta\neq0$). Its equation is
\[
y-5\sin\theta=-\frac{\cos\theta}{\sin\theta}\,(x-5\cos\theta).
\]
Multiplying by $\sin\theta$ and rearranging:
\[
x\cos\theta+y\sin\theta=5(\cos^{2}\theta+\sin^{2}\theta)=5.
\]
This elegant result, $x\cos\theta+y\sin\theta=5$, is the tangent at the point with parameter $\theta$. For $\theta=\dfrac{\pi}{3}$: $\cos\theta=\tfrac12$, $\sin\theta=\tfrac{\sqrt3}{2}$, giving
\[
\frac{x}{2}+\frac{\sqrt3\,y}{2}=5
\quad\Longrightarrow\quad
x+\sqrt3\,y=10.
\]
\end{exoresolu}

\begin{exoresolu}[Maximum distance from a point to a circle — general case]
A point $P$ moves on the circle $x^{2}+y^{2}=1$. Find the greatest and least values of the distance $PQ$, where $Q=(3,4)$.
\tcblower
\textbf{Solution.} Write $P=(\cos\theta,\,\sin\theta)$. Then
\begin{align*}
PQ^{2}&=(\cos\theta-3)^{2}+(\sin\theta-4)^{2}\\
&=\cos^{2}\theta-6\cos\theta+9+\sin^{2}\theta-8\sin\theta+16\\
&=(\cos^{2}\theta+\sin^{2}\theta)+25-6\cos\theta-8\sin\theta\\
&=26-6\cos\theta-8\sin\theta.
\end{align*}
We express $6\cos\theta+8\sin\theta$ in the form $R\cos(\theta-\alpha)$:
$R=\sqrt{6^{2}+8^{2}}=10$, $\cos\alpha=\frac{6}{10}=\frac{3}{5}$, $\sin\alpha=\frac{8}{10}=\frac{4}{5}$.
Thus $6\cos\theta+8\sin\theta=10\cos(\theta-\alpha)$, and since $-1\leq\cos(\theta-\alpha)\leq1$,
\[
-10 \leq 6\cos\theta+8\sin\theta \leq 10.
\]
Hence
\[
26-10 \leq PQ^{2} \leq 26+10 \quad\Longrightarrow\quad 16 \leq PQ^{2} \leq 36.
\]
Therefore $PQ_{\min}=4$ and $PQ_{\max}=6$.
\end{exoresolu}

\begin{attention}[Common mistakes with parametric forms]
\begin{itemize}
\item \textbf{Wrong centre.} For $x=2+5\cos\theta$, $y=-3+5\sin\theta$, the centre is $(2,-3)$, not $(-2,3)$ and not $(2,3)$. Read the signs directly from the parametric equations.
\item \textbf{Forgetting the radius in front of $\cos\theta$ and $\sin\theta$.} The pair $x=a+\cos\theta$, $y=b+\sin\theta$ describes a circle of radius $1$, whatever the values of $a$ and $b$.
\item \textbf{Mixing up the two coefficients.} If $x=1+3\cos\theta$ but $y=2+4\sin\theta$, the locus is \emph{not} a circle (it is an ellipse); a circle requires the \emph{same} coefficient $r$ in both equations.
\item \textbf{Measuring $\theta$ at the origin.} Remember: $\theta$ is always measured at the \emph{centre} of the circle.
\item \textbf{Sign slips when eliminating $\theta$.} From $y=b+r\sin\theta$ we get $\sin\theta=\dfrac{y-b}{r}$: subtract $b$ \emph{before} dividing by $r$.
\end{itemize}
\end{attention}

\begin{aretenir}[GCE exam technique: when parametric form wins]
In the GCE Advanced Level examination, watch for questions asking for the \textbf{greatest or least distance} from a fixed point or line to a circle, or for a \textbf{locus} involving a moving point on a circle. These are usually fastest by parametric form:
\begin{enumerate}
\item Write the moving point as $(a+r\cos\theta,\;b+r\sin\theta)$.
\item Express the required quantity (distance, midpoint, gradient) in terms of $\theta$.
\item Use $-1\leq\cos\theta\leq1$, $-1\leq\sin\theta\leq1$, or $\cos^{2}\theta+\sin^{2}\theta=1$ to finish.
\end{enumerate}
This route routinely saves several lines of algebra and avoids the discriminant computations of Section 4. Also memorise the tangent at a parametric point of $x^{2}+y^{2}=r^{2}$:
\[
x\cos\theta+y\sin\theta=r.
\]
For a circle $(x-a)^{2}+(y-b)^{2}=r^{2}$, the tangent at the point with parameter $\theta$ is
\[
(x-a)\cos\theta+(y-b)\sin\theta=r.
\]
\end{aretenir}

\begin{savaistu}[Why engineers love parameters]
Parametric descriptions are not just an exam trick. When a CNC machine in a Douala workshop cuts a circular metal plate, its controller does not use $x^{2}+y^{2}=r^{2}$; it drives the cutting head by computing $x=a+r\cos\theta$, $y=b+r\sin\theta$ as $\theta$ increases steadily with time. The parameter $\theta$ plays the role of a clock: it tells the machine \emph{where to be at each instant}. The same idea steers satellite orbits, car sat-nav curves and the rounded letters on this page.
\end{savaistu}

\begin{exercice}[Practice]
\begin{enumerate}
\item Write down the centre and radius of each circle given parametrically:
\begin{enumerate}
\item $x=4\cos\theta,\;y=4\sin\theta$;
\item $x=-1+3\cos\theta,\;y=5+3\sin\theta$.
\end{enumerate}
\item Convert to Cartesian form: $x=2+7\cos\theta$, $y=-3+7\sin\theta$.
\item Write parametric equations for the circle $x^{2}+y^{2}+6x-2y+6=0$.
\item A point $P$ moves on the circle $x^{2}+y^{2}=1$. Find the greatest and least values of the distance $PQ$, where $Q=(3,4)$.
\item $P=(\cos\theta,\sin\theta)$ moves on the unit circle and $M$ is the midpoint of $P$ and $(4,0)$. Find the Cartesian equation of the locus of $M$.
\item Find the equation of the tangent to the circle $x^{2}+y^{2}=16$ at the point corresponding to $\theta=\dfrac{\pi}{6}$.
\item A circle has parametric equations $x=1+2\cos\theta$, $y=-2+2\sin\theta$. Find the Cartesian equation of the tangent at the point where $\theta=\dfrac{\pi}{4}$.
\end{enumerate}
\end{exercice}

\begin{corrige}[Exercise 1]
\begin{enumerate}
\item (a) Centre $(0,0)$, radius $4$. (b) Centre $(-1,5)$, radius $3$.
\item $\cos\theta=\dfrac{x-2}{7}$, $\sin\theta=\dfrac{y+3}{7}$; hence $(x-2)^{2}+(y+3)^{2}=49$.
\item Completing the square: $(x+3)^{2}+(y-1)^{2}=4$, so $x=-3+2\cos\theta$, $y=1+2\sin\theta$.
\item $PQ^{2}=(\cos\theta-3)^{2}+(\sin\theta-4)^{2}=26-6\cos\theta-8\sin\theta$. Write $6\cos\theta+8\sin\theta=10\cos(\theta-\alpha)$ with $\cos\alpha=\frac{3}{5}$, $\sin\alpha=\frac{4}{5}$. Then $PQ^{2}_{\max}=36$, $PQ^{2}_{\min}=16$, so $PQ_{\max}=6$ and $PQ_{\min}=4$.
\item $M=\left(\dfrac{\cos\theta+4}{2},\dfrac{\sin\theta}{2}\right)$; eliminating $\theta$ gives $(2X-4)^{2}+(2Y)^{2}=1$, i.e.\ $(X-2)^{2}+Y^{2}=\tfrac14$: a circle of centre $(2,0)$, radius $\tfrac12$.
\item For $x^{2}+y^{2}=16$, tangent at $\theta$ is $x\cos\theta+y\sin\theta=4$. With $\theta=\frac{\pi}{6}$: $\frac{\sqrt{3}}{2}x+\frac{1}{2}y=4 \Rightarrow \sqrt{3}x+y=8$.
\item Circle centre $(1,-2)$, radius $2$. Tangent at $\theta$: $(x-1)\cos\theta+(y+2)\sin\theta=2$. With $\theta=\frac{\pi}{4}$: $\frac{\sqrt{2}}{2}(x-1)+\frac{\sqrt{2}}{2}(y+2)=2 \Rightarrow x+y+1=2\sqrt{2}$.
\end{enumerate}
\end{corrige}

With the parametric form in hand, we now return to Cartesian methods to study how a line meets a circle — the subject of Section 6 — where the number of intersection points will be decided by a discriminant condition.

\coursec{Section 6: Intersection of a Line and a Circle}

\courssub{Introduction: Intersection of a Line and a Circle — an overview}

In Section 5, we described a circle using parametric equations $x = a + r\cos\theta$, $y = b + r\sin\theta$. That description tells us \emph{where the points of the circle are}. We now ask a new question: \emph{where does a given straight line meet a given circle?} Think of a football field drawn on the ground at your school in Bamenda: the centre circle is painted, and the halfway line is a straight line crossing it. Where exactly does the halfway line enter and leave the circle? Answering that question precisely is the goal of this section. We shall see that the whole story is controlled by a single number: the \cle{discriminant} of a quadratic equation.

\courssub{Finding the Points of Intersection by Substitution}

\begin{experience}[A first encounter]
Consider the circle $x^{2} + y^{2} = 25$ (centre the origin, radius $5$) and the line $y = x + 1$. A point lies on \textbf{both} curves exactly when its coordinates satisfy \textbf{both} equations at the same time. So we substitute $y = x+1$ into the circle equation:
\[
x^2 + (x+1)^2 = 25 \;\Longrightarrow\; 2x^2 + 2x - 24 = 0 \;\Longrightarrow\; x^2 + x - 12 = 0.
\]
Factorising, $(x-3)(x+4) = 0$, so $x = 3$ or $x = -4$. Using $y = x+1$, the intersection points are $(3, 4)$ and $(-4, -3)$. Two points: the line is a \cle{secant}.
\end{experience}

This example contains the whole method. Geometrically, a line and a circle can meet in \textbf{two points} (the line is a \cle{secant}), in \textbf{one point} (the line is a \cle{tangent}), or in \textbf{no point at all} (the line misses the circle). Algebraically, substitution always produces a \textbf{quadratic equation} in one variable, and the three geometric situations correspond exactly to the three possible signs of its discriminant $\Delta = b^{2} - 4ac$.

\begin{methode}[Intersection of a line and a circle]
Given a circle $S = 0$ and a line $L = 0$:
\begin{enumerate}
\item From the line equation, express one variable in terms of the other (e.g.\ $y = mx + c$).
\item Substitute into the circle equation and simplify to a quadratic $ax^{2} + bx + c = 0$.
\item Compute the discriminant $\Delta = b^{2} - 4ac$.
\item Classify: $\Delta > 0$: two distinct intersection points (secant); $\Delta = 0$: one repeated point (tangent); $\Delta < 0$: no real intersection (line misses the circle).
\item If $\Delta \geq 0$, solve the quadratic and find the corresponding values of the other variable to obtain the coordinates of the intersection points.
\end{enumerate}
\end{methode}

\begin{popfigure}
\begin{center}
\begin{tikzpicture}
\begin{axis}[
  axis equal image,
  axis lines = middle,
  xlabel={$x$}, ylabel={$y$},
  xmin = -7, xmax = 7, ymin = -7, ymax = 7,
  xtick = {-6,-4,-2,2,4,6}, ytick = {-6,-4,-2,2,4,6},
  width = 12cm,
  clip = false
]
\draw[popblue, thick] (axis cs:0,0) circle [radius = 4];
\addplot[poporange, thick, domain=-6.5:6.5] {x + 1};
\addplot[popgreen, thick, domain=-1.5:6.5] {-x + 5.657};
\addplot[poppurple, thick, domain=-6.5:6.5] {0.5*x + 6.5};
\node[poporange, align=center] at (axis cs:-5.4,-3.1) {secant\\ $\Delta > 0$};
\node[popgreen, align=center] at (axis cs:5.6,1.2) {tangent\\ $\Delta = 0$};
\node[poppurple, align=center] at (axis cs:-4.6,5.6) {no meeting\\ $\Delta < 0$};
\end{axis}
\end{tikzpicture}
\legende{The circle $x^{2} + y^{2} = 16$ with a secant (two points), a tangent (one point) and a line that misses the circle (no point).}
\end{center}
\end{popfigure}

\begin{exoresolu}[Classifying three lines]
For the circle $x^{2} + y^{2} = 25$, determine the nature of the intersection with each line, and find the intersection points when they exist:
\begin{enumerate}
\item[(a)] $y = x + 1$;
\item[(b)] $3x + 4y = 25$;
\item[(c)] $y = x + 8$.
\end{enumerate}
\tcblower
\textbf{(a)} Substituting $y = x+1$ into $x^{2} + y^{2} = 25$:
\[
x^2 + (x+1)^2 = 25 \;\Longrightarrow\; 2x^2 + 2x - 24 = 0 \;\Longrightarrow\; x^2 + x - 12 = 0.
\]
$\Delta = 1 + 48 = 49 > 0$: \textbf{secant}. Roots $x = 3$ and $x = -4$, giving the points $(3, 4)$ and $(-4, -3)$.

\textbf{(b)} From $3x + 4y = 25$, $y = \dfrac{25 - 3x}{4}$. Substituting:
\[
x^2 + \frac{(25-3x)^2}{16} = 25
\;\Longrightarrow\; 16x^2 + 9x^2 - 150x + 625 = 400
\;\Longrightarrow\; 25x^2 - 150x + 225 = 0,
\]
i.e.\ $x^{2} - 6x + 9 = 0$, so $(x-3)^{2} = 0$. $\Delta = 0$: the line is a \textbf{tangent}, touching the circle at the single point $(3, 4)$.

\textbf{(c)} Substituting $y = x + 8$:
\[
x^2 + (x+8)^2 = 25 \;\Longrightarrow\; 2x^2 + 16x + 39 = 0.
\]
$\Delta = 256 - 312 = -56 < 0$: \textbf{no real intersection}; the line misses the circle completely.
\end{exoresolu}

\begin{attention}[Check before you conclude]
A negative discriminant does \textbf{not} mean you made an algebra error — it is a genuine geometric answer: the line and the circle simply do not meet. But do check your substitution arithmetic first; a single sign slip can flip the sign of $\Delta$.
\end{attention}

\courssub{The Distance Method: Comparing $d$ with $r$}

There is a second, purely geometric way to classify the intersection, which is often faster. Compute the perpendicular distance $d$ from the centre $C(a, b)$ of the circle to the line $px + qy + s = 0$:
\[
d = \frac{|pa + qb + s|}{\sqrt{p^2 + q^2}}.
\]
Then compare $d$ with the radius $r$:
\begin{itemize}
\item $d < r$: the line cuts the circle in \textbf{two points} (secant);
\item $d = r$: the line \textbf{touches} the circle (tangent);
\item $d > r$: the line \textbf{misses} the circle.
\end{itemize}
This is exactly the algebraic classification in disguise: one can prove that $\Delta$ (after substitution) has the same sign as $r^{2} - d^{2}$, up to a positive factor.

\begin{exemplebox}[Distance method in action]
Circle $(x-1)^{2} + (y-2)^{2} = 9$, so $C = (1, 2)$, $r = 3$. Line: $4x + 3y = 5$, i.e.\ $4x + 3y - 5 = 0$.
\[
d = \frac{|4(1) + 3(2) - 5|}{\sqrt{16+9}} = \frac{|10 - 5|}{5} = 1.
\]
Since $d = 1 < 3 = r$, the line is a secant: two intersection points.
\end{exemplebox}

\courssub{Chords: Length and Midpoint}

\begin{propriete}[The perpendicular from the centre bisects the chord]
Let $AB$ be a chord of a circle with centre $M$, and let $N$ be the foot of the perpendicular from $M$ to the line $AB$. Then $N$ is the \textbf{midpoint} of $AB$: the perpendicular from the centre of a circle to a chord bisects the chord. (Proof: triangles $MNA$ and $MNB$ are right-angled at $N$, share the side $MN$, and $MA = MB = r$; by Pythagoras $NA = NB$.)
\end{propriete}

\begin{propriete}[Length of a chord]
If a line is at perpendicular distance $d$ from the centre of a circle of radius $r$ (with $d < r$), then the length of the chord it cuts off is
\[
AB = 2\sqrt{r^2 - d^2}.
\]
This follows at once from Pythagoras in the right triangle $MNA$: $NA^{2} = MA^{2} - MN^{2} = r^{2} - d^{2}$, and $AB = 2\,NA$.
\end{propriete}

\begin{popfigure}
\begin{center}
\begin{tikzpicture}[scale = 1.3]
\coordinate (M) at (0,0);
\coordinate (A) at (-2.4,-1.4);
\coordinate (B) at (2.4,-1.4);
\coordinate (N) at (0,-1.4);
\draw[popblue, thick] (M) circle [radius = 2.777];
\draw[poporange, very thick] (A) -- (B);
\draw[popdark, thick] (M) -- (A);
\draw[popgreen, thick] (M) -- (N);
\draw[popmuted] (N) ++(0,0.22) -- ++(0.22,0.22) -- ++(0.22,0);
\fill (M) circle (1.5pt) node[above] {$M$};
\fill (A) circle (1.5pt) node[below left] {$A$};
\fill (B) circle (1.5pt) node[below right] {$B$};
\fill (N) circle (1.5pt) node[below] {$N$};
\node[popdark] at (-1.15,0.55) {$r$};
\node[popgreen] at (0.25,-0.7) {$d$};
\node[poporange] at (1.2,-1.15) {$\sqrt{r^{2} - d^{2}}$};
\end{tikzpicture}
\legende{The right triangle $MNA$: radius $r$, distance $d$ from centre to chord, and half-chord $\sqrt{r^{2} - d^{2}}$.}
\end{center}
\end{popfigure}

\begin{exoresolu}[Length of a chord]
The line $3x + 4y = 20$ cuts the circle $x^{2} + y^{2} = 25$. Find the length of the chord cut off, and the coordinates of its midpoint.
\tcblower
The circle has centre $M = (0,0)$ and radius $r = 5$. Distance from $M$ to the line $3x + 4y - 20 = 0$:
\[
d = \frac{|{-20}|}{\sqrt{9 + 16}} = \frac{20}{5} = 4.
\]
Chord length:
\[
AB = 2\sqrt{r^2 - d^2} = 2\sqrt{25 - 16} = 2 \times 3 = 6.
\]
The midpoint $N$ is the foot of the perpendicular from $M$ to the line. The perpendicular through $(0,0)$ has direction $(3, 4)$, so points on it are $(3t, 4t)$. Substituting into $3x + 4y = 20$: $9t + 16t = 20$, so $t = \dfrac{4}{5}$, giving $N = \left(\dfrac{12}{5}, \dfrac{16}{5}\right)$.
\end{exoresolu}

\begin{attention}[Do not confuse $d$ with the chord length]
The number $d$ is the distance from the \textbf{centre} to the \textbf{line}, not the distance between the intersection points. The chord length is $2\sqrt{r^{2} - d^{2}}$, and the \emph{half}-chord is $\sqrt{r^{2} - d^{2}}$. Forgetting the factor $2$ is one of the most frequent errors at the GCE.
\end{attention}

\courssub{Family of Circles Through the Intersection of a Circle and a Line}

Suppose the circle $S = 0$ and the line $L = 0$ meet at two points $P$ and $Q$. We want \textbf{all} circles passing through $P$ and $Q$ without solving for $P$ and $Q$ explicitly.

\begin{propriete}[The family $S + \lambda L = 0$]
Let $S = x^{2} + y^{2} + 2gx + 2fy + c = 0$ be a circle and $L = px + qy + s = 0$ a line which intersect it. Then for every real number $\lambda$, the equation
\[
S + \lambda L = 0
\]
represents a circle passing through the intersection points of $S = 0$ and $L = 0$. Conversely, every circle through those two points (except $S$ itself, in some formulations) is obtained for a suitable $\lambda$. \emph{Justification (examinable):} $S + \lambda L = 0$ has the form $x^{2} + y^{2} + \text{(linear terms)} = 0$, so it is a circle. If $(x_{0}, y_{0})$ lies on both $S = 0$ and $L = 0$, then $S(x_{0},y_{0}) + \lambda L(x_{0},y_{0}) = 0 + \lambda \cdot 0 = 0$, so the point lies on every member of the family.
\end{propriete}

\begin{attention}[Normalise $S$ first]
Before writing $S + \lambda L = 0$, make sure the circle equation is written with \textbf{coefficient $1$ for both $x^{2}$ and $y^{2}$}. If you are given $2x^{2} + 2y^{2} + 4x - 3 = 0$, divide by $2$ first. Otherwise $S + \lambda L$ will not have equal coefficients of $x^{2}$ and $y^{2}$ and will not even be a circle.
\end{attention}

\begin{methode}[Using the family $S + \lambda L = 0$]
\begin{enumerate}
\item Write $S = 0$ with unit coefficients of $x^{2}$ and $y^{2}$, and $L = 0$ in the form $px + qy + s = 0$.
\item Write the family $S + \lambda L = 0$.
\item Use the extra condition (the circle passes through a given point, has a given radius, a given centre line, etc.) to form an equation in $\lambda$.
\item Solve for $\lambda$ and substitute back to obtain the required circle.
\end{enumerate}
\end{methode}

\begin{exoresolu}[Circle through the intersection of a circle and a line]
Find the equation of the circle which passes through the points of intersection of the circle $x^{2} + y^{2} - 4x - 2y - 4 = 0$ and the line $x + 2y - 4 = 0$, and which also passes through the point $(2, 3)$.
\tcblower
The family of circles through the intersection is
\[
x^2 + y^2 - 4x - 2y - 4 + \lambda(x + 2y - 4) = 0.
\]
The point $(2, 3)$ lies on the required circle, so substitute $x = 2$, $y = 3$:
\[
4 + 9 - 8 - 6 - 4 + \lambda(2 + 6 - 4) = 0
\;\Longrightarrow\; -5 + 4\lambda = 0
\;\Longrightarrow\; \lambda = \frac{5}{4}.
\]
Substituting back and multiplying by $4$ to clear fractions:
\[
4x^2 + 4y^2 - 16x - 8y - 16 + 5x + 10y - 20 = 0,
\]
\[
\boxed{4x^2 + 4y^2 - 11x + 2y - 36 = 0.}
\]
\emph{Check:} the coefficients of $x^{2}$ and $y^{2}$ are equal, so this is indeed a circle, and one may verify that $(2,3)$ satisfies it: $16 + 36 - 22 + 6 - 36 = 0$. \checkmark
\end{exoresolu}

\begin{exoresolu}[Determining $\lambda$ from a radius condition]
Find the circle through the intersection of $x^{2} + y^{2} = 4$ and the line $x - y = 0$ which has radius $2$.
\tcblower
Family: $x^{2} + y^{2} - 4 + \lambda(x - y) = 0$. Completing the square:
\[
\left(x + \frac{\lambda}{2}\right)^2 + \left(y - \frac{\lambda}{2}\right)^2 = 4 + \frac{\lambda^2}{2}.
\]
The radius squared is $4 + \dfrac{\lambda^{2}}{2}$. Setting this equal to $2^{2} = 4$ gives $\lambda^{2} = 0$, so $\lambda = 0$: the only such circle is the original circle $x^{2} + y^{2} = 4$ itself. (Any other member of the family has a strictly larger radius — a nice geometric remark: the common chord is a diameter of the smallest circle of the family.)
\end{exoresolu}

\courssub{Consolidation Exercises}

\begin{exercice}[Exercise 1]
Find the coordinates of the points where the line $y = 2x - 3$ meets the circle $x^{2} + y^{2} - 4x + 2y - 5 = 0$.
\end{exercice}

\begin{exercice}[Exercise 2]
Show that the line $x + y = 5$ is a tangent to the circle $x^{2} + y^{2} - 2x - 4y + 3 = 0$, and find the point of contact.
\end{exercice}

\begin{exercice}[Exercise 3]
The line $4x - 3y = 10$ cuts the circle $x^{2} + y^{2} - 2x + 4y - 20 = 0$ at $A$ and $B$. Find the length $AB$ and the midpoint of $AB$.
\end{exercice}

\begin{exercice}[Exercise 4]
Find the equation of the circle passing through the intersection of $x^{2} + y^{2} - 6x - 8y + 9 = 0$ and the line $x - y + 1 = 0$, and through the point $(1, 2)$.
\end{exercice}

\begin{corrige}[Exercise 1]
Substitute $y = 2x - 3$: $x^{2} + (2x-3)^{2} - 4x + 2(2x-3) - 5 = 0$ gives $5x^{2} - 12x - 2 = 0$... expanding carefully: $x^{2} + 4x^{2} - 12x + 9 - 4x + 4x - 6 - 5 = 5x^{2} - 12x - 2 = 0$. Then $x = \dfrac{12 \pm \sqrt{144 + 40}}{10} = \dfrac{12 \pm \sqrt{184}}{10} = \dfrac{6 \pm \sqrt{46}}{5}$. The points are $\left(\dfrac{6 + \sqrt{46}}{5}, \dfrac{-3 + 2\sqrt{46}}{5}\right)$ and $\left(\dfrac{6 - \sqrt{46}}{5}, \dfrac{-3 - 2\sqrt{46}}{5}\right)$.
\end{corrige}

\begin{corrige}[Exercise 2]
Circle: centre $(1, 2)$, radius $\sqrt{1 + 4 - 3} = \sqrt{2}$. Distance from $(1,2)$ to $x + y - 5 = 0$: $d = \dfrac{|1 + 2 - 5|}{\sqrt{2}} = \dfrac{2}{\sqrt{2}} = \sqrt{2} = r$. Hence tangent. Point of contact: foot of perpendicular from $(1,2)$ along direction $(1,1)$: $(1 + t, 2 + t)$ on the line gives $1 + t + 2 + t = 5$, $t = 1$; contact point $(2, 3)$.
\end{corrige}

\begin{corrige}[Exercise 3]
Centre $(1, -2)$, radius $\sqrt{1 + 4 + 20} = 5$. Distance to line $4x - 3y - 10 = 0$: $d = \dfrac{|4 + 6 - 10|}{5} = 0$. The line passes through the centre, so the chord is a diameter: $AB = 10$, midpoint $= (1, -2)$.
\end{corrige}

\begin{corrige}[Exercise 4]
Family: $x^{2} + y^{2} - 6x - 8y + 9 + \lambda(x - y + 1) = 0$. At $(1, 2)$: $1 + 4 - 6 - 16 + 9 + \lambda(1 - 2 + 1) = -8 + 0 = -8 \neq 0$. Since $L(1,2) = 0$ but $S(1,2) = -8 \neq 0$, no value of $\lambda$ works: the point $(1,2)$ lies on the line $L$ but not on the circle $S$, so no circle of the family passes through it. The problem as stated has no solution — always check that the given point does not lie on the line $L = 0$ outside the circle!
\end{corrige}

\begin{aretenir}[Key points of Section 6]
\begin{itemize}
\item Substituting the line equation into the circle equation gives a quadratic; its discriminant classifies the intersection: $\Delta > 0$ secant, $\Delta = 0$ tangent, $\Delta < 0$ no meeting.
\item Equivalently, compare the distance $d$ from the centre to the line with the radius $r$: $d < r$, $d = r$, $d > r$.
\item The perpendicular from the centre to a chord bisects the chord; the chord length is $2\sqrt{r^{2} - d^{2}}$.
\item The family of circles through the intersection of $S = 0$ and $L = 0$ is $S + \lambda L = 0$ — with $S$ normalised so that $x^{2}$ and $y^{2}$ both have coefficient $1$.
\end{itemize}
\end{aretenir}

\coursec{Section 7: Intersection of Two Circles, Families of Circles and the Radical Axis}

In Section 6 we studied how a \emph{line} meets a circle: we substituted the equation of the line into the equation of the circle and analysed the resulting quadratic. We now go one step further and ask: what happens when \emph{two circles} meet? The beautiful surprise is that the hard work is again done by a straight line. By simply \emph{subtracting} the equations of the two circles, we obtain a line which contains their common points — the \cle{common chord}. This single idea unlocks three powerful tools of the GCE syllabus: the intersection points of two circles, the \cle{family of circles} through those points, and the \cle{radical axis}.

Throughout this section we write a circle in general form
\[
S:\; x^{2}+y^{2}+2gx+2fy+c=0,
\]
and we use the shorthand $S(x,y)=0$ or simply $S=0$. Recall that the centre is $(-g,-f)$ and the radius is $\sqrt{g^{2}+f^{2}-c}$.

\courssub{Points of Intersection of Two Circles and the Common Chord}

\textbf{Observation.} Let two circles
\[
S_{1}:\; x^{2}+y^{2}+2g_{1}x+2f_{1}y+c_{1}=0,
\qquad
S_{2}:\; x^{2}+y^{2}+2g_{2}x+2f_{2}y+c_{2}=0
\]
intersect at points $A$ and $B$. If $(x,y)$ is a common point, then both equations hold, so their \emph{difference} also holds:
\[
S_{1}-S_{2}:\; 2(g_{1}-g_{2})x+2(f_{1}-f_{2})y+(c_{1}-c_{2})=0.
\]
The $x^{2}$ and $y^{2}$ terms cancel, leaving the equation of a \emph{straight line}. Since $A$ and $B$ both satisfy $S_{1}-S_{2}=0$, this line passes through both intersection points: it is the line of the \cle{common chord} $AB$.

\begin{methode}[Intersecting two circles]
To find the points of intersection of $S_{1}=0$ and $S_{2}=0$:
\begin{enumerate}
\item Write both equations in general form and compute $S_{1}-S_{2}=0$: this is the line of the common chord.
\item Express $y$ (or $x$) from the chord equation and substitute into \emph{either} circle.
\item Solve the resulting quadratic for the remaining coordinate.
\item Write down both points of intersection and check them in the other circle.
\end{enumerate}
\end{methode}

\begin{exoresolu}[Intersection of two circles]
Find the points of intersection of the circles
\[
S_{1}:\; x^{2}+y^{2}-4x-2y-4=0
\qquad\text{and}\qquad
S_{2}:\; x^{2}+y^{2}+2x-6y+6=0.
\]
\tcblower
\textbf{Step 1: common chord.} Subtracting,
\[
S_{1}-S_{2}:\; -6x+4y-10=0
\quad\Longleftrightarrow\quad
3x-2y+5=0,
\qquad\text{so } y=\tfrac{3x+5}{2}.
\]
\textbf{Step 2: substitute into $S_{1}$.}
\[
x^{2}+\left(\tfrac{3x+5}{2}\right)^{2}-4x-2\left(\tfrac{3x+5}{2}\right)-4=0.
\]
Multiplying by $4$:
\[
4x^{2}+(9x^{2}+30x+25)-16x-(12x+20)-16=0
\;\Longrightarrow\;
13x^{2}+2x-11=0.
\]
\textbf{Step 3: solve.} $13x^{2}+2x-11=0$ factorises as $(13x-11)(x+1)=0$, giving $x=-1$ or $x=\tfrac{11}{13}$.

For $x=-1$: $y=\dfrac{-3+5}{2}=1$. For $x=\tfrac{11}{13}$: $y=\dfrac{1}{2}\left(\tfrac{33}{13}+5\right)=\dfrac{49}{13}$.

\textbf{Conclusion.} The circles meet at $\left(-1,\,1\right)$ and $\left(\dfrac{11}{13},\,\dfrac{49}{13}\right)$. Both points satisfy $S_{2}=0$, as required.
\end{exoresolu}

\begin{popfigure}
\begin{center}
\begin{tikzpicture}[scale=0.9]
\draw[popline,->] (-2.6,0)--(4.6,0) node[right]{\scriptsize $x$};
\draw[popline,->] (0,-2.4)--(0,3.4) node[above]{\scriptsize $y$};
\draw[popblue,thick] (0,0) circle (2);
\draw[poporange,thick] (2.2,0.6) circle (1.7);
\coordinate (A) at (1.02,1.72);
\coordinate (B) at (1.02,-1.72);
\draw[popdark,dashed] (0,0)--(2.2,0.6);
\draw[poppurple,very thick] (A)--(B);
\fill[popdark] (A) circle (1.6pt) node[above right]{\scriptsize $A$};
\fill[popdark] (B) circle (1.6pt) node[below right]{\scriptsize $B$};
\fill[popblue] (0,0) circle (1.6pt) node[below left]{\scriptsize $C_{1}$};
\fill[poporange] (2.2,0.6) circle (1.6pt) node[above right]{\scriptsize $C_{2}$};
\coordinate (M) at (1.02,0);
\draw[popgreen] (M) ++(0,0.18) -- ++(0.18,0.18) -- ++(0.18,0);
\node[poppurple] at (0.6,1.3) {\scriptsize chord};
\node[popdark] at (1.7,0.28) {\scriptsize centres};
\end{tikzpicture}
\end{center}
\legende{Two intersecting circles: the common chord $AB$ is perpendicular to the line of centres $C_{1}C_{2}$.}
\end{popfigure}

\courssub{The Family of Circles Through the Intersection of Two Circles}

\textbf{Observation.} Suppose we want \emph{the} circle through the two intersection points of $S_{1}=0$ and $S_{2}=0$ which also passes through a third given point. Solving from scratch (three points, three unknowns $g,f,c$) is long. There is a far more elegant route: build a whole \emph{family} of circles through $A$ and $B$ with one parameter, then choose the parameter to fit the third point.

\begin{propriete}[Family of circles through two intersection points — proof examinable]
Let $S_{1}=0$ and $S_{2}=0$ be two intersecting circles. Then for every real $\lambda\neq -1$, the equation
\[
S_{1}+\lambda S_{2}=0
\]
represents a circle passing through all the common points of $S_{1}$ and $S_{2}$. Conversely, every circle through those common points (other than $S_{2}$ itself) has this form for some $\lambda$.
\end{propriete}

\textbf{Proof.} First, $S_{1}+\lambda S_{2}=0$ is
\[
(1+\lambda)x^{2}+(1+\lambda)y^{2}+2(g_{1}+\lambda g_{2})x+2(f_{1}+\lambda f_{2})y+(c_{1}+\lambda c_{2})=0.
\]
Since $\lambda\neq -1$, we may divide by $1+\lambda$: the coefficients of $x^{2}$ and $y^{2}$ become $1$ and there is no $xy$ term, so this is indeed a circle. Secondly, if $P$ is a common point of the two circles, then $S_{1}(P)=0$ and $S_{2}(P)=0$, hence $S_{1}(P)+\lambda S_{2}(P)=0$: the point $P$ lies on every member of the family. $\blacksquare$

\begin{attention}[The case $\lambda=-1$ and the missing circle]
When $\lambda=-1$, the quadratic terms vanish and $S_{1}-S_{2}=0$ is the \emph{common chord} (a line, not a circle). Also, the family $S_{1}+\lambda S_{2}=0$ never produces $S_{2}$ itself (that would require $1+\lambda=0$). If a problem's answer turns out to be $S_{2}$, check it separately.
\end{attention}

\begin{methode}[Circle through the intersection of two circles and a given point]
\begin{enumerate}
\item Write the family $S_{1}+\lambda S_{2}=0$.
\item Substitute the coordinates of the given point: this gives a linear equation in $\lambda$.
\item Solve for $\lambda$ (reject $\lambda=-1$), substitute back and simplify.
\end{enumerate}
\end{methode}

\begin{exoresolu}[Circle through two intersection points and a given point]
Find the equation of the circle which passes through the points of intersection of
\[
S_{1}:\; x^{2}+y^{2}-4x-2y-4=0,
\qquad
S_{2}:\; x^{2}+y^{2}+2x-6y+6=0
\]
and through the point $(1,2)$.
\tcblower
The family is $S_{1}+\lambda S_{2}=0$:
\[
(1+\lambda)x^{2}+(1+\lambda)y^{2}+(-4+2\lambda)x+(-2-6\lambda)y+(-4+6\lambda)=0.
\]
Substituting $(1,2)$: note $S_{1}(1,2)=1+4-4-4-4=-7$ and $S_{2}(1,2)=1+4+2-12+6=1$. Hence
\[
-7+\lambda(1)=0 \quad\Longrightarrow\quad \lambda=7.
\]
Substituting back:
\[
8x^{2}+8y^{2}+10x-44y+38=0
\quad\Longleftrightarrow\quad
4x^{2}+4y^{2}+5x-22y+19=0.
\]
\textbf{Check:} at $(1,2)$: $4+16+5-44+19=0$. \checkmark
\end{exoresolu}

\begin{popfigure}
\begin{center}
\begin{tikzpicture}
\begin{axis}[
  width=9.2cm, height=6.4cm,
  axis equal image,
  axis lines=middle,
  xlabel={\scriptsize $x$}, ylabel={\scriptsize $y$},
  xmin=-4.5, xmax=4.5, ymin=-3.2, ymax=3.6,
  xtick=\empty, ytick=\empty,
  tick style={draw=none},
  axis line style={popline},
]
\addplot[popblue, thick, domain=0:360, samples=120] ({2*cos(x)-1},{2*sin(x)});
\addplot[poporange, thick, domain=0:360, samples=120] ({1.6*cos(x)+1},{1.6*sin(x)});
\addplot[poppurple, thick, domain=0:360, samples=120] ({2.9*cos(x)+0.05},{2.9*sin(x)+0.9});
\addplot[popgreen, thick, domain=0:360, samples=120] ({1.35*cos(x)-0.05},{1.35*sin(x)-0.55});
\addplot[only marks, mark=*, mark size=1.6pt, popdark] coordinates {(-0.05,1.99) (-0.05,-1.99)};
\node[popdark, anchor=west] at (axis cs:0.15,2.1) {\scriptsize $A$};
\node[popdark, anchor=west] at (axis cs:0.15,-2.3) {\scriptsize $B$};
\end{axis}
\end{tikzpicture}
\end{center}
\legende{Members of the family $S_{1}+\lambda S_{2}=0$: every circle passes through the same two points $A$ and $B$.}
\end{popfigure}

\courssub{The Radical Axis of Two Circles}

We now give the line $S_{1}-S_{2}=0$ a name and a meaning that remain valid \emph{even when the circles do not intersect}.

Recall from Section 4 that the length of the tangent from an external point $P(x_{1},y_{1})$ to the circle $S=0$ is $\sqrt{S(x_{1},y_{1})}$, provided $S$ is written with unit coefficients of $x^{2}$ and $y^{2}$.

\begin{definition}[Radical axis]
The \cle{radical axis} of two circles $S_{1}=0$ and $S_{2}=0$ is the locus of points $P$ from which the tangent lengths to the two circles are equal:
\[
\sqrt{S_{1}(P)}=\sqrt{S_{2}(P)}
\quad\Longleftrightarrow\quad
S_{1}(P)=S_{2}(P)
\quad\Longleftrightarrow\quad
S_{1}-S_{2}=0.
\]
Its equation is therefore
\[
2(g_{1}-g_{2})x+2(f_{1}-f_{2})y+(c_{1}-c_{2})=0.
\]
\end{definition}

\begin{propriete}[The radical axis is perpendicular to the line of centres — proof examinable]
For two non-concentric circles, the radical axis is perpendicular to the line joining their centres.
\end{propriete}

\textbf{Proof.} The centres are $C_{1}(-g_{1},-f_{1})$ and $C_{2}(-g_{2},-f_{2})$, so a direction vector of the line of centres is
\[
\vec{C_{1}C_{2}}=\begin{pmatrix} -(g_{2}-g_{1})\\ -(f_{2}-f_{1})\end{pmatrix}.
\]
The radical axis $2(g_{1}-g_{2})x+2(f_{1}-f_{2})y+(c_{1}-c_{2})=0$ has normal vector $\begin{pmatrix} g_{1}-g_{2}\\ f_{1}-f_{2}\end{pmatrix}$, which is parallel to $\vec{C_{1}C_{2}}$. A line whose normal is parallel to $C_{1}C_{2}$ is perpendicular to $C_{1}C_{2}$. $\blacksquare$

\begin{aretenir}[The radical axis in the three configurations]
\begin{itemize}
\item \textbf{Intersecting circles:} the radical axis is the line of the \cle{common chord}.
\item \textbf{Touching circles:} the radical axis is the \cle{common tangent} at the point of contact.
\item \textbf{Non-intersecting circles:} the radical axis is the line of points with equal tangent lengths; it lies between the circles, perpendicular to the line of centres.
\end{itemize}
\end{aretenir}

\begin{popfigure}
\begin{center}
\begin{tikzpicture}[scale=0.95]
\draw[popline,->] (-3.2,0)--(5.4,0) node[right]{\scriptsize $x$};
\draw[popline,->] (0,-2.6)--(0,2.8) node[above]{\scriptsize $y$};
\draw[popblue,thick] (-1.4,0) circle (1.3);
\draw[poporange,thick] (3.6,0) circle (1.6);
\draw[poppurple,very thick] (1.15,-2.4)--(1.15,2.6) node[above]{\scriptsize radical axis};
\fill[poppurple] (1.15,1.7) circle (1.8pt) node[right=2pt]{\scriptsize $P$};
\coordinate (T1) at (-0.45,1.14);
\coordinate (T2) at (2.55,1.42);
\draw[popgreen,thick] (1.15,1.7)--(T1);
\draw[popgreen,thick] (1.15,1.7)--(T2);
\fill[popgreen] (T1) circle (1.5pt) node[above left]{\scriptsize $T_{1}$};
\fill[popgreen] (T2) circle (1.5pt) node[above right]{\scriptsize $T_{2}$};
\fill[popblue] (-1.4,0) circle (1.6pt) node[below]{\scriptsize $C_{1}$};
\fill[poporange] (3.6,0) circle (1.6pt) node[below]{\scriptsize $C_{2}$};
\draw[popdark,dashed] (-1.4,0)--(3.6,0);
\draw[popdark] (1.15,0) ++(0,0.18) -- ++(0.18,0.18) -- ++(0.18,0);
\node[popgreen] at (0.1,2.15) {\scriptsize $PT_{1}=PT_{2}$};
\end{tikzpicture}
\end{center}
\legende{Non-intersecting circles: from any point $P$ of the radical axis, the tangent segments $PT_{1}$ and $PT_{2}$ have equal length.}
\end{popfigure}

\begin{exoresolu}[Radical axis and equal tangent lengths]
Find the radical axis of the circles
\[
S_{1}:\; x^{2}+y^{2}-2x-4y-4=0,
\qquad
S_{2}:\; x^{2}+y^{2}+4x-2y+1=0,
\]
and verify that the point $P(1,4)$ on it has equal tangent lengths to both circles.
\tcblower
\textbf{Radical axis:}
\[
S_{1}-S_{2}:\; -6x-2y-5=0
\quad\Longleftrightarrow\quad
6x+2y+5=0.
\]
\textbf{Check that $P$ lies on it:} $6(1)+2(4)+5=19\neq 0$ — so in fact $P(1,4)$ is \emph{not} on the axis. Take instead the point where the axis meets $x=1$: $2y=-11$, $y=-\tfrac{11}{2}$; call $Q\left(1,-\tfrac{11}{2}\right)$. Then
\[
S_{1}(Q)=1+\tfrac{121}{4}-2+22-4=\tfrac{189}{4},
\qquad
S_{2}(Q)=1+\tfrac{121}{4}+4+11+1=\tfrac{189}{4}.
\]
Hence the tangent lengths are equal: $\ell_{1}=\ell_{2}=\sqrt{\tfrac{189}{4}}=\tfrac{3\sqrt{21}}{2}$. This illustrates the defining property of the radical axis.
\end{exoresolu}

\begin{savaistu}[The radical centre of three circles]
Given three circles with non-collinear centres, the three radical axes taken in pairs are \emph{concurrent}: they meet at a single point called the \cle{radical centre}. Indeed, if $P$ satisfies $S_{1}=S_{2}$ and $S_{2}=S_{3}$, then automatically $S_{1}=S_{3}$. The radical centre is the unique point with equal tangent lengths to all three circles — a favourite one-line justification in GCE questions.
\end{savaistu}

\courssub{Locus Problems Involving Two Circles}

The radical axis answers a whole class of locus questions.

\begin{exoresolu}[Locus of points with equal tangent lengths]
A point $P$ moves so that the lengths of the tangents from $P$ to the circles
\[
x^{2}+y^{2}-6x-8y=0
\qquad\text{and}\qquad
x^{2}+y^{2}+2x-4y-3=0
\]
are equal. Find the locus of $P$.
\tcblower
Let $P(x,y)$. Equal tangent lengths means $S_{1}(P)=S_{2}(P)$:
\[
x^{2}+y^{2}-6x-8y = x^{2}+y^{2}+2x-4y-3
\;\Longrightarrow\;
8x+4y-3=0.
\]
The locus is the straight line $8x+4y-3=0$ — precisely the radical axis of the two circles.
\end{exoresolu}

A second classical type asks for a \emph{constant ratio or sum of squared distances}.

\begin{exoresolu}[Constant sum of squared distances]
Find the locus of a point $P$ such that the sum of the squares of its distances from $A(2,0)$ and $B(-2,0)$ is $20$.
\tcblower
Let $P(x,y)$. Then
\[
PA^{2}+PB^{2}=(x-2)^{2}+y^{2}+(x+2)^{2}+y^{2}=2x^{2}+2y^{2}+8.
\]
Setting this equal to $20$: $2x^{2}+2y^{2}=12$, i.e.
\[
x^{2}+y^{2}=6,
\]
a circle centred at the origin (the midpoint of $AB$) with radius $\sqrt{6}$.
\end{exoresolu}

\begin{attention}[Common mistakes in this section]
\begin{itemize}
\item Subtracting the equations when the coefficients of $x^{2}$ and $y^{2}$ are \emph{not both} $1$: first divide each circle through by the coefficient of $x^{2}$, otherwise $S_{1}-S_{2}$ is not the radical axis.
\item Forgetting that $S_{1}+\lambda S_{2}=0$ with $\lambda=-1$ gives a \emph{line}, not a circle: always state $\lambda\neq-1$.
\item Claiming the radical axis passes through the intersection points when the circles do \emph{not} intersect: it is still well defined, but it is no longer a chord.
\item Sign errors when computing $S_{1}-S_{2}$: subtract \emph{every} term, including the constants.
\end{itemize}
\end{attention}

\begin{exercice}
The circles $S_{1}: x^{2}+y^{2}-2x-6y+6=0$ and $S_{2}: x^{2}+y^{2}+4x+2y-4=0$ are given.
\begin{enumerate}
\item Find the equation of their radical axis.
\item Find the coordinates of their points of intersection.
\item Find the equation of the circle through these intersection points and the point $(0,2)$.
\end{enumerate}
\end{exercice}

\begin{corrige}[Exercise 1]
\textbf{1.} $S_{1}-S_{2}: -6x-8y+10=0$, i.e. $3x+4y-5=0$.

\textbf{2.} From the axis, $y=\dfrac{5-3x}{4}$. Substituting into $S_{1}$ and multiplying by $16$:
\[
16x^{2}+(25-30x+9x^{2})-32x-24(5-3x)+96=0
\]
\[
\Longrightarrow 25x^{2}+10x+1=0 \Longrightarrow (5x+1)^{2}=0.
\]
Thus $x=-\frac{1}{5}$ (double root). Then $y=\frac{5-3(-1/5)}{4}=\frac{5+3/5}{4}=\frac{28/5}{4}=\frac{7}{5}$.
The circles touch externally at $\left(-\frac{1}{5},\frac{7}{5}\right)$. (The discriminant is zero, confirming they touch.)

\textbf{3.} The family of circles through the common point (the point of contact) is $S_{1}+\lambda S_{2}=0$, $\lambda\neq-1$. Substituting $(0,2)$: $S_{1}(0,2)=4-12+6=-2$, $S_{2}(0,2)=4+4-4=4$, so $-2+4\lambda=0$, $\lambda=\frac{1}{2}$. Substituting back:
\[
S_{1}+\frac{1}{2}S_{2}=0 \;\Longrightarrow\; 2S_{1}+S_{2}=0 \;\Longrightarrow\; 3x^{2}+3y^{2}-10y+8=0.
\]
This circle passes through the point of contact $\left(-\frac{1}{5},\frac{7}{5}\right)$ and through $(0,2)$.
\end{corrige}

\begin{exercice}
Find the radical centre of the three circles
\[
x^{2}+y^{2}-4x-2y-4=0,\quad
x^{2}+y^{2}+2x-6y+6=0,\quad
x^{2}+y^{2}-8x+4y+4=0.
\]
\end{exercice}

\begin{corrige}[Exercise 2]
Radical axis of the first pair: $-6x+4y-10=0$, i.e. $3x-2y+5=0$.
Radical axis of circles 1 and 3: $4x-6y-8=0$, i.e. $2x-3y-4=0$.
Solving: from the first, $y=\tfrac{3x+5}{2}$; substituting: $2x-\tfrac{9x+15}{2}-4=0$, so $4x-9x-15-8=0$, $x=-\tfrac{23}{5}$, $y=\tfrac{1}{2}\left(-\tfrac{69}{5}+5\right)=-\tfrac{22}{5}$.
The radical centre is $\left(-\dfrac{23}{5},\,-\dfrac{22}{5}\right)$; one may verify it also lies on the radical axis of circles 2 and 3.
\end{corrige}

\begin{aretenir}[Key points of Section 7]
\begin{itemize}
\item Subtracting two circle equations gives a \emph{line}: the common chord (intersecting case), the common tangent (touching case), or the radical axis in general.
\item The family of circles through the intersections of $S_{1}=0$ and $S_{2}=0$ is $S_{1}+\lambda S_{2}=0$, $\lambda\neq-1$; find $\lambda$ by substituting a given point.
\item The radical axis $S_{1}-S_{2}=0$ is the locus of points with equal tangent lengths and is perpendicular to the line of centres.
\item The three radical axes of three circles concur at the radical centre.
\end{itemize}
\end{aretenir}

\coursec{Section 8: Conditions for Two Circles to Touch}

In Section 7 we studied what happens when two circles \emph{intersect}: we found their common points, the family of circles through those points, and the radical axis $S_1 - S_2 = 0$. But between ``two circles that cut each other in two points'' and ``two circles that do not meet at all'' there is a remarkable boundary case: the two circles meet in \emph{exactly one point}. We then say that the circles \cle{touch} each other, or that they are \cle{tangent} to one another. This section gives simple algebraic tests, in terms of the distance between the centres and the radii, to decide in which configuration two given circles lie, and how to find the point of contact and the common tangent when they touch. This is a classic GCE Advanced Level question.

\courssub{The Five Relative Positions of Two Circles}

Consider two circles $\mathcal{C}_1$ and $\mathcal{C}_2$, with centres $C_1$, $C_2$ and radii $r_1$, $r_2$. Let
\[ d = |C_1C_2| \]
be the distance between the centres. Imagine holding $\mathcal{C}_1$ fixed and sliding $\mathcal{C}_2$ towards it from far away. At first the circles are separate; then they just touch; then they overlap and cut in two points; then, if $\mathcal{C}_2$ is small enough, it slips inside $\mathcal{C}_1$, touches it from the inside, and finally lies entirely inside without touching. The whole geometry is controlled by comparing $d$ with the two numbers
\[ r_1 + r_2 \qquad \text{and} \qquad |r_1 - r_2|. \]

\begin{propriete}[Relative positions of two circles]
Let $\mathcal{C}_1(C_1, r_1)$ and $\mathcal{C}_2(C_2, r_2)$ be two circles and $d = |C_1C_2|$. Then:
\begin{enumerate}
\item $d > r_1 + r_2$: the circles are \cle{separate} (external to each other), no common point;
\item $d = r_1 + r_2$: the circles \cle{touch externally}, exactly one common point;
\item $|r_1 - r_2| < d < r_1 + r_2$: the circles \cle{intersect} in two distinct points;
\item $d = |r_1 - r_2|$ (with $d \neq 0$): the circles \cle{touch internally}, exactly one common point;
\item $d < |r_1 - r_2|$: one circle lies \cle{inside the other} without touching, no common point.
\end{enumerate}
(The degenerate case $d = 0$ with $r_1 = r_2$ gives two identical circles; if $d = 0$ and $r_1 \neq r_2$ the circles are concentric and have no common points.)
\end{propriete}

\begin{popfigure}
\begin{tikzpicture}[scale=0.45]
% 1: separate
\begin{scope}[shift={(0,0)}]
\draw[popblue, thick] (0,0) circle (1.5);
\draw[poporange, thick] (4.6,0) circle (1.2);
\fill (0,0) circle (0.09) node[below] {\scriptsize $C_1$};
\fill (4.6,0) circle (0.09) node[below] {\scriptsize $C_2$};
\draw[<->, popdark] (0,-1.9) -- (4.6,-1.9) node[midway, below] {\scriptsize $d > r_1+r_2$};
\node at (2.3,-2.9) {\scriptsize \textbf{Separate}};
\end{scope}
% 2: touch externally
\begin{scope}[shift={(8.6,0)}]
\draw[popblue, thick] (0,0) circle (1.5);
\draw[poporange, thick] (2.7,0) circle (1.2);
\fill (0,0) circle (0.09) node[below] {\scriptsize $C_1$};
\fill (2.7,0) circle (0.09) node[below] {\scriptsize $C_2$};
\fill[popdark] (1.5,0) circle (0.11);
\draw[<->, popdark] (0,-1.9) -- (2.7,-1.9) node[midway, below] {\scriptsize $d = r_1+r_2$};
\node at (1.35,-2.9) {\scriptsize \textbf{Touch externally}};
\end{scope}
% 3: intersect
\begin{scope}[shift={(16.4,0)}]
\draw[popblue, thick] (0,0) circle (1.5);
\draw[poporange, thick] (1.9,0) circle (1.2);
\fill (0,0) circle (0.09) node[below] {\scriptsize $C_1$};
\fill (1.9,0) circle (0.09) node[below] {\scriptsize $C_2$};
\draw[<->, popdark] (0,-1.9) -- (1.9,-1.9) node[midway, below] {\scriptsize $|r_1-r_2|<d<r_1+r_2$};
\node at (0.95,-2.9) {\scriptsize \textbf{Intersecting}};
\end{scope}
% 4: touch internally
\begin{scope}[shift={(24.2,0)}]
\draw[popblue, thick] (0,0) circle (1.8);
\draw[poporange, thick] (0.9,0) circle (0.9);
\fill (0,0) circle (0.09) node[below] {\scriptsize $C_1$};
\fill (0.9,0) circle (0.09) node[below left] {\scriptsize $C_2$};
\fill[popdark] (1.8,0) circle (0.11);
\draw[<->, popdark] (0,-2.2) -- (0.9,-2.2) node[midway, below] {\scriptsize $d = |r_1-r_2|$};
\node at (0.9,-3.2) {\scriptsize \textbf{Touch internally}};
\end{scope}
% 5: one inside the other
\begin{scope}[shift={(31.4,0)}]
\draw[popblue, thick] (0,0) circle (1.8);
\draw[poporange, thick] (0.4,0.3) circle (0.7);
\fill (0,0) circle (0.09) node[below] {\scriptsize $C_1$};
\fill (0.4,0.3) circle (0.09) node[above] {\scriptsize $C_2$};
\draw[<->, popdark] (0,-2.2) -- (0.4,-2.2) node[midway, below] {\scriptsize $d < |r_1-r_2|$};
\node at (0.2,-3.2) {\scriptsize \textbf{One inside the other}};
\end{scope}
\end{tikzpicture}
\legende{The five relative positions of two circles, classified by comparing $d$ with $r_1+r_2$ and $|r_1-r_2|$.}
\end{popfigure}

\courssub{Conditions for Two Circles to Touch}

\begin{propriete}[Touching conditions]
Two circles with centres $C_1, C_2$, radii $r_1, r_2$, and $d = |C_1C_2|$:
\begin{itemize}
\item \textbf{touch externally} if and only if $\;d = r_1 + r_2$;
\item \textbf{touch internally} if and only if $\;d = |r_1 - r_2|$ (with $d \neq 0$).
\end{itemize}
In both cases the two circles have \emph{exactly one} common point, called the \cle{point of contact}. (The proof of the ``if'' direction, by exhibiting the common point on the line of centres, is instructive and sketched below; the statement itself must be known for the examination.)
\end{propriete}

\textbf{Why these conditions are natural.} Suppose the circles touch externally at $T$. Since $T$ lies on both circles, $|C_1T| = r_1$ and $|C_2T| = r_2$, and $T$ lies between $C_1$ and $C_2$ on the line of centres, so
\[ d = |C_1C_2| = |C_1T| + |TC_2| = r_1 + r_2. \]
If they touch internally (say $\mathcal{C}_2$ inside $\mathcal{C}_1$), then $C_2$ and $T$ are on the same side of $C_1$ and $C_1, C_2, T$ are collinear with $T$ furthest from $C_1$, so
\[ d = |C_1C_2| = |C_1T| - |C_2T| = r_1 - r_2 = |r_1 - r_2|. \]
Conversely, if $d = r_1 + r_2$, the point $T$ on segment $[C_1C_2]$ with $|C_1T| = r_1$ satisfies $|TC_2| = d - r_1 = r_2$, so $T$ lies on both circles; any other common point $P$ would give a triangle $C_1PC_2$ with $|C_1P| + |PC_2| = r_1 + r_2 = d = |C_1C_2|$, forcing $P$ onto the segment $[C_1C_2]$ by the triangle inequality, hence $P = T$. The internal case is similar.

\begin{propriete}[The point of contact lies on the line of centres]
If two circles touch, their point of contact $T$ lies on the line joining their centres. Moreover $T$ divides $C_1C_2$ in the ratio of the radii:
\begin{itemize}
\item \textbf{external contact:} $T$ divides $[C_1C_2]$ \emph{internally} in the ratio $r_1 : r_2$, so
\[ T = \frac{r_2\,C_1 + r_1\,C_2}{r_1 + r_2}\quad\text{i.e.}\quad T = C_1 + \frac{r_1}{r_1+r_2}\,\overrightarrow{C_1C_2}; \]
\item \textbf{internal contact:} $T$ divides the line $C_1C_2$ \emph{externally} in the ratio $r_1 : r_2$; the vector formula
\[ T = C_1 + \frac{r_1}{r_1 - r_2}\,\overrightarrow{C_1C_2} \]
works in all cases (if $r_1 < r_2$ the coefficient is negative, placing $T$ on the ray opposite to $\overrightarrow{C_1C_2}$).
\end{itemize}
\end{propriete}

\begin{methode}[Finding the point of contact]
\begin{enumerate}
\item Compute the centres $C_1, C_2$ and radii $r_1, r_2$; compute $d = |C_1C_2|$.
\item Test $d = r_1 + r_2$ (external) and $d = |r_1 - r_2|$ (internal).
\item Use the vector formula:
   \begin{itemize}
   \item External: $T = C_1 + \dfrac{r_1}{r_1+r_2}\,\overrightarrow{C_1C_2}$.
   \item Internal: $T = C_1 + \dfrac{r_1}{r_1-r_2}\,\overrightarrow{C_1C_2}$ (this automatically handles whether $r_1 > r_2$ or $r_1 < r_2$).
   \end{itemize}
\end{enumerate}
\end{methode}

\begin{popfigure}
\begin{tikzpicture}[scale=1.05]
\draw[popblue, very thick] (0,0) circle (2);
\draw[poporange, very thick] (3.5,0) circle (1.5);
\fill (0,0) circle (0.07) node[below left] {$C_1$};
\fill (3.5,0) circle (0.07) node[below right] {$C_2$};
\fill[popdark] (2,0) circle (0.09) node[above right] {$T$};
\draw[popmuted, dashed, thick] (-2.4,0) -- (5.4,0) node[right] {\scriptsize line of centres};
\draw[popgreen, very thick] (2,-2.6) -- (2,2.6) node[above] {\scriptsize common tangent};
\draw[<->, popdark] (0,-0.35) -- (2,-0.35) node[midway, below] {\scriptsize $r_1$};
\draw[<->, popdark] (2,0.35) -- (3.5,0.35) node[midway, above] {\scriptsize $r_2$};
\draw[<->, popdark] (0,-2.35) -- (3.5,-2.35) node[midway, below] {\scriptsize $d = r_1 + r_2$};
\end{tikzpicture}
\legende{Two circles touching externally at $T$: $C_1, T, C_2$ are collinear, $d = r_1 + r_2$, and the common tangent at $T$ is perpendicular to the line of centres.}
\end{popfigure}

\courssub{The Common Tangent at the Point of Contact}

When two circles touch at $T$, the tangent to $\mathcal{C}_1$ at $T$ and the tangent to $\mathcal{C}_2$ at $T$ are both perpendicular to the line of centres $C_1C_2$ at the \emph{same} point $T$: they are therefore the \emph{same line}. Two touching circles always share a \cle{common tangent} at their point of contact.

Recall from Section 7 that the radical axis of two circles $S_1 = 0$ and $S_2 = 0$ is the line $S_1 - S_2 = 0$. When the circles intersect, this is the line through the two intersection points. When the circles touch, the two intersection points merge into one, and the radical axis becomes the common tangent at that point.

\begin{methode}[Common tangent of two touching circles]
Write the equations of the two circles in general form with the coefficients of $x^2$ and $y^2$ equal to $1$ in both (if necessary, divide the whole equation by the coefficient of $x^2$):
\[ S_1 = x^2 + y^2 + 2g_1x + 2f_1y + c_1 = 0, \qquad S_2 = x^2 + y^2 + 2g_2x + 2f_2y + c_2 = 0. \]
The common tangent at the point of contact is
\[ S_1 - S_2 = 0, \]
i.e. $\;2(g_1 - g_2)x + 2(f_1 - f_2)y + (c_1 - c_2) = 0$.
\end{methode}

\begin{exoresolu}[Two circles touching externally]
Show that the circles
\[ x^2 + y^2 - 4x - 6y + 9 = 0 \qquad \text{and} \qquad x^2 + y^2 + 2x + 2y - 7 = 0 \]
touch each other. Find the point of contact and the equation of the common tangent at that point.
\tcblower
\textbf{Centres and radii.} For $\mathcal{C}_1$: centre $C_1 = (2, 3)$, radius $r_1 = \sqrt{4 + 9 - 9} = 2$. For $\mathcal{C}_2$: centre $C_2 = (-1, -1)$, radius $r_2 = \sqrt{1 + 1 + 7} = 3$.

\textbf{Distance between centres.}
\[ d = |C_1C_2| = \sqrt{(2-(-1))^2 + (3-(-1))^2} = \sqrt{9 + 16} = \sqrt{25} = 5. \]
Since $d = 5 = 2 + 3 = r_1 + r_2$, the circles \textbf{touch externally}. $\blacksquare$

\textbf{Point of contact.} $T$ lies on $[C_1C_2]$ with $|C_1T| = r_1 = 2$ and $d = 5$:
\[ T = C_1 + \frac{r_1}{d}\,\overrightarrow{C_1C_2} = (2,3) + \frac{2}{5}(-3, -4) = \left(2 - \tfrac{6}{5},\; 3 - \tfrac{8}{5}\right) = \left(\tfrac{4}{5}, \tfrac{7}{5}\right). \]
Check: $|C_2T| = \sqrt{\left(\tfrac{4}{5}+1\right)^2 + \left(\tfrac{7}{5}+1\right)^2} = \sqrt{\tfrac{81}{25} + \tfrac{144}{25}} = \sqrt{9} = 3 = r_2$. \checkmark

\textbf{Common tangent.} Using $S_1 - S_2 = 0$:
\[ (-4x - 6y + 9) - (2x + 2y - 7) = 0 \;\Longrightarrow\; -6x - 8y + 16 = 0 \;\Longrightarrow\; 3x + 4y - 8 = 0. \]
Check that $T$ lies on it: $3\cdot\tfrac{4}{5} + 4\cdot\tfrac{7}{5} - 8 = \tfrac{12 + 28 - 40}{5} = 0$. \checkmark
\end{exoresolu}

\courssub{Determining Unknown Constants so that Two Circles Touch}

A typical GCE question gives two circles, one of which contains an unknown constant $k$, and asks for the value(s) of $k$ for which the circles touch. The strategy is always the same: express $d$, $r_1$, $r_2$ in terms of $k$ and solve \emph{both} equations $d = r_1 + r_2$ and $d = |r_1 - r_2|$.

\begin{exoresolu}[Finding $k$ so that two circles touch]
Find the values of $k$ for which the circle
\[ x^2 + y^2 - 6x - 8y + k = 0 \]
touches the circle $x^2 + y^2 = 4$.
\tcblower
$\mathcal{C}_2$: centre $C_2 = (0,0)$, radius $r_2 = 2$.
$\mathcal{C}_1$: centre $C_1 = (3, 4)$, radius $r_1 = \sqrt{9 + 16 - k} = \sqrt{25 - k}$ (we need $k < 25$ for a real circle).

Distance between centres: $d = \sqrt{9 + 16} = 5$.

\textbf{External contact:} $d = r_1 + r_2$ gives
\[ 5 = \sqrt{25 - k} + 2 \;\Longrightarrow\; \sqrt{25 - k} = 3 \;\Longrightarrow\; 25 - k = 9 \;\Longrightarrow\; k = 16. \]

\textbf{Internal contact:} $d = |r_1 - r_2|$ gives
\[ 5 = \left|\sqrt{25 - k} - 2\right|. \]
Either $\sqrt{25-k} - 2 = 5$, i.e. $\sqrt{25-k} = 7$, giving $25 - k = 49$, i.e. $k = -24$; or $\sqrt{25-k} - 2 = -5$, i.e. $\sqrt{25-k} = -3$, impossible.

Hence the circles touch when $\boxed{k = 16}$ (external contact) or $\boxed{k = -24}$ (internal contact, with $\mathcal{C}_2$ inside $\mathcal{C}_1$ since then $r_1 = 7 > r_2 = 2$).
\end{exoresolu}

\begin{attention}[``Touch'' may be internal or external]
Unless the question says \emph{externally} or \emph{internally}, the word \cle{touch} covers \textbf{both} cases. Always test $d = r_1 + r_2$ \emph{and} $d = |r_1 - r_2|$. In the example above, a candidate who only tries external contact finds $k = 16$ and loses the second answer $k = -24$ — typically half the marks.
\end{attention}

\begin{attention}[Internal contact is not ``one inside the other'']
Do not confuse $d = |r_1 - r_2|$ (the circles \emph{touch} internally, one common point) with $d < |r_1 - r_2|$ (one circle lies strictly inside the other, \emph{no} common point). Similarly $d > r_1 + r_2$ means separate, not touching. The equality sign is essential: touching is always an \emph{equality} between $d$ and one of $r_1 + r_2$, $|r_1 - r_2|$.
\end{attention}

\courssub{Number of Common Tangents (Useful for Multiple-Choice Questions)}

A quick way to visualise each configuration is to count the \cle{common tangents} — lines tangent to \emph{both} circles at once:

\begin{center}
\begin{tabularx}{\linewidth}{|l|X|c|}
\hline
\rowcolor{popblueL} \textbf{Configuration} & \textbf{Condition} & \textbf{Common tangents} \\
\hline
Separate & $d > r_1 + r_2$ & 4 \\
\hline
Touch externally & $d = r_1 + r_2$ & 3 \\
\hline
Intersecting & $|r_1 - r_2| < d < r_1 + r_2$ & 2 \\
\hline
Touch internally & $d = |r_1 - r_2|$ & 1 \\
\hline
One inside the other & $d < |r_1 - r_2|$ & 0 \\
\hline
\end{tabularx}
\end{center}

Notice the pattern: each time the circles pass from one configuration to the next, one pair of common tangents merges into a single tangent and then disappears. The tangent that survives the ``touching'' cases is precisely the common tangent at the point of contact, $S_1 - S_2 = 0$.

\begin{exercice}[Practice]
\begin{enumerate}
\item Determine the relative position of the circles $x^2 + y^2 = 9$ and $x^2 + y^2 - 8x + 6y + 9 = 0$. If they touch, find the point of contact and the common tangent.
\item Find the values of $a$ for which the circle $x^2 + y^2 - 2ax = 0$ touches the circle $x^2 + y^2 - 4y = 0$.
\item Two circles have centres $(1, 2)$ and $(4, 6)$ and touch externally. If the radius of the first is $2$, find the radius of the second and the equation of their common tangent.
\end{enumerate}
\end{exercice}

\begin{corrige}[Exercise 1]
$\mathcal{C}_1$: $C_1 = (0,0)$, $r_1 = 3$. $\mathcal{C}_2$: $C_2 = (4, -3)$, $r_2 = \sqrt{16 + 9 - 9} = 4$. Then $d = \sqrt{16 + 9} = 5$. Since $|r_1 - r_2| = 1 < 5 < 7 = r_1 + r_2$, the circles \textbf{intersect} in two points; they do not touch.
\end{corrige}

\begin{corrige}[Exercise 2]
$\mathcal{C}_1$: centre $(a, 0)$, radius $|a|$. $\mathcal{C}_2$: centre $(0, 2)$, radius $2$. Then $d = \sqrt{a^2 + 4}$.
External contact: $\sqrt{a^2+4} = |a| + 2 \Rightarrow a^2 + 4 = a^2 + 4|a| + 4 \Rightarrow |a| = 0 \Rightarrow a = 0$, but then $\mathcal{C}_1$ is the point circle $x^2+y^2=0$ (radius $0$) — rejected for a proper circle.
Internal contact: $\sqrt{a^2+4} = \big||a| - 2\big| \Rightarrow a^2 + 4 = a^2 - 4|a| + 4 \Rightarrow |a| = 0$, again rejected.
For $a \neq 0$, we have $(|a|-2)^2 < a^2+4 < (|a|+2)^2$, i.e. $|r_1-r_2| < d < r_1+r_2$, so the circles always intersect in two distinct points. Hence there is \textbf{no value of $a$} for which the two proper circles touch.
\end{corrige}

\begin{corrige}[Exercise 3]
$d = \sqrt{(4-1)^2 + (6-2)^2} = \sqrt{9 + 16} = 5$. External contact: $r_2 = d - r_1 = 5 - 2 = 3$.
Point of contact: $T = (1,2) + \tfrac{2}{5}(3, 4) = \left(\tfrac{11}{5}, \tfrac{18}{5}\right)$.
The common tangent is perpendicular to $C_1C_2$ (gradient $\tfrac{4}{3}$) at $T$, so its gradient is $-\tfrac{3}{4}$: $y - \tfrac{18}{5} = -\tfrac{3}{4}\left(x - \tfrac{11}{5}\right)$, i.e. $3x + 4y - 21 = 0$.
\end{corrige}

\begin{aretenir}[Key points of Section 8]
\begin{itemize}
\item Compare $d = |C_1C_2|$ with $r_1 + r_2$ and $|r_1 - r_2|$ to classify the position of two circles.
\item Touch externally $\iff d = r_1 + r_2$; touch internally $\iff d = |r_1 - r_2|$.
\item The point of contact lies on the line of centres and divides it in the ratio of the radii (internally for external contact, externally for internal contact).
\item The common tangent at the point of contact is $S_1 - S_2 = 0$ (with $x^2,y^2$ coefficients equal to $1$).
\item ``Touch'' without qualification means: test \emph{both} conditions.
\item Common tangents: $4, 3, 2, 1, 0$ as the circles move from separate to nested.
\end{itemize}
\end{aretenir}

\coursec{Section 9: Locus Problems and Synthesis for the GCE Exam}

In Section 8 we studied the conditions for two circles to touch, and throughout this chapter we have repeatedly answered the question: \emph{given an equation, what curve does it represent?} In this final section we reverse the question: \emph{given a geometric condition, what is the equation — and the nature — of the curve traced out?} This is the idea of a \cle{locus}, one of the most elegant and most frequently examined topics at the GCE Advanced Level. We close the chapter with synthesis problems in full GCE style, together with advice on exam technique.

\courssub{The Concept of a Locus}

Imagine a goat tethered by a rope of length $5$ m to a peg in a field near Bafoussam. As the goat moves while keeping the rope taut, the tip of its path traces a circle of radius $5$ m about the peg. The circle is the \emph{locus} of the goat: the set of all positions satisfying the condition ``distance from the peg equals $5$ m''.

\begin{definition}[Locus]
A \cle{locus} is the set of all points satisfying a given geometric condition. Finding a locus means determining the equation (and the geometric description) of the curve formed by all such points.
\end{definition}

The word \emph{locus} is Latin for ``place'': the locus is literally ``the place'' where the moving point is allowed to be. The power of coordinate geometry is that it converts a geometric condition into an algebraic equation.

\begin{methode}[The Four-Step Locus Strategy]
\begin{enumerate}
\item \textbf{Name the moving point.} Let $P(x, y)$ be a general point of the locus.
\item \textbf{Translate the condition.} Write the given geometric condition as an algebraic equation involving $x$ and $y$ (distance formula, gradient condition, etc.).
\item \textbf{Simplify.} Square where necessary, expand, and reduce to a recognisable form.
\item \textbf{Identify and describe.} State clearly what curve the equation represents (line, circle with centre and radius, etc.), and check for excluded points.
\end{enumerate}
\end{methode}

\begin{popfigure}
\begin{center}
\begin{tikzpicture}[scale=1, every node/.style={font=\small}]
\draw[rounded corners=4pt, draw=popblue, fill=popblueL, text width=3.1cm, align=center] (0,0) node[align=center]{\textbf{Step 1}\\ Let $P(x,y)$ be the moving point};
\draw[rounded corners=4pt, draw=poporange, fill=white, text width=3.1cm, align=center] (4.6,0) node[align=center]{\textbf{Step 2}\\ Write the condition algebraically};
\draw[rounded corners=4pt, draw=popgreen, fill=popgreenL, text width=3.1cm, align=center] (9.2,0) node[align=center]{\textbf{Step 3}\\ Simplify the equation};
\draw[rounded corners=4pt, draw=poppurple, fill=white, text width=3.1cm, align=center] (13.8,0) node[align=center]{\textbf{Step 4}\\ Identify and describe the curve};
\draw[-latex, thick, popdark] (1.7,0) -- (2.9,0);
\draw[-latex, thick, popdark] (6.3,0) -- (7.5,0);
\draw[-latex, thick, popdark] (10.9,0) -- (12.1,0);
\end{tikzpicture}
\end{center}
\legende{The four-step method for every locus problem.}
\end{popfigure}

\courssub{Elementary Loci: Fixed Distance from a Point or a Line}

The two simplest loci already contain an important contrast.

\textbf{Fixed distance from a fixed point.} The locus of $P$ with $|OP| = r$ for a fixed point $O$ is, by definition, a \cle{circle} of centre $O$ and radius $r$. With $O = (a,b)$ the equation is $(x-a)^2 + (y-b)^2 = r^2$.

\textbf{Fixed distance from a fixed line.} The locus of $P$ at a fixed distance $d$ from a fixed line $\ell$ is a \emph{pair of parallel lines}, one on each side of $\ell$, each at distance $d$ from $\ell$. For instance, the locus of points at distance $2$ from the line $y = 1$ is the pair $y = 3$ and $y = -1$.

\begin{attention}[A circle, not a single line]
Do not confuse the two situations: fixed distance from a \emph{point} gives one circle; fixed distance from a \emph{line} gives \emph{two} parallel lines (unless the problem restricts the point to one side). Forgetting one of the two lines is a classic loss of a mark.
\end{attention}

\courssub{Equidistant Points and the Circle of Apollonius}

\textbf{Equal distances: the perpendicular bisector.} Let $A(x_1,y_1)$ and $B(x_2,y_2)$ be fixed. The condition $PA = PB$ gives
\[ (x-x_1)^2+(y-y_1)^2 = (x-x_2)^2+(y-y_2)^2, \]
and after squaring out, the quadratic terms cancel, leaving a \emph{linear} equation: the locus is the \cle{perpendicular bisector} of $AB$.

\textbf{Constant ratio of distances.} Now replace equality by a ratio.

\begin{propriete}[Circle of Apollonius]
Let $A$ and $B$ be two fixed points and let $k > 0$, $k \neq 1$. The locus of points $P$ such that
\[ \frac{PA}{PB} = k \]
is a \cle{circle}, called the \emph{circle of Apollonius}. (For $k = 1$ the locus degenerates to the perpendicular bisector of $AB$.) The derivation below is instructive; the result itself is a GCE-useful extra.
\end{propriete}

\emph{Derivation.} Take $A = (-c, 0)$ and $B = (c, 0)$ (we may always choose axes this way). Then $PA^2 = k^2 PB^2$ gives
\[ (x+c)^2 + y^2 = k^2\big[(x-c)^2 + y^2\big]. \]
Expanding both sides,
\[ x^2 + 2cx + c^2 + y^2 = k^2 x^2 - 2k^2 c x + k^2 c^2 + k^2 y^2, \]
and collecting like terms:
\[ (1-k^2)x^2 + (1-k^2)y^2 + 2c(1+k^2)x + c^2(1-k^2) = 0. \]
Since $k \neq 1$ we may divide by $1-k^2$:
\[ x^2 + y^2 + \frac{2c(1+k^2)}{1-k^2}\,x + c^2 = 0, \]
which has the form $x^2+y^2+2gx+2fy+c'=0$ with $f=0$: a \emph{circle} whose centre lies on the line $AB$. $\blacksquare$

\begin{popfigure}
\begin{center}
\begin{tikzpicture}[scale=0.85]
\draw[popline] (-4.5,0) -- (6.5,0);
\fill[popdark] (-2,0) circle (2pt) node[below] {$A$};
\fill[popdark] (2,0) circle (2pt) node[below] {$B$};
\draw[thick, popblue] (3.333,0) circle (2.667);
\fill[poporange] (2,2.309) circle (2pt) node[above left] {$P_1$};
\fill[poporange] (3.333,2.667) circle (2pt) node[above right] {$P_2$};
\fill[poporange] (3.333,-2.667) circle (2pt) node[below] {$P_3$};
\draw[popgreen, dashed] (-2,0) -- (2,2.309);
\draw[poppurple, dashed] (2,0) -- (2,2.309);
\node[popblue] at (3.333,-3.3) {$PA = 2\,PB$};
\end{tikzpicture}
\end{center}
\legende{Circle of Apollonius for $A=(-2,0)$, $B=(2,0)$ with ratio $k=2$: centre $\left(\tfrac{10}{3},0\right)$, radius $\tfrac{8}{3}$; every point $P$ of the circle satisfies $PA = 2\,PB$.}
\end{popfigure}

\begin{exoresolu}[An Apollonius locus]
$A = (0,0)$ and $B = (6,0)$. Find the locus of $P$ such that $PA = 2\,PB$.
\tcblower
Let $P(x,y)$. The condition $PA^2 = 4\,PB^2$ gives
\[ x^2 + y^2 = 4\big[(x-6)^2 + y^2\big]. \]
Expanding: $x^2 + y^2 = 4x^2 - 48x + 144 + 4y^2$, so
\[ 3x^2 + 3y^2 - 48x + 144 = 0 \quad\Longrightarrow\quad x^2 + y^2 - 16x + 48 = 0. \]
Completing the square: $(x-8)^2 + y^2 = 16$.
\textbf{Description:} the locus is a circle of centre $(8,0)$ and radius $4$.
\end{exoresolu}

\courssub{Loci Defined by $PA^2 + PB^2$ or $PA^2 - PB^2$ Constant}

These two conditions appear regularly at the GCE and reward algebraic fluency.

\begin{exemplebox}[Sum of squares constant]
Let $A=(-a,0)$, $B=(a,0)$ and suppose $PA^2 + PB^2 = 2c$ (constant). Then
\[ (x+a)^2 + y^2 + (x-a)^2 + y^2 = 2c \;\Longrightarrow\; 2x^2 + 2y^2 + 2a^2 = 2c, \]
so $x^2 + y^2 = c - a^2$: a \textbf{circle centred at the midpoint of $AB$}, with radius $\sqrt{c-a^2}$ (real only if $c \ge a^2$).
\end{exemplebox}

\begin{exemplebox}[Difference of squares constant]
With the same points, suppose $PA^2 - PB^2 = k$ (constant). Then
\[ \big[(x+a)^2 + y^2\big] - \big[(x-a)^2 + y^2\big] = 4ax = k, \]
so $x = \dfrac{k}{4a}$: a \textbf{straight line perpendicular to $AB$}. Note the contrast: the \emph{sum} gives a circle, the \emph{difference} gives a line.
\end{exemplebox}

\courssub{Locus of the Midpoint of a Chord Through a Fixed Point}

A beautiful result links loci to the circle geometry of this chapter. Let $\mathcal{C}$ be a circle with centre $O$, and let $F$ be a fixed point inside it. Draw any chord of $\mathcal{C}$ passing through $F$, and let $M$ be its midpoint. Since the line from the centre to the midpoint of a chord is perpendicular to the chord, $\widehat{OMF} = 90^\circ$. By the converse of the angle-in-a-semicircle theorem, $M$ lies on the circle with diameter $OF$.

\begin{propriete}[Locus of midpoints of chords through a fixed point]
The locus of the midpoints of the chords of a circle $\mathcal{C}$ (centre $O$) which pass through a fixed interior point $F$ is the circle with diameter $OF$ (restricted to the interior of $\mathcal{C}$).
\end{propriete}

\begin{popfigure}
\begin{center}
\begin{tikzpicture}[scale=1.0]
\begin{axis}[width=10cm, height=7.5cm, axis equal image, axis lines=middle, xlabel={$x$}, ylabel={$y$},
xmin=-1.5, xmax=5.5, ymin=-3.5, ymax=3.5, xtick={0,1,2,3,4,5}, ytick={-3,-2,-1,1,2,3},
tick label style={font=\scriptsize}, label style={font=\small}]
\addplot[popblue, thick, domain=0:360, samples=100] ({2+3*cos(x)},{3*sin(x)});
\addplot[poppurple, thick, domain=0:360, samples=100] ({1.5+0.5*cos(x)},{0.5*sin(x)});
\addplot[only marks, mark=*, mark size=1.6pt, popdark] coordinates {(2,0) (1,0)};
\addplot[only marks, mark=*, mark size=1.4pt, poporange] coordinates {(1.5,0.5) (1.5,-0.5) (1.8,0.4)};
\addplot[popgreen, dashed, domain=-1.2:5.2] {x-1};
\addplot[popgreen, dashed, domain=-1.2:5.2] {1-x};
\addplot[popgreen, dashed, domain=-1.2:5.2] {0.5*(x-1)};
\node[popdark, font=\small] at (axis cs:2.35,-0.35) {$O$};
\node[popdark, font=\small] at (axis cs:0.75,-0.35) {$F$};
\node[poppurple, font=\small] at (axis cs:2.9,1.1) {locus of $M$};
\node[popblue, font=\small] at (axis cs:-0.9,2.9) {$\mathcal{C}$};
\node[poporange, font=\scriptsize] at (axis cs:1.05,0.75) {$M$};
\end{axis}
\end{tikzpicture}
\end{center}
\legende{Chords of $\mathcal{C}$ through the fixed point $F$ (dashed) have midpoints $M$ lying on the circle with diameter $OF$.}
\end{popfigure}

\begin{exoresolu}[Midpoint locus by algebra]
Find the locus of the midpoints of chords of the circle $x^2+y^2=25$ that pass through $F(3,0)$.
\tcblower
Let $M(h,k)$ be the midpoint of a chord. The chord with midpoint $(h,k)$ of the circle $x^2+y^2=25$ has equation $T = S_1$:
\[ hx + ky = h^2 + k^2. \]
Since the chord passes through $F(3,0)$: $3h = h^2 + k^2$. Replacing $(h,k)$ by $(x,y)$, the locus is
\[ x^2 + y^2 - 3x = 0 \quad\Longleftrightarrow\quad \left(x-\tfrac{3}{2}\right)^2 + y^2 = \tfrac{9}{4}, \]
the circle with diameter $OF$: centre $\left(\tfrac32,0\right)$, radius $\tfrac32$. (Only the part inside the given circle is attained.)
\end{exoresolu}

\courssub{Loci via Parametric Elimination}

Sometimes the moving point is given \emph{parametrically}: its coordinates depend on a parameter $t$ (or $\theta$). The locus is found by eliminating the parameter — exactly the technique of Lesson 10.

\begin{exoresolu}[Eliminating the parameter]
A point moves so that its position is $P = (3 + 2\cos\theta,\; -1 + 2\sin\theta)$. Find its locus.
\tcblower
We have $\cos\theta = \dfrac{x-3}{2}$ and $\sin\theta = \dfrac{y+1}{2}$. Using $\cos^2\theta + \sin^2\theta = 1$:
\[ (x-3)^2 + (y+1)^2 = 4. \]
The locus is the circle of centre $(3,-1)$ and radius $2$.
\end{exoresolu}

\begin{methode}[Parametric locus problems]
\begin{enumerate}
\item Express $\cos\theta$ and $\sin\theta$ (or the parameter $t$) in terms of $x$ and $y$.
\item Eliminate using an identity ($\cos^2+\sin^2=1$, or substitution for $t$).
\item Simplify and describe the resulting curve completely.
\end{enumerate}
\end{methode}

\courssub{Synthesis: GCE-Style Problems and Exam Technique}

GCE questions rarely test one idea in isolation. A typical Part B question might give a circle through three points, ask for a tangent, then impose a touching condition or a locus. The following exercise mixes several lessons of this chapter.

\begin{exercice}[GCE synthesis]
The circle $\mathcal{C}$ has equation $x^2 + y^2 - 4x - 6y + 9 = 0$.
\begin{enumerate}
\item Find the centre and radius of $\mathcal{C}$.
\item Show that the line $y = 1$ is a tangent to $\mathcal{C}$ and find the point of contact.
\item A point $P$ moves so that its distance from the centre of $\mathcal{C}$ equals its distance from the line $x = 0$. Show that the locus of $P$ is a parabola and find its equation.
\item Find the equation of the circle which passes through the points of intersection of $\mathcal{C}$ and the circle $x^2+y^2=9$ and also through the point $(0,0)$.
\end{enumerate}
\end{exercice}

\begin{corrige}[Exercise 1]
\begin{enumerate}
\item Complete the squares: $x^2-4x + y^2-6y + 9 = 0$ gives $(x-2)^2 - 4 + (y-3)^2 - 9 + 9 = 0$, i.e.\ $(x-2)^2 + (y-3)^2 = 4$. Centre $C(2,3)$, radius $r=2$.
\item The distance from $C(2,3)$ to the line $y = 1$ is $|3-1| = 2 = r$. Hence $y=1$ is a tangent to $\mathcal{C}$. The point of contact is the foot of the perpendicular from $C$ to the line, namely $(2,1)$.
\item Condition: $\sqrt{(x-2)^2+(y-3)^2} = |x|$. Squaring: $(x-2)^2 + (y-3)^2 = x^2$, i.e.\ $x^2 - 4x + 4 + (y-3)^2 = x^2$, so $(y-3)^2 = 4x - 4 = 4(x-1)$: a parabola with vertex $(1,3)$ and axis $y=3$.
\item Family through the intersections: $x^2+y^2-4x-6y+9 + \lambda(x^2+y^2-9) = 0$. Imposing $(0,0)$: $9 - 9\lambda = 0 \Rightarrow \lambda = 1$. Hence $2x^2+2y^2-4x-6y = 0$, i.e.\ $x^2+y^2-2x-3y=0$.
\end{enumerate}
\end{corrige}

\begin{attention}[State the full description — and check restrictions]
Two final warnings for the examination:
\begin{itemize}
\item Never leave the answer as a raw equation. Write: ``the locus is a \textbf{circle} of centre $(8,0)$ and radius $4$'' (or a line, a pair of lines, a parabola, with its key data).
\item Check for \textbf{restrictions and degenerate cases}: the ratio $k=1$ in Apollonius gives a \emph{line}, not a circle; $PA^2+PB^2 = 2c$ gives a real circle only if $c \ge a^2$; a midpoint locus may be limited to the interior of the given circle. Examiners award a specific mark for these observations.
\end{itemize}
\end{attention}

\begin{aretenir}[Key points of Section 9]
\begin{itemize}
\item A \cle{locus} is the set of all points satisfying a geometric condition; find it by the four-step method: name $P(x,y)$, translate, simplify, identify.
\item Fixed distance from a point $\Rightarrow$ circle; from a line $\Rightarrow$ two parallel lines; equidistant from two points $\Rightarrow$ perpendicular bisector.
\item $PA/PB = k \neq 1$ $\Rightarrow$ circle of Apollonius; $PA^2+PB^2$ constant $\Rightarrow$ circle; $PA^2-PB^2$ constant $\Rightarrow$ line.
\item Midpoints of chords of a circle through a fixed interior point $F$ lie on the circle of diameter $OF$.
\item Parametric loci: eliminate the parameter using $\cos^2\theta+\sin^2\theta = 1$ or substitution.
\item In the exam: describe the curve completely, check restrictions, and manage your time — a locus part is usually short once the method is launched.
\end{itemize}
\end{aretenir}

\begin{savaistu}[Apollonius of Perga]
The circle of Apollonius is named after the Greek geometer Apollonius of Perga (3rd century BC), whose treatise \emph{Conics} also gave us the names ellipse, parabola and hyperbola. Nearly two thousand years later, Descartes' coordinate geometry — the tool you have used throughout this chapter — turned his purely geometric results into simple algebraic manipulations. The same circle of Apollonius appears today in physics (field patterns between two charges) and in telecommunications (locating a receiver from signal-strength ratios).
\end{savaistu}

\textbf{Exam technique and time management.}\par
In the GCE Advanced Level paper, circle geometry questions are usually structured: early parts earn quick marks (centre, radius, tangent verification), while the final part is often a locus or a family-of-circles demand. Budget your time accordingly: secure the routine marks first, and do not spend more than a few minutes stuck on algebra — write down the condition $PA^2 = k^2 PB^2$ (or the family equation $S_1 + \lambda S_2 = 0$) immediately, since method marks are awarded for correct set-up even if the simplification is incomplete. Finally, re-read the question: if it asks for ``the locus'', a bare equation without a geometric description will cost you the last mark.

This completes our study of Circle Geometry. You now command every tool of the chapter: equations of circles in all their forms, tangents and their lengths, parametric representations, intersections of lines and circles, families of circles and the radical axis, touching conditions, and loci. Practise the mixed exercises of this section under timed conditions, and you will approach the GCE with confidence.

\newpage

\coursec{Integration activity}

\courssub{Aims of the activity}

By the end of this group activity, you should be able to:

\begin{itemize}
\item determine the equation of a circle passing through three given points, using both the general equation $x^{2}+y^{2}+2gx+2fy+c=0$ and the \cle{perpendicular bisector} method;
\item interpret the centre of that circle as the point \cle{equidistant} from three given locations, and its radius as a coverage distance;
\item use the \cle{perpendicular distance from a point to a line} to decide whether a line cuts, touches or misses a circle, and compute a \cle{chord length};
\item form the \cle{radical axis} of two circles by eliminating the quadratic terms, and give it a physical meaning;
\item identify the \cle{locus} of a point equidistant from two fixed points and compare it with the radical axis;
\item communicate a complete mathematical solution, supported by an accurate sketch on grid paper.
\end{itemize}

\begin{aretenir}[Skills mobilised]
Circle through three points $\bullet$ perpendicular bisectors $\bullet$ line--circle intersection $\bullet$ chord length $2\sqrt{r^{2}-d^{2}}$ $\bullet$ radical axis $S_{1}-S_{2}=0$ $\bullet$ locus (perpendicular bisector).
\end{aretenir}

\courssub{Situation}

The North-West Regional delegation of telecommunications wishes to improve mobile coverage between three villages situated on the highland plateau. On a simplified map drawn on squared paper, the planner chooses axes along two perpendicular dirt tracks and uses the kilometre as unit. The three villages are modelled by the points
\[ A(1,\,2), \qquad B(7,\,4), \qquad C(3,\,9) \qquad (\text{coordinates in km}). \]

A first relay antenna, called \textbf{Antenna 1}, must be placed so that its signal reaches the three villages equally well: its position must therefore be \emph{equidistant} from $A$, $B$ and $C$, and its coverage zone is the circle through the three villages. Recall the key geometric fact that makes this possible:

\begin{propriete}[Circle through three points]
Through any three \cle{non-collinear} points there passes exactly one circle. Its centre is the point equidistant from the three points, found as the intersection of the perpendicular bisectors of any two of the segments joining them.
\end{propriete}

A straight tarred road is modelled by the line of equation $y = x - 1$. For safety and engineering reasons, the delegation wants to know whether this road enters the coverage zone, and if so, over what length.

Later, a second relay, \textbf{Antenna 2}, is installed at the point $P(10,\,8)$ with a coverage radius of $4$ km. The engineers need to know along which line the signals of the two antennas \emph{balance} (a point where a phone could switch from one network to the other), and a technician walking so that he stays at equal \emph{distance} from the two antennas wants to describe his path.

Your group is the consultancy team. You must produce a full report: all equations, all decisions justified by calculation, and an accurate sketch on grid paper.

\begin{popfigure}
\begin{center}
\begin{tikzpicture}[scale=0.52]
\draw[popline, very thin] (-1,-2) grid (12,11);
\draw[->, thick] (-1,0) -- (12.4,0) node[below left] {$x$ (km)};
\draw[->, thick] (0,-2) -- (0,11.4) node[below left] {$y$ (km)};
\draw[popblue, thick] (3.29,5.13) circle (3.88);
\draw[poporange, thick] (10,8) circle (4);
\draw[popgreen, thick, domain=-1:10.5] plot (\x,{\x-1}) node[pos=0.88, above, sloped, font=\small] {road $y=x-1$};
\draw[poppurple, dashed, thick, domain=5.2:9.6] plot (\x,{(2392-255*\x)/109}) node[pos=0.12, above, sloped, font=\small] {radical axis};
\fill[popdark] (1,2) circle (2.2pt) node[below right] {$A$};
\fill[popdark] (7,4) circle (2.2pt) node[below right] {$B$};
\fill[popdark] (3,9) circle (2.2pt) node[above left] {$C$};
\fill[popblue] (3.29,5.13) circle (2.2pt) node[below left, font=\small] {$O_{1}$};
\fill[poporange] (10,8) circle (2.2pt) node[above right, font=\small] {$O_{2}$};
\end{tikzpicture}
\end{center}
\legende{Map of the situation: villages $A$, $B$, $C$, coverage circle of Antenna 1 (blue), Antenna 2 (orange), the road (green) and the radical axis (dashed).}
\end{popfigure}

\courssub{Tasks}

\textbf{Part A — Positioning Antenna 1.}
\begin{enumerate}
\item Write the general equation of a circle $x^{2}+y^{2}+2gx+2fy+c=0$ and explain why substituting $A$, $B$ and $C$ gives a system of three linear equations in $g$, $f$, $c$.
\item Solve the system and write the equation of the coverage circle.
\item Deduce the coordinates of the centre $O_{1}$ (the antenna position) and the coverage radius, to $2$ decimal places.
\item \emph{Check:} verify your centre by finding the perpendicular bisectors of $[AB]$ and $[AC]$ and intersecting them.
\end{enumerate}

\textbf{Part B — The road and the coverage zone.}
\begin{enumerate}
\setcounter{enumi}{4}
\item Compute the perpendicular distance $d$ from $O_{1}$ to the road $y=x-1$.
\item Comparing $d$ with the radius, decide whether the road cuts, touches or misses the coverage circle.
\item If the road cuts the circle, calculate the length of road lying inside the coverage zone.
\end{enumerate}

\textbf{Part C — Two antennas.}
\begin{enumerate}
\setcounter{enumi}{7}
\item Write the equation of the coverage circle of Antenna 2 (centre $(10,8)$, radius $4$).
\item Find the radical axis of the two circles and interpret it.
\item Find the locus of the technician who moves so that his distances to the two antennas are equal. Compare this locus with the radical axis: what do you notice, and why are they \emph{not} the same line?
\end{enumerate}

\textbf{Part D — Presentation.} Produce the final report with a TikZ-style sketch on grid paper showing the villages, both coverage circles, the road, the radical axis and the technician's path.

\courssub{Model answer}

\textbf{Part A.} The general equation $x^{2}+y^{2}+2gx+2fy+c=0$ contains three unknown constants $g$, $f$, $c$. Each point $(x_{0},y_{0})$ lying on the circle must satisfy the equation, and since $x_{0}^{2}+y_{0}^{2}$ is a known number, each substitution produces one \emph{linear} equation in $g$, $f$, $c$. Three non-collinear points therefore give three independent linear equations, which determine $g$, $f$, $c$ uniquely. Substituting:
\begin{align*}
A(1,2)&:\quad 1+4+2g+4f+c=0 &&\Rightarrow\quad 2g+4f+c=-5,\\
B(7,4)&:\quad 49+16+14g+8f+c=0 &&\Rightarrow\quad 14g+8f+c=-65,\\
C(3,9)&:\quad 9+81+6g+18f+c=0 &&\Rightarrow\quad 6g+18f+c=-90.
\end{align*}
Subtracting the first equation from the other two eliminates $c$:
\[ 12g+4f=-60 \;\Rightarrow\; 3g+f=-15, \qquad 4g+14f=-85. \]
From the first of these, $f=-15-3g$; substituting into the second:
\[ 4g+14(-15-3g)=-85 \;\Rightarrow\; -38g-210=-85 \;\Rightarrow\; -38g=125, \]
giving
\[ g=-\tfrac{125}{38}\approx -3.29, \qquad f=-15-3g=-\tfrac{195}{38}\approx -5.13, \]
\[ c=-5-2g-4f=-5+\tfrac{250}{38}+\tfrac{780}{38}=\tfrac{420}{19}\approx 22.11. \]
The coverage circle is
\[ x^{2}+y^{2}-\tfrac{125}{19}x-\tfrac{195}{19}y+\tfrac{420}{19}=0, \]
with centre $(-g,-f)$ and radius $\sqrt{g^{2}+f^{2}-c}$:
\[ O_{1}\left(\tfrac{125}{38},\,\tfrac{195}{38}\right)\approx(3.29,\,5.13), \qquad r=\sqrt{\tfrac{53650}{1444}-\tfrac{420}{19}}=\sqrt{\tfrac{10865}{722}}\approx 3.88\ \text{km}. \]

\emph{Check by bisectors.} Midpoint of $[AB]$: $(4,3)$; gradient of $AB$ is $\tfrac{4-2}{7-1}=\tfrac13$, so the perpendicular bisector has gradient $-3$ and equation $y-3=-3(x-4)$, i.e. $y=-3x+15$. Midpoint of $[AC]$: $\left(2,\tfrac{11}{2}\right)$; gradient of $AC$ is $\tfrac{9-2}{3-1}=\tfrac72$, so the bisector has gradient $-\tfrac27$ and equation $y-\tfrac{11}{2}=-\tfrac{2}{7}(x-2)$. Solving the two bisector equations simultaneously:
\[ -3x+15=\tfrac{11}{2}-\tfrac{2}{7}x+\tfrac{4}{7} \;\Rightarrow\; -\tfrac{19}{7}x=-\tfrac{125}{14} \;\Rightarrow\; x=\tfrac{125}{38}, \quad y=\tfrac{195}{38}. \]
Same centre as before. $\checkmark$

\textbf{Part B.} Write the road in the standard form $x-y-1=0$. The perpendicular distance from $O_{1}\left(\tfrac{125}{38},\tfrac{195}{38}\right)$ to the line is
\[ d=\frac{\left|\tfrac{125}{38}-\tfrac{195}{38}-1\right|}{\sqrt{1^{2}+(-1)^{2}}}=\frac{\left|-\tfrac{108}{38}\right|}{\sqrt2}=\frac{54}{19\sqrt2}\approx 2.01\ \text{km}. \]
Since $d\approx 2.01 < r\approx 3.88$, the road \textbf{cuts} the coverage circle in two points: a portion of the road lies inside the zone. That portion is a chord, and by Pythagoras' theorem applied to the right triangle formed by the centre, the foot of the perpendicular and an endpoint of the chord, half the chord has length $\sqrt{r^{2}-d^{2}}$. Hence
\[ L=2\sqrt{r^{2}-d^{2}}=2\sqrt{\tfrac{10865}{722}-\tfrac{2916}{722}}=2\sqrt{\tfrac{7949}{722}}\approx 6.64\ \text{km}. \]
About $6.64$ km of the road lies inside the coverage zone.

\textbf{Part C.} Antenna 2: $(x-10)^{2}+(y-8)^{2}=16$, which expands to
\[ S_{2}:\ x^{2}+y^{2}-20x-16y+148=0. \]
Writing $S_{1}: x^{2}+y^{2}-\tfrac{125}{19}x-\tfrac{195}{19}y+\tfrac{420}{19}=0$, the radical axis is $S_{1}-S_{2}=0$ (the quadratic terms cancel, leaving a straight line):
\[ \left(-\tfrac{125}{19}+20\right)x+\left(-\tfrac{195}{19}+16\right)y+\left(\tfrac{420}{19}-148\right)=0, \]
and since $-\tfrac{125}{19}+20=\tfrac{255}{19}$, $-\tfrac{195}{19}+16=\tfrac{109}{19}$ and $\tfrac{420}{19}-148=-\tfrac{2392}{19}$, multiplying through by $19$ gives
\[ 255x+109y-2392=0. \]
\emph{Interpretation:} along this line the powers of a point with respect to both circles are equal — equivalently, tangent segments drawn from a point of this line to the two circles have equal length. Physically, the signals of the two antennas balance along this line; a phone crossing it switches network.

For the technician, equidistance from $O_{1}$ and $O_{2}$ means his locus is the \textbf{perpendicular bisector} of $[O_{1}O_{2}]$. The midpoint is $\left(\tfrac{505}{76},\tfrac{499}{76}\right)$; the gradient of $O_{1}O_{2}$ is $\tfrac{8-\frac{195}{38}}{10-\frac{125}{38}}=\tfrac{109}{255}$, so the locus has gradient $-\tfrac{255}{109}$ and equation
\[ 255x+109y=255\cdot\tfrac{505}{76}+109\cdot\tfrac{499}{76}=\tfrac{91583}{38}\approx 2410.1. \]
\emph{Comparison:} the locus has the same coefficients of $x$ and $y$ as the radical axis, so the two lines are \textbf{parallel} — but distinct, since $2392\neq 2410.1$. The radical axis balances the \emph{powers} of a point with respect to the two circles (equal tangent lengths), whereas the technician's path balances the \emph{distances to the centres}; the two notions coincide only when the circles have equal radii.

\begin{attention}[Pitfalls to avoid in your report]
\begin{itemize}
\item Forgetting to halve the coefficients: the centre of $x^{2}+y^{2}+2gx+2fy+c=0$ is $(-g,-f)$, not $(-2g,-2f)$.
\item Using $d>r$ to conclude ``the line cuts the circle'' — it is the opposite: the line cuts the circle when $d<r$, touches it when $d=r$ and misses it when $d>r$.
\item Confusing the radical axis (equal powers) with the perpendicular bisector of the centres (equal distances).
\item Rounding too early: keep exact fractions throughout the working and round only the final answers.
\end{itemize}
\end{attention}

\begin{exercice}[Extension — for the fastest groups]
A third antenna of radius $3$ km is proposed at $(0,11)$. Show that the three radical axes (taken two circles at a time) are concurrent, and find their common point (the \cle{radical centre}). What could this point mean for the network engineer?
\end{exercice}

\begin{corrige}[Extension exercise]
The third circle is $S_{3}: x^{2}+(y-11)^{2}=9$, i.e. $S_{3}: x^{2}+y^{2}-22y+112=0$. The radical axis of $S_{1}$ and $S_{3}$ is $S_{1}-S_{3}=0$:
\[ -\tfrac{125}{19}x+\left(-\tfrac{195}{19}+22\right)y+\left(\tfrac{420}{19}-112\right)=0 \;\Rightarrow\; -125x+223y-1708=0, \]
after multiplying by $19$ (using $-\tfrac{195}{19}+22=\tfrac{223}{19}$ and $\tfrac{420}{19}-112=-\tfrac{1708}{19}$). The radical axis of $S_{1}$ and $S_{2}$ was $255x+109y-2392=0$. Solving these two linear equations simultaneously: from the second, $x=\dfrac{223y-1708}{125}$; substituting into the first,
\[ 255\cdot\frac{223y-1708}{125}+109y=2392 \;\Rightarrow\; 14098y=146908 \;\Rightarrow\; y=\tfrac{73454}{7049}\approx 10.42, \]
and then $x\approx 4.93$. The third radical axis, $S_{2}-S_{3}=0$, is $-20x+6y+36=0$, i.e. $10x-3y-18=0$; substituting $x\approx 4.93$, $y\approx 10.42$ gives $10(4.93)-3(10.42)-18\approx 0$, so the third axis passes through the same point: the three radical axes are \textbf{concurrent} at the radical centre $(4.93,\,10.42)$ (to $2$ d.p.). This concurrence is general: any point on two of the radical axes has equal powers with respect to all three circles, so it lies on the third axis too. For the engineer, the radical centre is the unique point where the three networks balance — a natural location for a shared switching station.
\end{corrige}

\newpage

\coursec{Methods and worked examples}

\coursec{Equations of a Circle — Centre at the Origin and General Centre}

\courssub{Introduction and Geometric Definition}
A \cle{circle} is the locus of a point $P(x,y)$ which moves in a plane such that its distance from a fixed point $C(a,b)$ (the \cle{centre}) is always equal to a constant $r$ (the \cle{radius}). Using the distance formula:
\[
PC = \sqrt{(x-a)^2 + (y-b)^2} = r \quad \implies \quad (x-a)^2 + (y-b)^2 = r^2.
\]
This is the \cle{standard form} (or centre-radius form) of the equation of a circle.

\begin{popfigure}
\begin{tikzpicture}[scale=0.8, >=stealth]
\draw[popline, thin] (-3,-2) grid (5,4);
\draw[->, thick] (-3,0) -- (5,0) node[right] {$x$};
\draw[->, thick] (0,-2) -- (0,4) node[above] {$y$};
\draw[popblue, thick] (1,1) circle (2.5);
\fill[popblue] (1,1) circle (2pt) node[below left] {$C(a,b)$};
\draw[poporange, dashed] (1,1) -- (3.5, 2.5) node[midway, above right] {$r$};
\fill[poporange] (3.5, 2.5) circle (2pt) node[above right] {$P(x,y)$};
\node at (2.3, 0.5) {$(x-a)^2+(y-b)^2=r^2$};
\end{tikzpicture}
\legende{The standard circle with centre $C(a,b)$ and radius $r$.}
\end{popfigure}

\begin{definition}[Circle Centred at the Origin]
If the centre is the origin $O(0,0)$, the equation reduces to:
\[
x^2 + y^2 = r^2.
\]
\end{definition}

\begin{propriete}[Centre and Radius from Standard Form]
For the circle $(x-a)^2 + (y-b)^2 = r^2$:
\begin{itemize}
\item Centre: $C(a,b)$.
\item Radius: $r = \sqrt{r^2}$ (must have $r^2 > 0$ for a real circle).
\end{itemize}
\end{propriete}

\begin{exemplebox}[Identifying Centre and Radius]
Find the centre and radius of: (i) $x^2+y^2=16$; (ii) $(x+2)^2+(y-3)^2=10$.
\tcblower
\textbf{Solution.}
(i) Centre $(0,0)$, radius $r=4$.
(ii) Rewrite as $(x-(-2))^2+(y-3)^2=(\sqrt{10})^2$. Centre $(-2,3)$, radius $\sqrt{10}$.
\end{exemplebox}

\courssub{Equation Given Centre and a Point on the Circle}
If the centre $C(a,b)$ and a point $P(x_1,y_1)$ on the circumference are known, the radius is $r = CP = \sqrt{(x_1-a)^2+(y_1-b)^2}$. Substitute into standard form.

\begin{exoresolu}[Centre and Point on Circumference]
Find the equation of the circle with centre $C(2,-3)$ passing through $P(5,1)$.
\tcblower
\textbf{Solution.}
$r^2 = (5-2)^2 + (1-(-3))^2 = 3^2 + 4^2 = 25$.
Equation: $(x-2)^2 + (y+3)^2 = 25$ or $x^2+y^2-4x+6y-12=0$.
\end{exoresolu}

\coursec{The General Equation and Circles on a Given Diameter}

\courssub{General Equation of a Circle}
Expanding the standard form $(x-a)^2+(y-b)^2=r^2$ gives:
\[
x^2+y^2 -2ax -2by + (a^2+b^2-r^2) = 0.
\]
Letting $-2a = 2g$, $-2b = 2f$, and $a^2+b^2-r^2 = c$, we obtain the \cle{general equation}:
\[
x^2 + y^2 + 2gx + 2fy + c = 0 \qquad (1)
\]
with centre $(-g, -f)$ and radius $r = \sqrt{g^2+f^2-c}$, provided $g^2+f^2-c > 0$.

\begin{propriete}[Conditions for a Real Circle]
The equation $x^2+y^2+2gx+2fy+c=0$ represents:
\begin{itemize}
\item A real circle if $g^2+f^2-c > 0$ (radius $r = \sqrt{g^2+f^2-c}$).
\item A point circle (radius zero) if $g^2+f^2-c = 0$.
\item No real locus (imaginary circle) if $g^2+f^2-c < 0$.
\end{itemize}
\end{propriete}

\begin{methode}[Completing the Square: General $\to$ Standard]
To find the centre and radius from $x^2+y^2+2gx+2fy+c=0$:
\begin{enumerate}
\item Group $x$ and $y$ terms: $(x^2+2gx) + (y^2+2fy) = -c$.
\item Complete squares: Add $g^2$ and $f^2$ to both sides:
$(x^2+2gx+g^2) + (y^2+2fy+f^2) = g^2+f^2-c$.
\item Write as: $(x+g)^2 + (y+f)^2 = g^2+f^2-c$.
\item Centre: $(-g, -f)$; Radius: $\sqrt{g^2+f^2-c}$.
\end{enumerate}
\end{methode}

\begin{exoresolu}[General to Standard Form]
Find the centre and radius of $x^2+y^2-6x+4y-12=0$.
\tcblower
\textbf{Solution.}
$(x^2-6x) + (y^2+4y) = 12$.
$(x-3)^2 - 9 + (y+2)^2 - 4 = 12$.
$(x-3)^2 + (y+2)^2 = 25$.
Centre $(3,-2)$, Radius $5$.
\end{exoresolu}

\courssub{Circle on a Given Diameter}
If $A(x_1,y_1)$ and $B(x_2,y_2)$ are endpoints of a diameter, the centre is the midpoint and the radius is half the length $AB$.

\begin{propriete}[Equation of Circle on Diameter $AB$]
Centre: $\left(\frac{x_1+x_2}{2}, \frac{y_1+y_2}{2}\right)$.
Radius: $\frac{1}{2}\sqrt{(x_2-x_1)^2+(y_2-y_1)^2}$.
Equation: $(x-x_1)(x-x_2) + (y-y_1)(y-y_2) = 0$ (diameter form).
\end{propriete}

\begin{exoresolu}[Circle on Diameter]
Find the equation of the circle with $A(2,-1)$ and $B(4,3)$ as diameter endpoints.
\tcblower
\textbf{Solution.}
Centre $C\left(\frac{2+4}{2}, \frac{-1+3}{2}\right) = (3,1)$.
$AB = \sqrt{(4-2)^2+(3+1)^2} = \sqrt{20} = 2\sqrt{5} \implies r = \sqrt{5}$.
Standard form: $(x-3)^2+(y-1)^2 = 5$.
General form: $x^2+y^2-6x-2y+5=0$.
\emph{Check using diameter form:} $(x-2)(x-4)+(y+1)(y-3)=0 \implies x^2-6x+8+y^2-2y-3=0 \implies x^2+y^2-6x-2y+5=0$. \checkmark
\end{exoresolu}

\begin{attention}[Common Mistake: Sign of Centre]
In $x^2+y^2+2gx+2fy+c=0$, the centre is $(-g, -f)$, \textbf{not} $(g,f)$.
Example: $x^2+y^2-4x+6y=0 \implies g=-2, f=3 \implies$ Centre $(2, -3)$.
\end{attention}

\coursec{Circle Through Three Points and Circles Under Geometric Conditions}

\courssub{Circle Through Three Non-Collinear Points}
Three non-collinear points uniquely determine a circle. We use the general form $x^2+y^2+2gx+2fy+c=0$ and substitute the coordinates to get three linear equations in $g, f, c$.

\begin{methode}[Circle Through Three Points $P_1, P_2, P_3$]
\begin{enumerate}
\item Let the circle be $x^2+y^2+2gx+2fy+c=0$.
\item Substitute each point $(x_i,y_i)$:
$2gx_i + 2fy_i + c = -(x_i^2+y_i^2)$ for $i=1,2,3$.
\item Solve the $3\times 3$ system for $g, f, c$ (elimination or determinants).
\item Write the equation. Centre $=(-g,-f)$, Radius $=\sqrt{g^2+f^2-c}$.
\item \textbf{Alternative (Perpendicular Bisectors):} Find equations of perpendicular bisectors of chords $P_1P_2$ and $P_2P_3$. Their intersection is the centre. Radius = distance to any point. Use as a check.
\end{enumerate}
\end{methode}

\begin{exoresolu}[Circle Through Three Points]
A circle passes through $P(1,2)$, $Q(3,-4)$, $R(5,2)$. Find its equation, centre and radius.
\tcblower
\textbf{Solution.}
Let circle be $x^2+y^2+2gx+2fy+c=0$.
$P(1,2): 1+4+2g+4f+c=0 \implies 2g+4f+c=-5$ \quad (1)
$Q(3,-4): 9+16+6g-8f+c=0 \implies 6g-8f+c=-25$ \quad (2)
$R(5,2): 25+4+10g+4f+c=0 \implies 10g+4f+c=-29$ \quad (3)

(3) - (1): $8g = -24 \implies g = -3$.
Sub $g=-3$ in (1): $-6+4f+c=-5 \implies 4f+c=1$ \quad (4)
Sub $g=-3$ in (2): $-18-8f+c=-25 \implies -8f+c=-7$ \quad (5)

(4) - (5): $12f = 8 \implies f = \frac{2}{3}$.
From (4): $4(\frac{2}{3})+c=1 \implies c = -\frac{5}{3}$.

Equation: $x^2+y^2-6x+\frac{4}{3}y-\frac{5}{3}=0$ or $\mathbf{3x^2+3y^2-18x+4y-5=0}$.
Centre: $(3, -\frac{2}{3})$.
Radius: $\sqrt{9+\frac{4}{9}+\frac{5}{3}} = \sqrt{\frac{100}{9}} = \frac{10}{3}$.
\end{exoresolu}

\courssub{Circles Touching the Coordinate Axes}
A circle touches the $x$-axis if $|y_{\text{centre}}| = r$. It touches the $y$-axis if $|x_{\text{centre}}| = r$.
If it touches both axes, centre is $(\pm r, \pm r)$ (four quadrants). Equation: $(x \mp r)^2 + (y \mp r)^2 = r^2$.

\begin{exoresolu}[Circle Touching Both Axes]
Find the equation of the circle of radius $3$ touching both axes in the second quadrant.
\tcblower
\textbf{Solution.}
Second quadrant: centre $(-3, 3)$.
Equation: $(x+3)^2 + (y-3)^2 = 9 \implies x^2+y^2+6x-6y+9=0$.
\end{exoresolu}

\courssub{Circle Through Two Points with Centre on a Given Line}
The centre lies on the perpendicular bisector of the chord joining the two points \emph{and} on the given line. Intersection gives the centre.

\begin{exoresolu}[Two Points, Centre on Line]
Find the circle through $A(1,2)$, $B(3,4)$ with centre on $y=2x-1$.
\tcblower
\textbf{Solution.}
Midpoint of $AB$: $M(2,3)$. Slope $AB=1 \implies$ perp. bisector slope $=-1$.
Eq. perp. bisector: $y-3 = -(x-2) \implies y = -x+5$.
Intersect with $y=2x-1$: $-x+5 = 2x-1 \implies 3x=6 \implies x=2, y=3$.
Centre $C(2,3)$. Radius $r = CA = \sqrt{(2-1)^2+(3-2)^2} = \sqrt{2}$.
Equation: $(x-2)^2+(y-3)^2=2$ or $x^2+y^2-4x-6y+11=0$.
\end{exoresolu}

\courssub{Circle Through Two Points and Touching a Given Line}
Let the line be $L: ax+by+c=0$. Centre $(h,k)$ satisfies:
\begin{enumerate}
\item Equidistant from the two points (lies on perp. bisector).
\item Distance to line $= r = \text{distance to either point}$.
\end{enumerate}
$\frac{|ah+bk+c|}{\sqrt{a^2+b^2}} = \sqrt{(h-x_1)^2+(k-y_1)^2}$.
Solve for $(h,k)$ (usually two solutions).

\begin{exoresolu}[Two Points, Touching a Line]
Find circles through $A(1,1)$, $B(2,2)$ touching the line $x+y=0$.
\tcblower
\textbf{Solution.}
Midpoint $M(1.5, 1.5)$. Slope $AB=1 \implies$ perp. bisector: $y-1.5 = -1(x-1.5) \implies x+y=3$.
Centre $(h,k)$ lies on $h+k=3 \implies k=3-h$.
Distance to line $x+y=0$: $r = \frac{|h+k|}{\sqrt{2}} = \frac{3}{\sqrt{2}}$.
Distance to $A(1,1)$: $r^2 = (h-1)^2+(k-1)^2 = (h-1)^2+(2-h)^2 = 2h^2-6h+5$.
Equate $r^2$: $\frac{9}{2} = 2h^2-6h+5 \implies 4h^2-12h+1=0$.
$h = \frac{12\pm\sqrt{144-16}}{8} = \frac{3\pm\sqrt{8}}{2} = \frac{3\pm2\sqrt{2}}{2}$.
Two circles exist. Centres $\left(\frac{3+2\sqrt{2}}{2}, \frac{3-2\sqrt{2}}{2}\right)$ and $\left(\frac{3-2\sqrt{2}}{2}, \frac{3+2\sqrt{2}}{2}\right)$, radius $\frac{3}{\sqrt{2}}$.
\end{exoresolu}

\courssub{Circle Touching a Given Line (at a Given Point or Not)}
If the point of contact $T(x_1,y_1)$ on line $L$ is given: Centre lies on normal to $L$ through $T$. Radius $= CT$.
If only the line is given (and other conditions like radius or centre on line), use distance formula $r = \frac{|ah+bk+c|}{\sqrt{a^2+b^2}}$.

\coursec{Tangents to a Circle}

\courssub{Equation of Tangent at a Point on the Circle}
For the circle $S: x^2+y^2+2gx+2fy+c=0$ and point $P(x_1,y_1)$ \emph{on} the circle ($S_{11}=0$), the tangent is:
\[
T = 0 \quad \text{where} \quad T: xx_1+yy_1+g(x+x_1)+f(y+y_1)+c=0.
\]
For circle centred at origin $x^2+y^2=r^2$: $xx_1+yy_1=r^2$.

\begin{popfigure}
\begin{tikzpicture}[scale=0.8, >=stealth]
\draw[popline, thin] (-2,-2) grid (5,4);
\draw[->, thick] (-2,0) -- (5,0) node[right] {$x$};
\draw[->, thick] (0,-2) -- (0,4) node[above] {$y$};
\draw[popblue, thick] (1,1) circle (2);
\fill[popblue] (1,1) circle (2pt) node[below left] {$C$};
\fill[poporange] (3,1) circle (2pt) node[below right] {$P(x_1,y_1)$};
\draw[poporange, thick] (-1,1) -- (5,1) node[right] {Tangent $T=0$};
\draw[popgreen, dashed] (1,1) -- (3,1) node[midway, above] {$r$};
\end{tikzpicture}
\legende{Tangent at $P$ is perpendicular to radius $CP$.}
\end{popfigure}

\begin{propriete}[Tangent-Radius Perpendicularity]
The tangent at any point $P$ on a circle is perpendicular to the radius $CP$.
Gradient of radius $m_{CP} = \frac{y_1-b}{x_1-a} \implies$ Gradient of tangent $m_T = -\frac{1}{m_{CP}} = -\frac{x_1-a}{y_1-b}$ (if $y_1 \neq b$).
\end{propriete}

\courssub{Condition for a Line to be Tangent}
Line $y = mx + c$ (or $y=mx+d$ to avoid confusion with circle constant $c$) and circle $x^2+y^2=r^2$ (centre origin).
Substitute: $x^2+(mx+d)^2=r^2 \implies (1+m^2)x^2 + 2md x + (d^2-r^2)=0$.
Tangency $\iff$ Discriminant $\Delta = 0$:
\[
(2md)^2 - 4(1+m^2)(d^2-r^2) = 0 \implies d^2 = r^2(1+m^2).
\]
For circle $(x-a)^2+(y-b)^2=r^2$, substitute line equation, set $\Delta=0$.

\begin{methode}[Tangent Condition and Equation]
\begin{enumerate}
\item \textbf{Given line $y=mx+d$ and circle centre $(a,b)$ radius $r$:} Substitute $y$, get quadratic in $x$. Set $\Delta=0$ to find $d$ (or $m$).
\item \textbf{Given point $P(x_1,y_1)$ on circle:} Use $T=0$ formula or gradient method.
\item \textbf{Given external point $P(x_1,y_1)$:} Find tangents by solving for $m$ in $y-y_1=m(x-x_1)$ with $\Delta=0$, or use chord of contact $T=0$ (where $T$ uses external point).
\end{enumerate}
\end{methode}

\begin{exoresolu}[Tangent Condition: Line $y=2x+k$ tangent to $x^2+y^2-4x+2y-4=0$]
Find $k$ and points of contact.
\tcblower
\textbf{Solution.}
Circle: $(x-2)^2+(y+1)^2=9$. Centre $C(2,-1)$, $r=3$.
Substitute $y=2x+k$:
$(x-2)^2+(2x+k+1)^2=9$.
$5x^2 + 4kx + (k^2+2k-4) = 0$.
$\Delta = (4k)^2 - 4(5)(k^2+2k-4) = 16k^2 - 20k^2 - 40k + 80 = -4k^2 - 40k + 80 = 0$.
$k^2 + 10k - 20 = 0 \implies k = -5 \pm 3\sqrt{5}$.
Points of contact (double root $x = -\frac{4k}{10} = -\frac{2k}{5}$):
For $k_1 = -5+3\sqrt{5}$: $x_1 = 2 - \frac{6\sqrt{5}}{5}$, $y_1 = -1 + \frac{3\sqrt{5}}{5}$.
For $k_2 = -5-3\sqrt{5}$: $x_2 = 2 + \frac{6\sqrt{5}}{5}$, $y_2 = -1 - \frac{3\sqrt{5}}{5}$.
\end{exoresolu}

\courssub{Length of Tangent from an External Point}
For point $P(x_1,y_1)$ outside circle $S: x^2+y^2+2gx+2fy+c=0$, the length of tangent $PT$ is:
\[
PT = \sqrt{x_1^2+y_1^2+2gx_1+2fy_1+c} = \sqrt{S_{11}}.
\]
This is the square root of the \cle{Power of the Point} $S_{11}$.
If $S_{11} > 0$: Point outside, two real tangents.
If $S_{11} = 0$: Point on circle, tangent length $0$.
If $S_{11} < 0$: Point inside, no real tangents.

\begin{exoresolu}[Length of Tangent]
Find the length of the tangent from $P(3,4)$ to the circle $x^2+y^2-2x-4y-4=0$.
\tcblower
\textbf{Solution.}
$S_{11} = 3^2+4^2-2(3)-4(4)-4 = 9+16-6-16-4 = -1 < 0$.
Point is \emph{inside} the circle. No real tangent exists.
\end{exoresolu}

\coursec{Parametric Equations of a Circle}

\courssub{Standard Parametrisation}
A circle with centre $(a,b)$ and radius $r$ can be described by the parameter $\theta$ (angle with positive $x$-axis):
\[
x = a + r\cos\theta, \quad y = b + r\sin\theta, \qquad \theta \in [0, 2\pi).
\]

\begin{popfigure}
\begin{tikzpicture}[scale=0.8, >=stealth]
\draw[popline, thin] (-1,-1) grid (5,4);
\draw[->, thick] (-1,0) -- (5,0) node[right] {$x$};
\draw[->, thick] (0,-1) -- (0,4) node[above] {$y$};
\draw[popblue, thick] (2,1.5) circle (2);
\fill[popblue] (2,1.5) circle (2pt) node[below left] {$C(a,b)$};
\draw[poporange, dashed] (2,1.5) -- (4, 3.2) node[midway, above right] {$r$};
\fill[poporange] (4, 3.2) circle (2pt) node[above right] {$P(x,y)$};
\draw[popgreen, ->] (2,1.5) -- (4,1.5) node[midway, below] {$r\cos\theta$};
\draw[popgreen, ->] (4,1.5) -- (4,3.2) node[midway, right] {$r\sin\theta$};
\draw[poppurple, ->] (3,1.5) arc (0:40:1) node[midway, right] {$\theta$};
\node at (3, 3) {$x=a+r\cos\theta$};
\node at (3, 2.7) {$y=b+r\sin\theta$};
\end{tikzpicture}
\legende{Parametric coordinates: $\theta$ is the angle $CP$ makes with the horizontal.}
\end{popfigure}

\courssub{Conversion: Parametric $\to$ Cartesian}
From $x = a + r\cos\theta$, $y = b + r\sin\theta$:
$\cos\theta = \frac{x-a}{r}$, $\sin\theta = \frac{y-b}{r}$.
Using $\cos^2\theta+\sin^2\theta=1$:
\[
\left(\frac{x-a}{r}\right)^2 + \left(\frac{y-b}{r}\right)^2 = 1 \implies (x-a)^2+(y-b)^2=r^2.
\]

\begin{exoresolu}[Parametric to Cartesian and Tangent]
Circle: $x = 1 + 3\cos\theta$, $y = -2 + 3\sin\theta$.
(i) Cartesian equation. (ii) Tangent at $\theta = \frac{\pi}{3}$.
\tcblower
\textbf{Solution.}
(i) $\cos\theta = \frac{x-1}{3}$, $\sin\theta = \frac{y+2}{3}$.
$\left(\frac{x-1}{3}\right)^2 + \left(\frac{y+2}{3}\right)^2 = 1 \implies \mathbf{(x-1)^2+(y+2)^2=9}$.
Centre $(1,-2)$, $r=3$.

(ii) At $\theta=\frac{\pi}{3}$: $\cos\frac{\pi}{3}=\frac{1}{2}$, $\sin\frac{\pi}{3}=\frac{\sqrt{3}}{2}$.
Point $P$: $x_1 = 1+3\cdot\frac{1}{2} = \frac{5}{2}$, $y_1 = -2+3\cdot\frac{\sqrt{3}}{2} = \frac{3\sqrt{3}-4}{2}$.
\textbf{Method 1 (Gradient):} $m_{CP} = \frac{y_1+2}{x_1-1} = \frac{3\sqrt{3}/2}{3/2} = \sqrt{3}$.
$m_T = -\frac{1}{\sqrt{3}} = -\frac{\sqrt{3}}{3}$.
Tangent: $y - \frac{3\sqrt{3}-4}{2} = -\frac{\sqrt{3}}{3}\left(x - \frac{5}{2}\right)$.
Multiply by 6: $6y - 9\sqrt{3}+12 = -2\sqrt{3}(2x-5) = -4\sqrt{3}x+10\sqrt{3}$.
$\mathbf{4\sqrt{3}x + 6y + 12 - 19\sqrt{3} = 0}$ (or simplified $\mathbf{x+\sqrt{3}y = 7-2\sqrt{3}}$).

\textbf{Method 2 (Formula $T=0$):} Circle: $x^2+y^2-2x+4y-4=0$ ($g=-1, f=2, c=-4$).
$xx_1+yy_1 -1(x+x_1) +2(y+y_1) -4 = 0$.
$x\frac{5}{2}+y\frac{3\sqrt{3}-4}{2} -x-\frac{5}{2} +2y+3\sqrt{3}-4 -4 = 0$.
Multiply by 2: $5x + (3\sqrt{3}-4)y -2x-5 +4y+6\sqrt{3}-16 = 0$.
$3x + 3\sqrt{3}y + 6\sqrt{3} -21 = 0 \implies \mathbf{x+\sqrt{3}y = 7-2\sqrt{3}}$. \checkmark
\end{exoresolu}

\coursec{Intersection of a Line and a Circle}

\courssub{Algebraic Method and Discriminant Analysis}
Substitute line equation into circle equation $\to$ quadratic in $x$ (or $y$).
Let quadratic be $Ax^2+Bx+C=0$, discriminant $\Delta = B^2-4AC$.

\begin{propriete}[Line-Circle Intersection]
\begin{itemize}
\item $\Delta > 0$: Line cuts circle in \textbf{two distinct points} (chord).
\item $\Delta = 0$: Line \textbf{touches} circle (tangent), one repeated root (point of contact).
\item $\Delta < 0$: Line \textbf{does not meet} the circle.
\end{itemize}
\end{propriete}

\begin{methode}[Chord Length and Midpoint]
If line $y=mx+d$ cuts circle giving roots $x_1, x_2$:
\begin{enumerate}
\item Chord length $AB = \sqrt{1+m^2}\,|x_1-x_2| = \frac{\sqrt{1+m^2}\sqrt{\Delta}}{|A|}$.
\item Midpoint $M$: $x_M = \frac{x_1+x_2}{2} = -\frac{B}{2A}$, $y_M = mx_M+d$.
\item \textbf{Geometric property:} The perpendicular from centre to chord bisects the chord.
Distance from centre to line $d = \frac{|Ax_c+By_c+C|}{\sqrt{A^2+B^2}}$.
Half-chord length $= \sqrt{r^2-d^2}$.
\end{enumerate}
\end{methode}

\begin{exoresolu}[Chord Length and Midpoint]
Line $y=x-2$ cuts circle $x^2+y^2-4x+6y-12=0$ at $A,B$. Find $A,B$, length $AB$, midpoint $M$.
\tcblower
\textbf{Solution.}
Circle: $(x-2)^2+(y+3)^2=25$. Centre $C(2,-3)$, $r=5$.
Substitute $y=x-2$:
$x^2+(x-2)^2-4x+6(x-2)-12=0 \implies 2x^2-2x-20=0 \implies x^2-x-10=0$.
$\Delta = 1+40=41 > 0$ (chord).
$x = \frac{1\pm\sqrt{41}}{2}$, $y = x-2 = \frac{-3\pm\sqrt{41}}{2}$.
$A\left(\frac{1+\sqrt{41}}{2}, \frac{-3+\sqrt{41}}{2}\right)$, $B\left(\frac{1-\sqrt{41}}{2}, \frac{-3-\sqrt{41}}{2}\right)$.
Length $AB = \sqrt{(\sqrt{41})^2+(\sqrt{41})^2} = \sqrt{82}$.
Midpoint $M$: $x_M = \frac{1}{2}$, $y_M = -\frac{3}{2}$.
Check: $CM$ gradient $= \frac{-1.5+3}{0.5-2} = -1 \perp$ line gradient $1$. \checkmark
\end{exoresolu}

\begin{savaistu}[Application: Satellite Coverage]
A satellite at height $h$ above Earth (radius $R$) covers a circular cap. The boundary circle on Earth's surface is the intersection of the Earth sphere with a cone from the satellite. In 2D cross-section, it reduces to line-circle intersection problems. Cameroon's SAT-1 satellite uses such geometry for footprint calculation.
\end{savaistu}

\coursec{Intersection of Two Circles, Families of Circles and the Radical Axis}

\courssub{Intersection of Two Circles}
Two circles $S_1=0$, $S_2=0$. Subtract to get the \cle{Radical Axis} (a straight line):
\[
S_1 - S_2 = 0 \implies 2(g_1-g_2)x + 2(f_1-f_2)y + (c_1-c_2) = 0.
\]
This line passes through the intersection points (if they intersect). Solve the system: Radical Axis + one Circle.

\begin{popfigure}
\begin{tikzpicture}[scale=0.7, >=stealth]
\draw[popline, thin] (-3,-3) grid (7,5);
\draw[->, thick] (-3,0) -- (7,0) node[right] {$x$};
\draw[->, thick] (0,-3) -- (0,5) node[above] {$y$};
\draw[popblue, thick] (0,0) circle (3);
\draw[poporange, thick] (4,1) circle (2.5);
\fill[popblue] (0,0) circle (2pt) node[below left] {$C_1$};
\fill[poporange] (4,1) circle (2pt) node[below right] {$C_2$};
\draw[popgreen, thick, dashed] (-2.5, -1.5) -- (5.5, 3.5) node[above right] {Radical Axis};
\fill[poppurple] (1.5, 2.2) circle (2pt) node[above] {$P$};
\fill[poppurple] (2.8, -1.2) circle (2pt) node[below] {$Q$};
\end{tikzpicture}
\legende{Two intersecting circles and their radical axis (line through $P$ and $Q$).}
\end{popfigure}

\courssub{Family of Circles Through Intersection of Two Circles}
The family of circles passing through the intersection of $S_1=0$ and $S_2=0$ is:
\[
S_1 + \lambda S_2 = 0 \quad (\lambda \neq -1) \qquad \text{or} \qquad S_1 + \lambda(S_1-S_2) = 0.
\]
For $\lambda = -1$ (first form) or $\lambda \to \infty$, we get the radical axis (a line).
To find a particular circle through the intersection and a given point $P$, substitute $P$ to find $\lambda$.

\begin{exoresolu}[Radical Axis and Family of Circles]
$S_1: x^2+y^2-4x+2y-7=0$, $S_2: x^2+y^2+2x-6y+5=0$.
(i) Radical axis. (ii) Circle through intersection and $(1,2)$.
\tcblower
\textbf{Solution.}
(i) $S_1-S_2 = -6x+8y-12=0 \implies \mathbf{3x-4y+6=0}$.
(ii) Family: $S_1 + \lambda S_2 = 0$.
At $(1,2)$: $S_1(1,2) = 1+4-4+4-7 = -2$. $S_2(1,2) = 1+4+2-12+5 = 0$.
Equation: $-2 + \lambda(0) = 0 \implies -2=0$ (impossible for finite $\lambda$).
This means $(1,2)$ lies on $S_2$ but not $S_1$. The required circle is $S_2$ itself (limit $\lambda\to\infty$).
Equation: $\mathbf{x^2+y^2+2x-6y+5=0}$.
\end{exoresolu}

\courssub{Orthogonal Circles}
Two circles $S_1, S_2$ cut orthogonally if tangents at intersection are perpendicular.
Condition: $2g_1g_2 + 2f_1f_2 = c_1 + c_2$.
Equivalently: $d^2 = r_1^2 + r_2^2$ (Pythagoras at intersection point).

\coursec{Conditions for Two Circles to Touch}

\courssub{Classification by Centre Distance}
Let circles have centres $C_1, C_2$, radii $r_1, r_2$ ($r_1 \ge r_2$), distance $d = C_1C_2$.

\begin{center}
\begin{tabular}{|c|c|c|l|}
\hline
\textbf{Position} & \textbf{Condition} & \textbf{Common Tangents} & \textbf{Point of Contact} \\
\hline
Separate (outside) & $d > r_1+r_2$ & 4 & None \\
Touch externally & $d = r_1+r_2$ & 3 & Divides $C_1C_2$ internally in ratio $r_1:r_2$ \\
Intersect (2 pts) & $|r_1-r_2| < d < r_1+r_2$ & 2 & N/A (chord = radical axis) \\
Touch internally & $d = r_1-r_2$ & 1 & Divides $C_1C_2$ externally in ratio $r_1:r_2$ \\
One inside other & $d < r_1-r_2$ & 0 & None \\
Concentric & $d=0$ & 0 (or $\infty$ if $r_1=r_2$) & N/A \\
\hline
\end{tabular}
\end{center}

\begin{methode}[Point of Contact Coordinates]
For centres $C_1(x_1,y_1)$, $C_2(x_2,y_2)$:
\begin{enumerate}
\item \textbf{External touch:} $T$ divides $C_1C_2$ internally $r_1:r_2$.
$T = \left( \frac{r_2 x_1 + r_1 x_2}{r_1+r_2}, \frac{r_2 y_1 + r_1 y_2}{r_1+r_2} \right)$.
\item \textbf{Internal touch ($r_1>r_2$):} $T$ divides $C_1C_2$ externally $r_1:r_2$ (on extension past $C_2$).
$T = \left( \frac{r_2 x_1 - r_1 x_2}{r_2-r_1}, \frac{r_2 y_1 - r_1 y_2}{r_2-r_1} \right)$.
\end{enumerate}
\end{methode}

\begin{exoresolu}[Touching Circles]
$C_1: x^2+y^2-6x+4y-12=0$ (Centre $O_1(3,-2)$, $r_1=5$).
$C_2$ centre $O_2(10,-5)$ touches $C_1$ externally. Find $r_2$, eq. of $C_2$.
If $C_2$ touches $C_1$ internally with $r_2=10$, find its centre.
\tcblower
\textbf{Solution.}
$d = O_1O_2 = \sqrt{(10-3)^2+(-5+2)^2} = \sqrt{58}$.
\textbf{External:} $d = r_1+r_2 \implies \sqrt{58} = 5+r_2 \implies r_2 = \sqrt{58}-5$.
Eq: $(x-10)^2+(y+5)^2 = (\sqrt{58}-5)^2 = 83-10\sqrt{58}$.

\textbf{Internal ($r_2=10$):} $d = r_2-r_1 = 5$.
$O_2$ lies on circle centre $O_1$ radius $5$: $(h-3)^2+(k+2)^2=25$.
If $O_2$ lies on line $O_1O_2$ (direction $(7,-3)$):
$O_2 = (3,-2) \pm \frac{5}{\sqrt{58}}(7,-3)$.
Two possible centres: $\left(3+\frac{35}{\sqrt{58}}, -2-\frac{15}{\sqrt{58}}\right)$ and $\left(3-\frac{35}{\sqrt{58}}, -2+\frac{15}{\sqrt{58}}\right)$.
\end{exoresolu}

\coursec{Circle Through Intersection of a Line and a Circle}

\courssub{Family of Circles Through Line-Circle Intersection}
Let circle $S: x^2+y^2+2gx+2fy+c=0$ and line $L: lx+my+n=0$.
The family of circles passing through their intersection points is:
\[
S + \lambda L = 0 \implies x^2+y^2+2gx+2fy+c + \lambda(lx+my+n) = 0.
\]
This represents a circle for all $\lambda$ (coefficient of $x^2,y^2$ remains 1). It passes through the intersection of $S$ and $L$ (real or imaginary).

\begin{exoresolu}[Circle Through Line-Circle Intersection]
Find the circle through the intersection of $x^2+y^2-4x+2y-4=0$ and $x+y-1=0$ that passes through $(2,3)$.
\tcblower
\textbf{Solution.}
Family: $x^2+y^2-4x+2y-4 + \lambda(x+y-1) = 0$.
Passes through $(2,3)$: $4+9-8+6-4 + \lambda(2+3-1) = 0 \implies 7 + 4\lambda = 0 \implies \lambda = -\frac{7}{4}$.
Equation: $4(x^2+y^2-4x+2y-4) - 7(x+y-1) = 0$.
$4x^2+4y^2-16x+8y-16 -7x-7y+7 = 0$.
$\mathbf{4x^2+4y^2-23x+y-9=0}$.
\end{exoresolu}

\coursec{Locus Problems and Synthesis for the GCE Exam}

\courssub{General Strategy for Locus Problems}
\begin{enumerate}
\item Let the moving point be $P(x,y)$.
\item Translate the geometric condition into an algebraic equation in $x,y$.
\item Simplify to a standard form (circle, line, parabola, etc.).
\item Identify the curve and state any restrictions (e.g., excluded points, arcs only).
\end{enumerate}

\courssub{Common Loci Leading to Circles}
\begin{itemize}
\item \textbf{Constant distance from fixed point:} Circle.
\item \textbf{Constant ratio of distances to two fixed points $A,B$ ($k \neq 1$):} Circle (Apollonius Circle).
\item \textbf{Angle in a semicircle / Constant angle subtended by segment $AB$:} Arc of a circle.
\item \textbf{Equal tangent lengths to two circles:} Radical Axis (Line).
\item \textbf{Sum of squares of distances to two fixed points constant:} Circle (centre at midpoint).
\end{itemize}

\begin{exoresolu}[Locus: Ratio of Distances (Apollonius Circle)]
$P(x,y)$ moves such that $PA:PB = 2:1$ for $A(1,0)$, $B(4,0)$. Find locus.
\tcblower
\textbf{Solution.}
$PA = 2 PB \implies PA^2 = 4 PB^2$.
$(x-1)^2+y^2 = 4[(x-4)^2+y^2]$.
$x^2-2x+1+y^2 = 4x^2-32x+64+4y^2$.
$3x^2-30x+63+3y^2=0 \implies x^2-10x+21+y^2=0$.
$(x-5)^2+y^2 = 4$.
\textbf{Locus:} Circle centre $(5,0)$, radius $2$.
\end{exoresolu}

\begin{exoresolu}[Locus: Equal Tangents to Two Circles]
Find the locus of points from which tangents to circles $x^2+y^2=4$ and $x^2+y^2-6x+8=0$ have equal length.
\tcblower
\textbf{Solution.}
Lengths: $\sqrt{x^2+y^2-4}$ and $\sqrt{x^2+y^2-6x+8}$.
Equate squares: $x^2+y^2-4 = x^2+y^2-6x+8 \implies 6x = 12 \implies \mathbf{x=2}$.
Locus is the \textbf{radical axis} (vertical line $x=2$).
\end{exoresolu}

\courssub{Synthesis: GCE Exam Style Questions}

\begin{exercice}[GCE Style Practice]
\begin{enumerate}
\item A circle has equation $x^2+y^2-6x+8y+k=0$. Given that the radius is $5$, find $k$ and the coordinates of the centre.
\item The line $y=2x-1$ is a tangent to the circle $x^2+y^2+4x-2y+c=0$. Find $c$ and the point of contact.
\item Find the equation of the circle passing through the points $(1,2)$, $(3,-4)$ and $(5,2)$.
\item Two circles $C_1: x^2+y^2=25$ and $C_2: x^2+y^2-10x-10y+25=0$ intersect. Find the equation of their common chord and its length.
\item A point $P$ moves so that the sum of the squares of its distances from $A(0,0)$ and $B(4,0)$ is $20$. Find the equation of the locus of $P$ and identify the curve.
\item Find the equations of the tangents from the point $(4,0)$ to the circle $x^2+y^2=4$.
\item A circle touches the $x$-axis at $(3,0)$ and passes through $(1,2)$. Find its equation.
\end{enumerate}
\end{exercice}

\begin{corrige}[Exercise 1]
$g=-3, f=4$. $r^2 = g^2+f^2-c = 9+16-k = 25-k$.
Given $r=5 \implies 25-k=25 \implies k=0$.
Centre $(3,-4)$.
\end{corrige}

\begin{corrige}[Exercise 2]
Circle: $(x+2)^2+(y-1)^2 = 5-c$. Centre $(-2,1)$, $r^2=5-c$.
Line: $2x-y-1=0$. Distance from centre to line $= r$.
$r = \frac{|2(-2)-1(1)-1|}{\sqrt{5}} = \frac{6}{\sqrt{5}} \implies r^2 = \frac{36}{5}$.
$5-c = \frac{36}{5} \implies c = 5 - \frac{36}{5} = -\frac{11}{5}$.
Point of contact: Foot of perpendicular from $(-2,1)$ to line.
Line perp: $x+2y=0$. Intersection with $2x-y=1$: $x+2(2x-1)=0 \implies 5x=2 \implies x=0.4, y=-0.2$.
Contact: $(0.4, -0.2)$.
\end{corrige}

\begin{corrige}[Exercise 3]
Using method from Section 3 (or perpendicular bisectors).
Centre $(3, -2/3)$, Radius $10/3$.
Equation: $3x^2+3y^2-18x+4y-5=0$.
\end{corrige}

\begin{corrige}[Exercise 4]
$C_1$: Centre $(0,0)$, $r_1=5$. $C_2$: $(x-5)^2+(y-5)^2=25$, Centre $(5,5)$, $r_2=5$.
Radical axis (common chord): $S_1-S_2=0 \implies 10x+10y-50=0 \implies x+y-5=0$.
Distance from $C_1(0,0)$ to chord: $d = 5/\sqrt{2}$.
Half-chord $= \sqrt{25 - 25/2} = 5/\sqrt{2}$.
Chord length $= 2 \times 5/\sqrt{2} = 5\sqrt{2}$.
\end{corrige}

\begin{corrige}[Exercise 5]
$PA^2+PB^2=20 \implies x^2+y^2 + (x-4)^2+y^2 = 20$.
$2x^2-8x+16+2y^2=20 \implies x^2-4x+y^2=2 \implies (x-2)^2+y^2=6$.
Circle centre $(2,0)$, radius $\sqrt{6}$.
\end{corrige}

\begin{corrige}[Exercise 6]
Circle $x^2+y^2=4$. Point $(4,0)$.
Tangent form $y=m(x-4)$. Substitute: $x^2+m^2(x-4)^2=4$.
$(1+m^2)x^2 - 8m^2 x + 16m^2-4=0$.
$\Delta=0 \implies 64m^4 - 4(1+m^2)(16m^2-4)=0 \implies 16m^4 - (16m^4+12m^2-4)=0 \implies 12m^2=4 \implies m=\pm 1/\sqrt{3}$.
Tangents: $y = \pm \frac{1}{\sqrt{3}}(x-4)$ or $x \pm \sqrt{3}y - 4 = 0$.
Points of contact: $(1, \pm \sqrt{3})$.
\end{corrige}

\begin{corrige}[Exercise 7]
Touches $x$-axis at $(3,0) \implies$ Centre $(3, r)$ or $(3, -r)$. Radius $r$.
Passes through $(1,2)$: Distance from centre to $(1,2) = r$.
Case 1: Centre $(3, r)$. $(3-1)^2+(r-2)^2 = r^2 \implies 4 + r^2-4r+4 = r^2 \implies 8=4r \implies r=2$.
Centre $(3,2)$. Eq: $(x-3)^2+(y-2)^2=4$.
Case 2: Centre $(3, -r)$. $(3-1)^2+(-r-2)^2 = r^2 \implies 4 + r^2+4r+4 = r^2 \implies 8=-4r \implies r=-2$ (impossible, $r>0$).
Only one circle: $\mathbf{(x-3)^2+(y-2)^2=4}$.
\end{corrige}

\begin{aretenir}[Key Formulas Summary for GCE Revision]
\begin{itemize}
\item \textbf{Standard Form:} $(x-a)^2+(y-b)^2=r^2$. Centre $(a,b)$, Radius $r$.
\item \textbf{General Form:} $x^2+y^2+2gx+2fy+c=0$. Centre $(-g,-f)$, Radius $\sqrt{g^2+f^2-c}$.
\item \textbf{Diameter Form (endpoints $A,B$):} $(x-x_1)(x-x_2)+(y-y_1)(y-y_2)=0$.
\item \textbf{Tangent at $P(x_1,y_1)$ on $S=0$:} $xx_1+yy_1+g(x+x_1)+f(y+y_1)+c=0$.
\item \textbf{Tangent Condition ($y=mx+d$ to $x^2+y^2=r^2$):} $d^2=r^2(1+m^2)$.
\item \textbf{Length of Tangent from $P(x_1,y_1)$:} $\sqrt{S_{11}} = \sqrt{x_1^2+y_1^2+2gx_1+2fy_1+c}$.
\item \textbf{Parametric:} $x=a+r\cos\theta, y=b+r\sin\theta$.
\item \textbf{Radical Axis of $S_1, S_2$:} $S_1-S_2=0$.
\item \textbf{Family through $S_1 \cap S_2$:} $S_1+\lambda S_2=0$.
\item \textbf{Touching Circles:} $d = r_1+r_2$ (external), $d = |r_1-r_2|$ (internal).
\end{itemize}
\end{aretenir}

\newpage

\coursec{Exercises}

\courssub{I check my knowledge}

\begin{exercice}[Multiple choice]
Choose the correct answer, justifying your choice briefly.
\begin{enumerate}
\item The circle $x^2 + y^2 - 4x + 6y - 12 = 0$ has centre and radius:
\begin{enumerate}
\item centre $(2, -3)$, radius $5$;
\item centre $(-2, 3)$, radius $5$;
\item centre $(2, -3)$, radius $1$;
\item centre $(4, -6)$, radius $12$.
\end{enumerate}
\item The line $y = 2$ and the circle $x^2 + y^2 = 4$:
\begin{enumerate}
\item meet in two distinct points;
\item are tangent to each other;
\item do not meet;
\item meet at the origin.
\end{enumerate}
\item The length of the tangent from the point $(0, 0)$ to the circle $x^2 + y^2 - 6x - 8y + 9 = 0$ is:
\begin{enumerate}
\item $9$; \item $3$; \item $5$; \item $\sqrt{41}$.
\end{enumerate}
\end{enumerate}
\end{exercice}

\begin{exercice}[True or false?]
State whether each assertion is true or false, justifying your answer.
\begin{enumerate}
\item The equation $x^2 + y^2 + 2x - 4y + 9 = 0$ represents a circle of radius $2$.
\item The point $(3, 4)$ lies on the circle $x^2 + y^2 = 25$.
\item Two circles touch externally if and only if the distance between their centres equals the sum of their radii.
\item The parametric equations $x = 1 + 3\cos\theta$, $y = -2 + 3\sin\theta$ represent a circle of centre $(1, -2)$ and radius $9$.
\end{enumerate}
\end{exercice}

\begin{exercice}[Definitions]
\begin{enumerate}
\item Define the \cle{radical axis} of two circles and state one of its geometric properties.
\item State the condition on the distance $d$ between the centres of two circles of radii $r_1$ and $r_2$ ($r_1 \geq r_2$) for the circles to touch \emph{internally}.
\item Write down the general equation of a circle and state the conditions on its coefficients for it to represent a real circle.
\end{enumerate}
\end{exercice}

\begin{exercice}[Direct calculations]
\begin{enumerate}
\item Write down the equation of the circle of centre $(0, 0)$ and radius $7$.
\item Write down the equation of the circle of centre $(-1, 4)$ and radius $\sqrt{5}$.
\item Find the centre and radius of the circle $x^2 + y^2 + 8x - 10y + 5 = 0$.
\item Find the equation of the circle having the segment joining $A(1, 3)$ and $B(5, 7)$ as a diameter.
\end{enumerate}
\end{exercice}

\courssub{I practise}

\begin{exercice}[Centre and a point]
Find the equation of the circle with centre $(2, -3)$ passing through the point $(5, 1)$. Determine whether the point $(6, -2)$ lies inside, on or outside this circle.
\end{exercice}

\begin{exercice}[Tangent at a point]
A circle has equation $x^2 + y^2 - 6x + 4y - 12 = 0$.
\begin{enumerate}
\item Find its centre and its radius.
\item Verify that the point $(6, 2)$ lies on the circle.
\item Find the equation of the tangent to the circle at the point $(6, 2)$.
\end{enumerate}
\end{exercice}

\begin{exercice}[Tangency of a line]
Show that the line $3x + 4y = 25$ is a tangent to the circle $x^2 + y^2 = 25$, and find the coordinates of the point of contact.
\end{exercice}

\begin{exercice}[Length of a tangent]
Find the length of the tangent from the point $(7, 1)$ to the circle $x^2 + y^2 - 4x + 6y - 3 = 0$.
\end{exercice}

\courssub{I prepare for the GCE}

\begin{exercice}[Circle through three points and a line condition]
\begin{enumerate}
\item Find the equation of the circle passing through the points $A(1, 1)$, $B(2, -1)$ and $C(3, 2)$. \hfill (6 marks)
\item Find the equation of the tangent to this circle at the point $A(1, 1)$. \hfill (4 marks)
\item A second circle passes through the points $(2, 0)$ and $(6, 0)$ and has radius $5$. Show that there are two such circles and find the equation of each. \hfill (6 marks)
\end{enumerate}
\end{exercice}

\begin{exercice}[Two circles: intersection, common chord and touching]
The circles $C_1$ and $C_2$ have equations
\[ C_1 : x^2 + y^2 = 10, \qquad C_2 : x^2 + y^2 - 6x - 2y = 10. \]
\begin{enumerate}
\item Find the equation of the common chord of $C_1$ and $C_2$. \hfill (3 marks)
\item Find the coordinates of the points of intersection of the two circles. \hfill (5 marks)
\item Find the equation of the circle passing through the points of intersection of $x^2 + y^2 - 4 = 0$ and $x^2 + y^2 - 2x - 4y = 0$ and through the point $(1, 2)$. \hfill (5 marks)
\item Show that the circles $x^2 + y^2 - 2x - 4y - 20 = 0$ and $x^2 + y^2 + 6x + 2y - 90 = 0$ touch internally, and find the coordinates of their point of contact. \hfill (5 marks)
\end{enumerate}
\end{exercice}

\begin{exercice}[Locus and circles under geometric conditions]
\begin{enumerate}
\item A point $P$ moves so that the sum of the squares of its distances from the points $A(-2, 0)$ and $B(2, 0)$ is equal to $20$. Show that the locus of $P$ is a circle, and find its centre and radius. \hfill (5 marks)
\item Find the equations of the circles which touch both coordinate axes and pass through the point $(2, 1)$. \hfill (5 marks)
\item Find the equation of the circle which passes through the point $(4, 1)$ and touches the line $x + y = 1$ at the point $(2, -1)$. \hfill (6 marks)
\end{enumerate}
\end{exercice}



\newpage

\coursec{Solutions to the exercises}

\begin{corrige}[Exercise 1 — Multiple choice]
\begin{enumerate}
\item Completing the square for $x^2 - 4x + y^2 + 6y = 12$:
\[ (x^2 - 4x + 4) + (y^2 + 6y + 9) = 12 + 4 + 9 \;\Longrightarrow\; (x-2)^2 + (y+3)^2 = 25. \]
Centre $(2, -3)$, radius $5$. Answer \textbf{(a)}.
\emph{Note}: Answers (b) and (d) arise from sign errors on $-g$ and $-f$; answer (c) forgets to take the square root of $25$.
\item Circle $x^2 + y^2 = 4$ (centre $O(0,0)$, radius $r=2$). Line $y=2$ (i.e. $y-2=0$). Perpendicular distance from $O$ to the line:
\[ d = \frac{|0 - 2|}{\sqrt{1}} = 2 = r. \]
The line is a tangent. Answer \textbf{(b)}. Point of contact is $(0,2)$, which satisfies $0^2+2^2=4$.
\item Circle $x^2 + y^2 - 6x - 8y + 9 = 0$. Centre $C(3,4)$, radius $r = \sqrt{9+16-9}=4$. Point $P(0,0)$.
Power of $P$: $S(0,0) = 0 + 0 - 0 - 0 + 9 = 9$. Length of tangent $\ell = \sqrt{9} = 3$. Answer \textbf{(b)}.
\emph{Note}: Answer (a) gives the power (squared length) instead of the length.
\end{enumerate}
\end{corrige}

\begin{corrige}[Exercise 2 — True or false?]
\begin{enumerate}
\item \textbf{False.} For $x^2+y^2+2x-4y+9=0$, we have $g=1$, $f=-2$, $c=9$.
$g^2+f^2-c = 1+4-9 = -4 < 0$. The equation represents no real circle (the ``radius squared'' would be negative).
\item \textbf{True.} Substituting $(3,4)$ into $x^2+y^2=25$: $9+16=25$. The point lies on the circle.
\item \textbf{True.} Two circles touch externally if and only if the distance between centres equals the sum of the radii: $d = r_1 + r_2$. They then have exactly one common point, which lies on the line of centres between the two centres.
\item \textbf{False.} Parametric equations $x = 1 + 3\cos\theta$, $y = -2 + 3\sin\theta$.
Centre is $(1, -2)$. The radius is the common coefficient of $\cos\theta$ and $\sin\theta$, namely $3$ — not $3^2=9$.
Indeed $(x-1)^2+(y+2)^2 = 9\cos^2\theta+9\sin^2\theta = 9$, so $r=3$.
\end{enumerate}
\end{corrige}

\begin{corrige}[Exercise 3 — Definitions]
\begin{enumerate}
\item The \cle{radical axis} of two circles is the locus of points having equal powers with respect to the two circles. Algebraically, if the circles are $S_1=0$ and $S_2=0$, the radical axis is $S_1 - S_2 = 0$, which is always a \textbf{straight line}.
Geometric properties:
\begin{itemize}
\item The radical axis is perpendicular to the line of centres of the two circles.
\item If the circles intersect in two points, the radical axis is the line through these points (the common chord).
\item If the circles touch, the radical axis is their common tangent at the point of contact.
\end{itemize}
\item Two circles with radii $r_1 \ge r_2$ touch \emph{internally} if and only if the distance between their centres equals the difference of the radii:
\[ d = r_1 - r_2. \]
\item The general equation of a circle is
\[ x^2 + y^2 + 2gx + 2fy + c = 0. \]
It represents a real circle if and only if $g^2 + f^2 - c > 0$; the centre is then $(-g, -f)$ and the radius $\sqrt{g^2 + f^2 - c}$.
(If $g^2+f^2-c = 0$ the equation reduces to a single point — a point circle; if it is negative there is no real locus.)
\end{enumerate}
\end{corrige}

\begin{corrige}[Exercise 4 — Direct calculations]
\begin{enumerate}
\item Centre $(0,0)$, radius $7$: $x^2 + y^2 = 49$.
\item Centre $(-1, 4)$, radius $\sqrt{5}$:
\[ (x+1)^2 + (y-4)^2 = 5 \;\Longrightarrow\; x^2 + 2x + 1 + y^2 - 8y + 16 = 5, \]
\[ \boxed{x^2 + y^2 + 2x - 8y + 12 = 0.} \]
\item $x^2 + y^2 + 8x - 10y + 5 = 0$. Completing the square:
\[ (x^2+8x+16) + (y^2-10y+25) = -5 + 16 + 25, \]
\[ (x+4)^2 + (y-5)^2 = 36. \]
Centre $(-4, 5)$, radius $6$.
\item Diameter form $(x - x_1)(x - x_2) + (y - y_1)(y - y_2) = 0$ with $A(1,3)$, $B(5,7)$:
\[ (x-1)(x-5) + (y-3)(y-7) = 0 \;\Longrightarrow\; x^2 - 6x + 5 + y^2 - 10y + 21 = 0, \]
\[ \boxed{x^2 + y^2 - 6x - 10y + 26 = 0.} \]
\emph{Check}: Centre $(3, 5)$ is the midpoint of $AB$; radius $\sqrt{9 + 25 - 26} = \sqrt{8} = 2\sqrt{2}$.
$AB = \sqrt{(5-1)^2 + (7-3)^2} = \sqrt{16+16} = \sqrt{32} = 4\sqrt{2}$, so radius $= \frac{1}{2}AB$. \checkmark
\end{enumerate}
\end{corrige}

\begin{corrige}[Exercise 5 — Centre and a point on the circle]
Centre $C(2, -3)$, point on circle $P(5, 1)$.
Radius squared:
\[ r^2 = CP^2 = (5-2)^2 + (1-(-3))^2 = 3^2 + 4^2 = 9 + 16 = 25, \quad r = 5. \]
Equation in centre-radius form:
\[ (x-2)^2 + (y+3)^2 = 25. \]
Expanded form:
\[ x^2 - 4x + 4 + y^2 + 6y + 9 = 25 \;\Longrightarrow\; \boxed{x^2 + y^2 - 4x + 6y - 12 = 0.} \]

For the point $Q(6, -2)$, compute its squared distance to the centre:
\[ CQ^2 = (6-2)^2 + (-2+3)^2 = 16 + 1 = 17. \]
Since $17 < 25 = r^2$, the point $Q$ lies \textbf{inside} the circle.
\end{corrige}

\begin{corrige}[Exercise 6 — Tangent at a given point]
\begin{enumerate}
\item Circle: $x^2 + y^2 - 6x + 4y - 12 = 0$.
Completing the square:
\[ (x^2-6x+9) + (y^2+4y+4) = 12 + 9 + 4 \;\Longrightarrow\; (x-3)^2 + (y+2)^2 = 25. \]
Centre $C(3, -2)$, radius $r = 5$.
\item Check if $T(6, 2)$ lies on the circle:
\[ 6^2 + 2^2 - 6(6) + 4(2) - 12 = 36 + 4 - 36 + 8 - 12 = 0. \]
Yes, $T$ lies on the circle.
\item Gradient of radius $CT$:
\[ m_{CT} = \frac{2 - (-2)}{6 - 3} = \frac{4}{3}. \]
The tangent at $T$ is perpendicular to $CT$, so its gradient is $m_T = -\frac{3}{4}$.
Equation of tangent:
\[ y - 2 = -\frac{3}{4}(x - 6) \;\Longrightarrow\; 4y - 8 = -3x + 18 \;\Longrightarrow\; \boxed{3x + 4y - 26 = 0.} \]
\emph{Verification using the ``$T=0$'' formula} for $x^2+y^2+2gx+2fy+c=0$ at $(x_1,y_1)$:
$xx_1+yy_1+g(x+x_1)+f(y+y_1)+c=0$. Here $g=-3$, $f=2$, $c=-12$, $(x_1,y_1)=(6,2)$:
\[ 6x + 2y - 3(x+6) + 2(y+2) - 12 = 0 \;\Longrightarrow\; 3x + 4y - 18 + 4 - 12 = 0 \;\Longrightarrow\; 3x + 4y - 26 = 0. \quad \checkmark \]
\end{enumerate}
\end{corrige}

\begin{corrige}[Exercise 7 — Tangency of a line to a circle]
Circle $C: x^2 + y^2 = 25$ (centre $O(0,0)$, radius $r=5$).
Line $L: 3x + 4y - 25 = 0$.
Perpendicular distance from centre $O$ to line $L$:
\[ d = \frac{|3(0) + 4(0) - 25|}{\sqrt{3^2 + 4^2}} = \frac{25}{5} = 5. \]
Since $d = r = 5$, the line is a \textbf{tangent} to the circle.

The point of contact $T$ is the foot of the perpendicular from $O$ to $L$.
The normal vector to $L$ is $(3,4)$. Parametric form of perpendicular through $O$: $(x,y) = (3t, 4t)$.
Substitute into $L$: $3(3t) + 4(4t) - 25 = 0 \;\Longrightarrow\; 25t = 25 \;\Longrightarrow\; t = 1$.
Hence $T = (3, 4)$.
\emph{Check}: $T$ lies on $L$: $3(3)+4(4)=25$. $T$ lies on $C$: $3^2+4^2=25$. \checkmark
\end{corrige}

\begin{corrige}[Exercise 8 — Length of a tangent from an external point]
Circle: $x^2 + y^2 - 4x + 6y - 3 = 0$. External point $P(7, 1)$.
Length of tangent $\ell = \sqrt{S(7,1)}$, where $S(x,y) = x^2+y^2-4x+6y-3$.
\[ \ell = \sqrt{7^2 + 1^2 - 4(7) + 6(1) - 3} = \sqrt{49 + 1 - 28 + 6 - 3} = \sqrt{25} = \boxed{5}. \]

\emph{Geometric check}: Centre $C(2, -3)$, radius $r = \sqrt{4+9+3} = 4$.
Distance $CP = \sqrt{(7-2)^2 + (1+3)^2} = \sqrt{25 + 16} = \sqrt{41}$.
By Pythagoras in right triangle $CPT$ (where $T$ is point of contact): $\ell^2 = CP^2 - r^2 = 41 - 16 = 25$, so $\ell = 5$. \checkmark
\end{corrige}

\begin{corrige}[Exercise 9 — Circle through three points and tangent condition]
\begin{enumerate}
\item Let the circle be $x^2 + y^2 + 2gx + 2fy + c = 0$. Substituting the three points:
\begin{align*}
A(1,1):&\quad 1+1+2g+2f+c = 0 \;\Longrightarrow\; 2g+2f+c = -2 \quad (1)\\
B(2,-1):&\quad 4+1+4g-2f+c = 0 \;\Longrightarrow\; 4g-2f+c = -5 \quad (2)\\
C(3,2):&\quad 9+4+6g+4f+c = 0 \;\Longrightarrow\; 6g+4f+c = -13 \quad (3)
\end{align*}
$(2)-(1)$: $2g - 4f = -3 \quad (4)$\\
$(3)-(2)$: $2g + 6f = -8 \quad (5)$\\
$(5)-(4)$: $10f = -5 \;\Longrightarrow\; f = -\frac{1}{2}$.
From $(4)$: $2g = 4f - 3 = -2 - 3 = -5 \;\Longrightarrow\; g = -\frac{5}{2}$.
From $(1)$: $c = -2 - 2g - 2f = -2 + 5 + 1 = 4$.
Equation:
\[ \boxed{x^2 + y^2 - 5x - y + 4 = 0.} \]
Centre $\left(\frac{5}{2}, \frac{1}{2}\right)$, radius $r = \sqrt{\frac{25}{4} + \frac{1}{4} - 4} = \sqrt{\frac{10}{4}} = \sqrt{\frac{5}{2}}$.
\emph{Examiner's expectation}: Any correct method (general equation or perpendicular bisectors) earns full credit; the three equations must be set up correctly before any method mark for solving is awarded.

\item Tangent at $A(1,1)$. Centre $O\left(\frac{5}{2}, \frac{1}{2}\right)$.
Gradient of radius $OA$:
\[ m_{OA} = \frac{1 - \frac{1}{2}}{1 - \frac{5}{2}} = \frac{\frac{1}{2}}{-\frac{3}{2}} = -\frac{1}{3}. \]
Gradient of tangent (perpendicular to radius): $m_T = 3$.
Equation:
\[ y - 1 = 3(x - 1) \;\Longrightarrow\; \boxed{3x - y - 2 = 0.} \]

\item Circle of radius $5$ passing through $(2,0)$ and $(6,0)$.
The centre is equidistant from $(2,0)$ and $(6,0)$, so it lies on the perpendicular bisector of this segment.
Midpoint is $(4,0)$, segment is horizontal, so perpendicular bisector is the vertical line $x = 4$.
Let centre be $(4, k)$. Then
\[ r^2 = (4-2)^2 + (k-0)^2 = 4 + k^2 = 25 \;\Longrightarrow\; k^2 = 21 \;\Longrightarrow\; k = \pm\sqrt{21}. \]
Since $k$ takes \textbf{two} distinct values, there are exactly \textbf{two} such circles.
Centres: $(4, \sqrt{21})$ and $(4, -\sqrt{21})$.
Equations:
\[ \boxed{(x-4)^2 + (y - \sqrt{21})^2 = 25 \quad\text{and}\quad (x-4)^2 + (y + \sqrt{21})^2 = 25.} \]
Expanded: $x^2 + y^2 - 8x \mp 2\sqrt{21}\,y + 16 + 21 - 25 = 0$, i.e.
\[ \boxed{x^2 + y^2 - 8x \mp 2\sqrt{21}\,y + 12 = 0.} \]
\emph{Examiner's expectation}: The explicit statement ``two values of $k$, hence two circles'' is required for the justification mark.
\end{enumerate}
\end{corrige}

\begin{corrige}[Exercise 10 — Intersection of two circles, common chord and family of circles]
\textbf{Given circles:}
\[ C_1:\; x^2 + y^2 - 10 = 0 \qquad\text{(Centre $O_1(0,0)$, radius $r_1=\sqrt{10}$)} \]
\[ C_2:\; x^2 + y^2 - 6x - 2y - 10 = 0 \qquad\text{(Centre $O_2(3,1)$, radius $r_2=\sqrt{9+1+10}=\sqrt{20}=2\sqrt{5}$)} \]

\begin{enumerate}
\item The common chord (radical axis) is $C_1 - C_2 = 0$:
\[ (x^2+y^2-10) - (x^2+y^2-6x-2y-10) = 0 \;\Longrightarrow\; 6x + 2y = 0 \;\Longrightarrow\; \boxed{3x + y = 0.} \]

\item On the chord, $y = -3x$. Substitute into $C_1$:
\[ x^2 + (-3x)^2 = 10 \;\Longrightarrow\; 10x^2 = 10 \;\Longrightarrow\; x^2 = 1 \;\Longrightarrow\; x = \pm 1. \]
Corresponding $y$: $x=1 \Rightarrow y=-3$; $x=-1 \Rightarrow y=3$.
Points of intersection:
\[ \boxed{(1, -3) \quad\text{and}\quad (-1, 3).} \]
\emph{Check in $C_2$}: $1+9-6+6-10=0$ and $1+9+6-6-10=0$. \checkmark

\item Family of circles through the intersections of $C_1$ and $C_2$:
\[ C_1 + \lambda C_2 = 0 \;\Longrightarrow\; (x^2+y^2-10) + \lambda(x^2+y^2-6x-2y-10) = 0. \]
This circle passes through $(1,2)$:
\[ (1+4-10) + \lambda(1+4-6-4-10) = 0 \;\Longrightarrow\; -5 + \lambda(-15) = 0 \;\Longrightarrow\; \lambda = -\frac{1}{3}. \]
Substitute $\lambda = -\frac{1}{3}$ and multiply by $3$:
\[ 3(x^2+y^2-10) - (x^2+y^2-6x-2y-10) = 0 \]
\[ 3x^2+3y^2-30 - x^2 - y^2 + 6x + 2y + 10 = 0 \]
\[ \boxed{2x^2 + 2y^2 + 6x + 2y - 20 = 0 \quad\text{or}\quad x^2 + y^2 + 3x + y - 10 = 0.} \]
\end{enumerate}
\end{corrige}

\begin{corrige}[Exercise 10 (continued) — Conditions for circles to touch]
\textbf{Given circles:}
\[ C_1:\; x^2 + y^2 - 2x - 4y - 20 = 0 \quad\Rightarrow\quad \text{Centre } O_1(1,2),\; r_1 = \sqrt{1+4+20} = 5. \]
\[ C_2:\; x^2 + y^2 + 6x + 2y - 90 = 0 \quad\Rightarrow\quad \text{Centre } O_2(-3,-1),\; r_2 = \sqrt{9+1+90} = 10. \]

Distance between centres:
\[ d = O_1O_2 = \sqrt{(1+3)^2 + (2+1)^2} = \sqrt{16+9} = 5. \]
Compare $d$ with radii:
\[ r_2 - r_1 = 10 - 5 = 5 = d \quad\text{and}\quad d \neq 0. \]
Since $d = |r_2 - r_1|$, the circles touch \textbf{internally}.

The point of contact $T$ lies on the line of centres $O_1O_2$, on the side of $O_1$ opposite to $O_2$ (since $r_2 > r_1$).
Vector $\overrightarrow{O_2O_1} = (4, 3)$. Unit vector $\left(\frac{4}{5}, \frac{3}{5}\right)$.
Starting from $O_1$ and moving a distance $r_1=5$ along this direction:
\[ T = O_1 + 5\left(\frac{4}{5}, \frac{3}{5}\right) = (1,2) + (4,3) = \boxed{(5, 5).} \]
\emph{Check}: $O_2T = \sqrt{(5+3)^2 + (5+1)^2} = \sqrt{64+36} = 10 = r_2$.
Also $T$ satisfies $C_1$: $25+25-10-20-20=0$. \checkmark
\end{corrige}

\begin{corrige}[Exercise 11 — Locus problems and circles under geometric conditions]
\begin{enumerate}
\item Let $P(x,y)$ be the moving point. Fixed points $A(-2,0)$, $B(2,0)$.
Condition: $PA^2 + PB^2 = 20$.
\[ (x+2)^2 + y^2 + (x-2)^2 + y^2 = 20. \]
Expand:
\[ (x^2+4x+4) + y^2 + (x^2-4x+4) + y^2 = 20 \]
\[ 2x^2 + 2y^2 + 8 = 20 \;\Longrightarrow\; 2x^2 + 2y^2 = 12 \;\Longrightarrow\; \boxed{x^2 + y^2 = 6.} \]
This is a circle with centre $(0,0)$ — the midpoint of $AB$ — and radius $\sqrt{6}$.

\item A circle in the first quadrant touching both coordinate axes has centre $(a,a)$ and radius $a$ ($a>0$).
It passes through $(2,1)$:
\[ (2-a)^2 + (1-a)^2 = a^2 \]
\[ 4 - 4a + a^2 + 1 - 2a + a^2 = a^2 \]
\[ a^2 - 6a + 5 = 0 \;\Longrightarrow\; (a-1)(a-5) = 0 \;\Longrightarrow\; a = 1 \text{ or } a = 5. \]
Two circles satisfy the conditions:
\[ \boxed{(x-1)^2 + (y-1)^2 = 1 \qquad\text{and}\qquad (x-5)^2 + (y-5)^2 = 25.} \]
\emph{Check for $a=5$}: $(2-5)^2+(1-5)^2 = 9+16=25$. \checkmark
\end{enumerate}
\end{corrige}

\begin{popfigure}
\begin{tikzpicture}[scale=0.7, >=stealth]
\draw[popmuted, thin] (-1,-1) grid (7,7);
\draw[->, thick] (-1,0) -- (7,0) node[right] {$x$};
\draw[->, thick] (0,-1) -- (0,7) node[above] {$y$};
% Circle 1: centre (1,1), r=1
\draw[popblue, thick] (1,1) circle (1);
\fill[popblue] (1,1) circle (2pt) node[above right] {$C_1(1,1)$};
% Circle 2: centre (5,5), r=5
\draw[poporange, thick] (5,5) circle (5);
\fill[poporange] (5,5) circle (2pt) node[above right] {$C_2(5,5)$};
% Point (2,1)
\fill[popdark] (2,1) circle (2pt) node[below right] {$P(2,1)$};
% Axes labels
\foreach \x in {1,2,3,4,5,6} \draw (\x,0.1) -- (\x,-0.1) node[below] {\x};
\foreach \y in {1,2,3,4,5,6} \draw (0.1,\y) -- (-0.1,\y) node[left] {\y};
\end{tikzpicture}
\legende{Exercise 11.2: Two circles touching both axes and passing through $(2,1)$.}
\end{popfigure}

\begin{corrige}[Exercise 11 (continued) — Circle touching a line at a given point]
\begin{itemize}
\item Circle touches the line $L: x + y = 1$ at $T(2, -1)$ and passes through $P(4, 1)$.
Since the circle touches $L$ at $T$, the centre $C$ lies on the perpendicular to $L$ at $T$.
Line $L$ has gradient $-1$, so the perpendicular has gradient $1$.
Parametric form of perpendicular through $T(2,-1)$: $C = (2+t, -1+t)$.
Radius squared is the squared distance $CT^2$:
\[ r^2 = t^2 + t^2 = 2t^2. \]
The circle passes through $P(4,1)$, so $CP^2 = r^2$:
\[ (4 - (2+t))^2 + (1 - (-1+t))^2 = 2t^2 \]
\[ (2-t)^2 + (2-t)^2 = 2t^2 \;\Longrightarrow\; 2(2-t)^2 = 2t^2 \;\Longrightarrow\; (2-t)^2 = t^2. \]
Hence $2-t = \pm t$.
Case $2-t = -t \Rightarrow 2=0$ (impossible).
Case $2-t = t \Rightarrow 2t=2 \Rightarrow t=1$.
Centre $C = (3, 0)$, $r^2 = 2(1)^2 = 2$.
Equation:
\[ \boxed{(x-3)^2 + y^2 = 2 \quad\text{i.e.}\quad x^2 + y^2 - 6x + 7 = 0.} \]

\emph{Verification}:
Distance from $C(3,0)$ to line $x+y-1=0$: $\frac{|3+0-1|}{\sqrt{2}} = \frac{2}{\sqrt{2}} = \sqrt{2} = r$. \checkmark
Foot of perpendicular from $C$ to $L$:
$C - \frac{3+0-1}{2}(1,1) = (3,0) - (1,1) = (2,-1) = T$. \checkmark
$P(4,1)$: $(4-3)^2+1^2 = 2 = r^2$. \checkmark
\end{itemize}
\end{corrige}

\begin{popfigure}
\begin{tikzpicture}[scale=0.8, >=stealth]
\draw[popmuted, thin] (-1,-3) grid (6,3);
\draw[->, thick] (-1,0) -- (6,0) node[right] {$x$};
\draw[->, thick] (0,-3) -- (0,3) node[above] {$y$};
% Line x+y=1
\draw[popgreen, thick] (-1,2) -- (4,-3) node[below right] {$x+y=1$};
% Circle centre (3,0) r=sqrt(2) approx 1.414
\draw[popblue, thick] (3,0) circle (1.414);
\fill[popblue] (3,0) circle (2pt) node[below] {$C(3,0)$};
% Point of contact T(2,-1)
\fill[popdark] (2,-1) circle (2pt) node[below left] {$T(2,-1)$};
% Point P(4,1)
\fill[popdark] (4,1) circle (2pt) node[above right] {$P(4,1)$};
% Radius CT
\draw[popblue, dashed] (3,0) -- (2,-1);
% Perpendicular from C to line
\draw[popmuted, dashed] (3,0) -- (2,-1);
% Axes labels
\foreach \x in {-1,1,2,3,4,5} \draw (\x,0.1) -- (\x,-0.1) node[below] {\x};
\foreach \y in {-2,-1,1,2} \draw (0.1,\y) -- (-0.1,\y) node[left] {\y};
\end{tikzpicture}
\legende{Exercise 11.3: Circle touching $x+y=1$ at $(2,-1)$ and passing through $(4,1)$.}
\end{popfigure}

\newpage

\coursec{Summary sheet}

\begin{popfigure}
\begin{tikzpicture}[mindmap, grow cyclic, every node/.style={concept, text=white, font=\small},
  root concept/.append style={concept color=popink, font=\bfseries},
  level 1/.append style={level distance=3.4cm, sibling angle=60},
  level 1 concept/.append style={font=\footnotesize}]
\node[root concept]{Circle Geometry}
  child[concept color=popblue]{ node {Equations: $x^2+y^2=r^2$; $(x-a)^2+(y-b)^2=r^2$; general form} }
  child[concept color=popgreen]{ node {Diameter form; 3 points; centre on a line} }
  child[concept color=poporange]{ node {Tangents: $xx_1+yy_1=r^2$; perpendicular radius; length $\sqrt{S_1}$} }
  child[concept color=poppurple]{ node {Parametric: $x=a+r\cos\theta$, $y=b+r\sin\theta$} }
  child[concept color=popgold]{ node {Line $\cap$ circle: discriminant; chord; family $S+\lambda L=0$} }
  child[concept color=poppink]{ node {Two circles: $S_1-S_2=0$ radical axis; touch: $d=r_1\pm r_2$; family $S_1+\lambda S_2=0$} }
  child[concept color=popmuted]{ node {Locus: distance conditions, $PA=PB$, $PA=k\,PB$} };
\end{tikzpicture}
\legende{Mind map of the chapter Circle Geometry.}
\end{popfigure}

\begin{aretenir}[Circle Geometry — Summary Sheet]
\textbf{Key definitions.}
\begin{itemize}
\item A \cle{circle} is the locus of a point $P$ whose distance from a fixed point (the \cle{centre}) is constant (the \cle{radius}).
\item A \cle{tangent} is a line touching the circle at exactly one point; it is perpendicular to the radius at the point of contact.
\item The \cle{radical axis} of two circles is the locus of points with equal tangent lengths (equal powers) to both circles.
\end{itemize}

\textbf{Equations of a circle (know them all!).}
\begin{itemize}
\item Centre $O(0,0)$, radius $r$: $x^2+y^2=r^2$.
\item Centre $C(a,b)$, radius $r$: $(x-a)^2+(y-b)^2=r^2$.
\item General form: $x^2+y^2+2gx+2fy+c=0$, with centre $(-g,-f)$ and radius $r=\sqrt{g^2+f^2-c}$ (real circle iff $g^2+f^2-c>0$).
\item Diameter form, ends $A(x_1,y_1)$, $B(x_2,y_2)$: $(x-x_1)(x-x_2)+(y-y_1)(y-y_2)=0$ (from $\vec{PA}\cdot\vec{PB}=0$).
\item Parametric form: $x=a+r\cos\theta$, $y=b+r\sin\theta$; convert to Cartesian by $\cos^2\theta+\sin^2\theta=1$.
\item Circle touching both axes: centre $(\pm r,\pm r)$, equation $(x\pm r)^2+(y\pm r)^2=r^2$.
\end{itemize}

\textbf{Tangents.}
\begin{itemize}
\item Tangent at $(x_1,y_1)$ on $x^2+y^2=r^2$: $xx_1+yy_1=r^2$; on the general circle: $xx_1+yy_1+g(x+x_1)+f(y+y_1)+c=0$.
\item Line $y=mx+c$ touches $x^2+y^2=r^2$ iff $c^2=r^2(1+m^2)$; generally: perpendicular distance from centre to line equals $r$.
\item Length of tangent from external point $P(x_1,y_1)$: $\ell=\sqrt{x_1^2+y_1^2+2gx_1+2fy_1+c}=\sqrt{CP^2-r^2}$.
\end{itemize}

\textbf{Intersections and families.}
\begin{itemize}
\item Line $\cap$ circle: substitute, get a quadratic; $\Delta>0$ two points, $\Delta=0$ tangent, $\Delta<0$ no intersection.
\item Family through intersection of circle $S=0$ and line $L=0$: $S+\lambda L=0$.
\item Family through intersections of circles $S_1=0$, $S_2=0$: $S_1+\lambda S_2=0$; radical axis: $S_1-S_2=0$ (a straight line, the common chord if they intersect).
\item Two circles touch externally iff $d=r_1+r_2$; internally iff $d=|r_1-r_2|$; intersect iff $|r_1-r_2|<d<r_1+r_2$; separate iff $d>r_1+r_2$.
\end{itemize}

\textbf{Methods.}
\begin{itemize}
\item Circle through 3 points: substitute into the general equation, solve the $3\times3$ system for $g,f,c$ — or use perpendicular bisectors of two chords.
\item Centre on a line: write the centre as $(h,k)$ with the line's relation, then equate radii to the given points.
\item Locus: let $P(x,y)$, translate the geometric condition into an equation, simplify.
\end{itemize}

\textbf{Common mistakes.}
\begin{itemize}
\item Sign errors: centre is $(-g,-f)$, not $(g,f)$; radius is $\sqrt{g^2+f^2-c}$, not $\sqrt{g^2+f^2+c}$.
\item Forgetting to divide by the coefficient of $x^2$ before reading off $g,f,c$.
\item Using $d=r_1+r_2$ for internal touching.
\item Mixing up point of tangency and external point in tangent formulas.
\end{itemize}

\textbf{GCE tips.}
\begin{itemize}
\item Always state centre and radius first — most questions build on them.
\item Check your final circle passes through all given points.
\item Show the discriminant or distance condition clearly: method marks are awarded for the condition, not just the answer.
\end{itemize}
\end{aretenir}

\findecours

\end{document}
